% Generated mechanically from the frozen schemes.tex target.
% Do not edit this generated file; edit the bound locale source instead.

% OK, start here.
%

\setcounter{chapter}{25}
\chapter{概形}



\phantomsection
\label{schemes-section-phantom}


\section{引言}
\label{schemes-section-introduction}

\noindent
本文定义概形。
一部基本参考文献是 \cite{EGA}。









\section{局部环层空间}
\label{schemes-section-locally-ringed-spaces}

\noindent
回顾《层》第 \ref{sheaves-section-ringed-spaces} 节中定义的环层空间。
简言之，环层空间是一个二元组 $(X, \mathcal{O}_X)$，其中 $X$ 是拓扑空间，
$\mathcal{O}_X$ 是环层。环层空间之间的态射
$f : (X, \mathcal{O}_X) \to (Y, \mathcal{O}_Y)$ 由一个连续映射
$f : X \to Y$ 和一个环层的 $f$-映射
$f^\sharp : \mathcal{O}_Y \to \mathcal{O}_X$ 给出。也可以把 $f^\sharp$
看作映射 $\mathcal{O}_Y \to f_*\mathcal{O}_X$；参见《层》的定义
\ref{sheaves-definition-f-map} 和引理 \ref{sheaves-lemma-f-map}。

\medskip\noindent
一个值得牢记的几何例子是 $\mathcal{C}^\infty$-流形及
$\mathcal{C}^\infty$-流形之间的态射。
具体地，若 $M$ 是 $\mathcal{C}^\infty$-流形，则光滑函数层
$\mathcal{C}^\infty_M$ 是 $M$ 上的环层。流形间的任意映射
$f : M \to N$ 是光滑的，当且仅当对每个局部截面 $h$——这里指
$\mathcal{C}^\infty_N$ 的局部截面——复合 $h \circ f$ 都是
$\mathcal{C}^\infty_M$ 的局部截面。
因此，光滑映射 $f$ 自然给出环层空间之间的态射
$$
f : (M , \mathcal{C}^\infty_M) \longrightarrow (N, \mathcal{C}^\infty_N)
$$
；参见《层》的例 \ref{sheaves-example-continuous-map-ringed}。
考察茎上的情形很有启发性。设 $m \in M$，且其像为 $f(m) = n \in N$。
回顾茎 $\mathcal{C}^\infty_{M, m}$ 是 $m$ 处光滑函数芽组成的环；
参见《层》的例 \ref{sheaves-example-germs-functions}。
$(M, m)$ 上的函数芽代数是一个局部环，其极大理想由在 $m$ 处消失的函数组成。
$\mathcal{C}^\infty_{N, n}$ 的情形相同。茎上的映射
$f^\sharp : \mathcal{C}^\infty_{N, n} \to \mathcal{C}^\infty_{M, m}$
把极大理想映入极大理想；这直接源于 $f(m) = n$。

\medskip\noindent
代数几何研究概形。概形上的环层并不由其底空间的某种内蕴性质决定。
给环 $R$ 的谱（见《代数》第 \ref{algebra-section-spectrum-ring} 节）
配上由 $R$ 构造出的环层（见下文），便得到我们的基本构件。
稍后将看到，$\mathcal{O}$ 在 $\Spec(R)$ 上的各个茎正是 $R$ 在相应素理想处的局部环。
在这一背景下引入局部环层空间有两个原因：
(1) 一般而言，没有一种机制能由谱之间的连续映射自动给出相应环之间的映射，
因此我们把 $f^\sharp$ 作为附加数据；
(2) 若在环层空间范畴中考虑这些谱之间的态射，则其茎上的映射未必是局部同态，
而几何直观要求它们应当如此。为此作如下定义。

\begin{definition}
\label{schemes-definition-locally-ringed-space}
局部环层空间。
\begin{enumerate}
\item {\it 局部环层空间 $(X, \mathcal{O}_X)$} 是一个二元组，其中
$X$ 是拓扑空间，$\mathcal{O}_X$ 是环层，并且它的每个茎都是局部环。
\item 给定局部环层空间 $(X, \mathcal{O}_X)$，称
$\mathcal{O}_{X, x}$ 为 {\it $X$ 在 $x$ 处的局部环}。
以 $\mathfrak{m}_{X, x}$（或简记为 $\mathfrak{m}_x$）表示
$\mathcal{O}_{X, x}$ 的极大理想。此外，{\it $X$ 在 $x$ 处的剩余域} 是
$\kappa(x) = \mathcal{O}_{X, x}/\mathfrak{m}_x$。
\item {\it 局部环层空间之间的态射}
$(f, f^\sharp) : (X, \mathcal{O}_X) \to (Y, \mathcal{O}_Y)$
是满足下述条件的环层空间态射：对每个 $x \in X$，诱导的环同态
$\mathcal{O}_{Y, f(x)} \to \mathcal{O}_{X, x}$ 都是局部环同态。
\end{enumerate}
\end{definition}

\noindent
讨论局部环层空间时，我们通常在记号中略去环层 $\mathcal{O}_X$，
而只说“局部环层空间 $X$”。我们还滥用记号，以 $X$ 同时表示其底拓扑空间。
相应的环层 $\mathcal{O}_X$ 称为 {\it $X$ 的结构层}。
此外，局部环 $\mathcal{O}_{X, x}$ 的极大理想通常记作
$\mathfrak{m}_{X, x}$，或简记作 $\mathfrak{m}_x$。
我们会说“设 $f : X \to Y$ 是局部环层空间之间的态射”，
从而也略去结构层。此时又滥用记号，以 $f : X\to Y$ 同时表示底拓扑空间之间的连续映射。
与 $f$ 对应的 $f$-映射通常记作 $f^\sharp$。
$f$ 是局部环层空间态射这一条件，可以表述为：对每个 $x\in X$，茎上的映射
$$
f^\sharp_x : \mathcal{O}_{Y, f(x)} \longrightarrow \mathcal{O}_{X, x}
$$
把极大理想 $\mathfrak m_{Y, f(x)}$ 映入 $\mathfrak m_{X, x}$。

\medskip\noindent
下面利用这些记号约定说明，局部环层空间及其态射构成一个范畴。
为此只须验证，两个局部环层空间态射的复合仍是局部环层空间态射。
设 $f : X \to Y$ 和 $g : Y \to Z$ 是局部环层空间态射。
$f$ 与 $g$ 的复合定义于《层》的定义
\ref{sheaves-definition-composition-maps-ringed-spaces}。
取 $x \in X$。由《层》的引理 \ref{sheaves-lemma-compose-f-maps-stalks}，复合映射
$$
\mathcal{O}_{Z, g(f(x))}
\xrightarrow{g^\sharp}
\mathcal{O}_{Y, f(x)}
\xrightarrow{f^\sharp}
\mathcal{O}_{X, x}
$$
就是态射 $g \circ f$ 所诱导的茎上映射。由于局部环同态的复合仍是局部环同态，
结论随即成立。

\medskip\noindent
这个定义有一个令人满意的性质：函子
$$
\textit{局部环化空间}
\longrightarrow
\textit{环化空间}
$$
反映同构（事实上还不止于此）。下面给出一个较为具体的表述。

\begin{lemma}
\label{schemes-lemma-isomorphism-locally-ringed}
\begin{slogan}
局部环层空间之间的环层空间同构也是局部环层空间同构。
\end{slogan}
设 $X$、$Y$ 为局部环层空间。若 $f : X \to Y$ 是环层空间同构，
则 $f$ 是局部环层空间同构。
\end{lemma}

\begin{proof}
这直接归结为代数中的相应事实：若 $A$、$B$ 是局部环，
则任意环同构 $A \to B$ 都是局部环同构。
\end{proof}













\section{局部环层空间的开浸入}
\label{schemes-section-open-immersion}

\begin{definition}
\label{schemes-definition-immersion-locally-ringed-spaces}
设 $f : X \to Y$ 是局部环层空间态射。称 $f$ 为 {\it 开浸入}，如果
$f$ 将 $X$ 同胚地映到 $Y$ 的一个开子集上，并且映射
$f^{-1}\mathcal{O}_Y \to \mathcal{O}_X$ 是同构。
\end{definition}

\noindent
下述构造与《层》的定义 \ref{sheaves-definition-restriction} (3) 平行。

\begin{example}
\label{schemes-example-open-subspace}
设 $X$ 为局部环层空间，$U \subset X$ 为开子集，并令
$\mathcal{O}_U = \mathcal{O}_X|_U$ 为 $\mathcal{O}_X$ 在 $U$ 上的限制。
对 $u \in U$，茎 $\mathcal{O}_{U, u}$ 等于茎 $\mathcal{O}_{X, u}$，因而是局部环。
所以 $(U, \mathcal{O}_U)$ 是局部环层空间，且态射
$j : (U, \mathcal{O}_U) \to (X, \mathcal{O}_X)$ 是开浸入。
\end{example}

\begin{definition}
\label{schemes-definition-open-subspace}
设 $X$ 为局部环层空间，$U \subset X$ 为开子集。
上面例 \ref{schemes-example-open-subspace} 中的局部环层空间
$(U, \mathcal{O}_U)$ 称为 {\it $X$ 的与 $U$ 相伴的开子空间}。
\end{definition}

\begin{lemma}
\label{schemes-lemma-open-immersion}
设 $f : X \to Y$ 是局部环层空间的开浸入，并令
$j : V = f(X) \to Y$ 为 $Y$ 的与 $f$ 的像相伴的开子空间。
存在唯一的局部环层空间同构 $f' : X \cong V$，使得 $f = j \circ f'$。
\end{lemma}

\begin{proof}
令 $f'$ 为 $X$ 与 $V$ 之间由 $f$ 诱导的同胚。作为拓扑空间之间的映射，
$f = j \circ f'$。由于存在层同构
$f^\sharp : f^{-1}(\mathcal{O}_Y) \to \mathcal{O}_X$，故有环同构
$f^\sharp : \Gamma(U, f^{-1}(\mathcal{O}_Y)) \to \Gamma(U, \mathcal{O}_X)$，
其中 $U \subset X$ 为任意开子集
。由于 $\mathcal{O}_V = j^{-1}\mathcal{O}_Y$ 且 $f^{-1} = f'^{-1} j^{-1}$
（《层》的引理 \ref{sheaves-lemma-pullback-composition}），可得
$f^{-1}\mathcal{O}_Y = f'^{-1}\mathcal{O}_V$，从而映射
$\Gamma(U, f'^{-1}(\mathcal{O}_V)) \to \Gamma(U, f^{-1}(\mathcal{O}_Y))$
对每个开集 $U \subset X$ 都是同构。将这些同构复合，便对每个开子集
$U \subset X$ 得到环同构
$$
\Gamma(U, f'^{-1}(\mathcal{O}_V)) \to \Gamma(U, \mathcal{O}_X)
$$
，因而得到层同构 $f^{-1}(\mathcal{O}_V) \to \mathcal{O}_X$。
换言之，我们得到同构 $f'^{\sharp} : f'^{-1}(\mathcal{O}_V) \to \mathcal{O}_X$，
从而得到局部环层空间同构
$(f', f'^{\sharp}) : (X, \mathcal{O}_X) \to (V, \mathcal{O}_V)$
（使用引理 \ref{schemes-lemma-isomorphism-locally-ringed}）。按构造，
作为局部环层空间态射也有 $f = j \circ f'$。

\medskip\noindent
再设有态射 $f'' : (X, \mathcal{O}_X) \to (V, \mathcal{O}_V)$，满足
$f = j \circ f''$。对任意 $x \in X$，有 $j(f'(x)) = j(f''(x))$，由此推出
$f'(x) = f''(x)$，因为 $j$ 是包含映射。因此 $f'$ 与 $f''$ 作为拓扑空间态射相同。
在结构层上，对每个开子集 $U \subset X$，有交换图
$$
\xymatrix @R=5em{
\Gamma(U, f^{-1}(\mathcal{O}_Y)) \ar[d]_\cong\ar[r]^\cong &
\Gamma(U, \mathcal{O}_X) \\
\Gamma(U, f'^{-1}(\mathcal{O}_V)) \ar@/^/[ru]^{f'^\sharp}
\ar@/_/[ru]_{f''^\sharp} &
}
$$
由此可见，$f'^\sharp$ 与 $f''^\sharp$ 定义同一个层态射。
\end{proof}

\noindent
从现在起，我们不再区分开子集与其相伴的开子空间。

\begin{lemma}
\label{schemes-lemma-restrict-map-to-opens}
设 $f : X \to Y$ 是局部环层空间态射，$U \subset X$、$V \subset Y$ 为开子集，
并设 $f(U) \subset V$。存在唯一的局部环层空间态射 $f|_U : U \to V$，
使下图成为局部环层空间的交换方块：
$$
\xymatrix{
U \ar[d]_{f|_U} \ar[r] & X \ar[d]^f \\
V \ar[r] & Y
}
$$
\end{lemma}

\begin{proof}
略。
\end{proof}

\noindent
下文将不加说明地使用由上引理得到的如下事实。给定任意局部环层空间态射
$f : Y \to X$，以及开子集 $U \subset X$，并设它满足 $f(Y) \subset U$，
存在唯一的局部环层空间态射 $Y \to U$，使复合 $Y \to U \to X$ 等于 $f$。
事实上，我们甚至会滥用记号直接写成 $f : Y \to U$，因为这很少引起混淆。









\section{局部环层空间的闭浸入}
\label{schemes-section-closed-immersion}

\noindent
沿用《模》的定义 \ref{modules-definition-closed-immersion} 中引入的约定。

\begin{definition}
\label{schemes-definition-closed-immersion-locally-ringed-spaces}
设 $i : Z \to X$ 是局部环层空间态射。若满足下列条件，则称 $i$ 为
{\it 闭浸入}：
\begin{enumerate}
\item 映射 $i$ 将 $Z$ 同胚地映到 $X$ 的一个闭子集上。
\item 映射 $\mathcal{O}_X \to i_*\mathcal{O}_Z$ 是满射；以
$\mathcal{I}$ 表示其核。
\item $\mathcal{O}_X$-模 $\mathcal{I}$ 局部地由截面生成。
\end{enumerate}
\end{definition}

\begin{lemma}
\label{schemes-lemma-closed-local-target}
设 $f : Z \to X$ 是局部环层空间态射。要使 $f$ 为闭浸入，只需存在开覆盖
$X = \bigcup U_i$，使每个 $f : f^{-1}U_i \to U_i$ 都是闭浸入。
\end{lemma}

\begin{proof}
略。
\end{proof}

\begin{example}
\label{schemes-example-closed-subspace}
设 $X$ 是局部环层空间，$\mathcal{I} \subset \mathcal{O}_X$ 是理想层，
并且作为 $\mathcal{O}_X$-模层局部地由截面生成。令 $Z$ 为环层
$\mathcal{O}_X/\mathcal{I}$ 的支撑。由《模》的引理
\ref{modules-lemma-support-sheaf-rings-closed}，它是 $X$ 的闭子集。
以 $i : Z \to X$ 表示包含映射。由《模》的引理
\ref{modules-lemma-i-star-exact}，存在唯一的环层 $\mathcal{O}_Z$ 定义在 $Z$ 上，
满足 $i_*\mathcal{O}_Z = \mathcal{O}_X/\mathcal{I}$。
对任意 $z \in Z$，茎 $\mathcal{O}_{Z, z}$ 等于局部环的非零商
$\mathcal{O}_{X, i(z)}/\mathcal{I}_{i(z)}$，因而仍是局部环。
所以 $i : (Z, \mathcal{O}_Z) \to (X, \mathcal{O}_X)$ 是局部环层空间的闭浸入。
\end{example}

\begin{definition}
\label{schemes-definition-closed-subspace}
设 $X$ 是局部环层空间，$\mathcal{I}$ 是 $X$ 上局部地由截面生成的理想层。
上面例 \ref{schemes-example-closed-subspace} 中的局部环层空间
$(Z, \mathcal{O}_Z)$ 称为 {\it $X$ 的与理想层 $\mathcal{I}$ 相伴的闭子空间}。
\end{definition}

\begin{lemma}
\label{schemes-lemma-closed-immersion}
设 $f : X \to Y$ 是局部环层空间的闭浸入，$\mathcal{I}$ 是映射
$\mathcal{O}_Y \to f_*\mathcal{O}_X$ 的核，并令 $i : Z \to Y$ 为
$Y$ 的与 $\mathcal{I}$ 相伴的闭子空间。存在唯一的局部环层空间同构
$f' : X \cong Z$，使得 $f = i \circ f'$。
\end{lemma}

\begin{proof}
略。
\end{proof}

\begin{lemma}
\label{schemes-lemma-characterize-closed-subspace}
设 $X$、$Y$ 为局部环层空间，$\mathcal{I} \subset \mathcal{O}_X$ 是局部地由截面生成的理想层，
且 $i : Z \to X$ 是相伴的闭子空间。态射 $f : Y \to X$ 经 $Z$ 分解，
当且仅当映射 $f^*\mathcal{I} \to f^*\mathcal{O}_X = \mathcal{O}_Y$ 为零。
在这种情形下，态射 $g : Y \to Z$ 若满足 $f = i \circ g$，则它是唯一的。
\end{lemma}

\begin{proof}
显然，若 $f$ 分解为 $Y \to Z \to X$，则映射
$f^*\mathcal{I} \to \mathcal{O}_Y$ 为零。反过来，设
$f^*\mathcal{I} \to \mathcal{O}_Y$ 为零。任取 $y \in Y$，考虑环同态
$f^\sharp_y : \mathcal{O}_{X, f(y)} \to \mathcal{O}_{Y, y}$。
由于复合映射
$\mathcal{I}_{f(y)} \to \mathcal{O}_{X, f(y)} \to \mathcal{O}_{Y, y}$
按假设为零，且 $f^\sharp_y(1) = 1$，可知
$1 \not \in \mathcal{I}_{f(y)}$，即
$\mathcal{I}_{f(y)} \not = \mathcal{O}_{X, f(y)}$。因此
$f(Y) \subset Z = \text{Supp}(\mathcal{O}_X/\mathcal{I})$。
于是 $f = i \circ g$，其中 $g : Y \to Z$ 连续。考虑映射
$f^\sharp : \mathcal{O}_X \to f_*\mathcal{O}_Y$。
假设 $f^*\mathcal{I} \to \mathcal{O}_Y$ 为零，意味着复合
$\mathcal{I} \to \mathcal{O}_X \to f_*\mathcal{O}_Y$ 为零；
这是由 $f_*$ 与 $f^*$ 的伴随性得到的。
换言之，我们得到环层态射
$\overline{f^\sharp} : \mathcal{O}_X/\mathcal{I} \to f_*\mathcal{O}_Y$。
注意 $f_*\mathcal{O}_Y = i_*g_*\mathcal{O}_Y$，且
$\mathcal{O}_X/\mathcal{I} = i_*\mathcal{O}_Z$。
由《层》的引理 \ref{sheaves-lemma-equivalence-categories-closed-structures}，
得到唯一的环层态射 $g^\sharp : \mathcal{O}_Z \to g_*\mathcal{O}_Y$，
其沿 $i$ 的正像为 $\overline{f^\sharp}$。我们略去如下验证：
$(g, g^\sharp)$ 定义局部环层空间态射，并且作为局部环层空间态射有
$f = i \circ g$。$(g, g^\sharp)$ 的唯一性已在上面指出。
\end{proof}

\begin{lemma}
\label{schemes-lemma-restrict-map-to-closed}
设 $f : X \to Y$ 是局部环层空间态射，
$\mathcal{I} \subset \mathcal{O}_Y$ 是局部地由截面生成的理想层，
且 $i : Z \to Y$ 是与 $\mathcal{I}$ 相伴的闭子空间。
令 $\mathcal{J}$ 为映射
$f^*\mathcal{I} \to f^*\mathcal{O}_Y = \mathcal{O}_X$ 的像。
则这个理想局部地由截面生成。再令 $i' : Z' \to X$ 为相伴的 $X$ 的闭子空间。
存在唯一的局部环层空间态射 $f' : Z' \to Z$，使下图成为局部环层空间的交换方块：
$$
\xymatrix{
Z' \ar[d]_{f'} \ar[r]_{i'} & X \ar[d]^f \\
Z \ar[r]^{i} & Y
}
$$
此外，这个图在局部环层空间范畴中是纤维方块。
\end{lemma}

\begin{proof}
由《模》的引理 \ref{modules-lemma-pullback-locally-generated}，
理想 $\mathcal{J}$ 局部地由截面生成。引理的其余部分由上面的引理
\ref{schemes-lemma-characterize-closed-subspace} 对“态射经闭子空间分解”的刻画得到。
\end{proof}














\section{仿射概形}
\label{schemes-section-affine-schemes}

\noindent
设 $R$ 为环。考虑拓扑空间 $\Spec(R)$，它与 $R$ 相伴；参见《代数》第
\ref{algebra-section-spectrum-ring} 节。我们将在这个空间上赋予一个环层
$\mathcal{O}_{\Spec(R)}$，所得二元组
$(\Spec(R), \mathcal{O}_{\Spec(R)})$ 将是一个仿射概形。

\medskip\noindent
回顾 $\Spec(R)$ 有一个由开集 $D(f)$（$f \in R$）组成的基；这些开集称为
标准开集，参见《代数》的定义 \ref{algebra-definition-Zariski-topology}。
此外，两个标准开集的交仍为标准开集：
$D(f) \cap D(g) = D(fg)$，其中 $f, g\in R$。

\begin{lemma}
\label{schemes-lemma-standard-open}
设 $R$ 为环，$f \in R$。
\begin{enumerate}
\item 若 $g\in R$ 且 $D(g) \subset D(f)$，则
\begin{enumerate}
\item $f$ 在 $R_g$ 中可逆；
\item $g^e = af$，其中 $e \geq 1$ 且 $a \in R$；
\item 存在典范环同态 $R_f \to R_g$；
\item 存在典范的 $R_f$-模同态 $M_f \to M_g$，适用于任意 $R$-模 $M$。
\end{enumerate}
\item $D(f)$ 的任意开覆盖都可加细为形如
$D(f) = \bigcup_{i = 1}^n D(g_i)$ 的有限开覆盖。
\item 若 $g_1, \ldots, g_n \in R$，则 $D(f) \subset \bigcup D(g_i)$
当且仅当 $g_1, \ldots, g_n$ 在 $R_f$ 中生成单位理想。
\end{enumerate}
\end{lemma}

\begin{proof}
回顾 $D(g) = \Spec(R_g)$（见《代数》的引理
\ref{algebra-lemma-standard-open}）。(a) 成立，是因为 $f$ 在 $R_g$ 中的像
不属于任何素理想，因而可逆；参见《代数》的引理
\ref{algebra-lemma-Zariski-topology}。把 $f$ 在 $R_g$ 中的逆元写成
$a/g^d$。这意味着 $g^d - af$ 被 $g$ 的某个幂零化，从而得到 (b)。
对于 (c)，由 (a) 及局部化的泛性质，映射 $R_f \to R_g$ 存在；
也可以把 $b/f^n$ 映到 $a^nb/g^{ne}$ 来定义它。利用等式
$M_f = M \otimes_R R_f$ 可以得到模之间的映射 $M_f \to M_g$；
也可以把 $x/f^n$ 映到 $a^nx/g^{ne}$ 来定义它。

\medskip\noindent
回顾 $D(f)$ 是拟紧的，参见《代数》的引理
\ref{algebra-lemma-qc-open}。由于标准开集构成这个拓扑的一组基，
第二个断言立即成立。

\medskip\noindent
第三个断言直接来自《代数》的引理
\ref{algebra-lemma-Zariski-topology}。
\end{proof}

\noindent
在《层》第 \ref{sheaves-section-bases} 节中，我们定义了基上的层，
并说明这一概念与空间上的层本质上等价；参见《层》的引理
\ref{sheaves-lemma-extend-off-basis} 和
\ref{sheaves-lemma-extend-off-basis-structures}。此外，《层》的引理
\ref{sheaves-lemma-cofinal-systems-coverings-standard-case} 表明，
对每个标准开集，只须在一个开覆盖的共尾系上检验层条件。
由上面的引理，只须对标准开集构成的有限覆盖检验这一条件。

\begin{definition}
\label{schemes-definition-standard-covering}
设 $R$ 为环。
\begin{enumerate}
\item $\Spec(R)$ 的一个 {\it 标准开覆盖} 是形如
$\Spec(R) = \bigcup_{i = 1}^n D(f_i)$ 的覆盖，其中
$f_1, \ldots, f_n \in R$。
\item 设 $D(f) \subset \Spec(R)$ 是标准开集。$D(f)$ 的一个
{\it 标准开覆盖} 是形如 $D(f) = \bigcup_{i = 1}^n D(g_i)$ 的覆盖，
其中 $g_1, \ldots, g_n \in R$。
\end{enumerate}
\end{definition}

\noindent
设 $R$ 为环，$M$ 为 $R$-模。我们将在标准开集构成的基上定义预层
$\widetilde M$。设 $U \subset \Spec(R)$ 是标准开集。若 $f, g \in R$
满足 $D(f) = D(g)$，则由上面的引理 \ref{schemes-lemma-standard-open}，
存在互为逆映射的典范映射 $M_f \to M_g$ 和 $M_g \to M_f$。
因此可以任选 $f$，使其满足 $U = D(f)$，并定义
$$
\widetilde M(U) = M_f.
$$
注意，若 $D(g) \subset D(f)$，则由上面的引理
\ref{schemes-lemma-standard-open}，存在典范映射
$$
\widetilde M(D(f)) = M_f \longrightarrow M_g = \widetilde M(D(g)).
$$
显然，这在标准开集构成的基上定义了一个阿贝尔群预层。若 $M = R$，
则 $\widetilde R$ 是该基上的环预层。

\medskip\noindent
下面计算 $\widetilde M$ 在点 $x \in \Spec(R)$ 处的茎。设 $x$ 对应于素理想
$\mathfrak p \subset R$。由茎的定义可得
$$
\widetilde M_x = \colim_{f\in R, f\not\in \mathfrak p} M_f
$$
这里，集合 $\{f \in R, f \not \in \mathfrak p\}$ 由规则
$f \geq f' \Leftrightarrow D(f) \subset D(f')$ 赋予预序。
若 $f_1, f_2 \in R \setminus \mathfrak p$，则在这个预序中
$f_1f_2 \geq f_1$。因此，由《代数》的引理
\ref{algebra-lemma-localization-colimit} 可得
$$
\widetilde M_x = M_{\mathfrak p}.
$$

\medskip\noindent
接下来检验标准开覆盖的层条件。若
$D(f) = \bigcup_{i = 1}^n D(g_i)$，则这个覆盖的层条件等价于下列序列的正合性：
$$
0 \to M_f \to \bigoplus M_{g_i} \to \bigoplus M_{g_ig_j}.
$$
注意 $D(g_i) = D(fg_i)$，因此可将这个序列改写为
$$
0 \to M_f \to \bigoplus M_{fg_i} \to \bigoplus M_{fg_ig_j}.
$$
此外，由上面的引理 \ref{schemes-lemma-standard-open}，
$g_1, \ldots, g_n$ 在 $R_f$ 中生成单位理想。因此可把《代数》的引理
\ref{algebra-lemma-cover-module} 应用于模 $M_f$、环 $R_f$ 以及元素
$g_1, \ldots, g_n$。由此可知该序列正合。根据前面的说明，
$\widetilde M$ 是标准开集构成的基上的一个层。

\medskip\noindent
由《层》第 \ref{sheaves-section-bases} 节的内容可知，存在唯一的环层
$\mathcal{O}_{\Spec(R)}$，它在标准开集上与 $\widetilde R$ 一致。
注意，由上面对茎的计算，这个环层的所有茎都是局部环。

\medskip\noindent
类似地，对任意 $R$-模 $M$，存在唯一的 $\mathcal{O}_{\Spec(R)}$-模层
$\mathcal{F}$，它在标准开集上与 $\widetilde M$ 一致；参见《层》的引理
\ref{sheaves-lemma-extend-off-basis-module}。

\begin{definition}
\label{schemes-definition-structure-sheaf}
设 $R$ 为环。
\begin{enumerate}
\item {\it 结构层 $\mathcal{O}_{\Spec(R)}$，即 $R$ 的谱的结构层}，是唯一的环层
$\mathcal{O}_{\Spec(R)}$，它在标准开集构成的基上与 $\widetilde R$ 一致。
\item 局部环层空间 $(\Spec(R), \mathcal{O}_{\Spec(R)})$ 称为 $R$ 的
{\it 谱}，记作 $\Spec(R)$。
\item 这个 $\mathcal{O}_{\Spec(R)}$-模层把 $\widetilde M$ 扩张到
$\Spec(R)$ 的全部开集上；它称为一个 $\mathcal{O}_{\Spec(R)}$-模层，
且与 $M$ 相伴。这个层仍记作 $\widetilde M$。
\end{enumerate}
\end{definition}

\noindent
下面汇总至此所得的结果。

\begin{lemma}
\label{schemes-lemma-spec-sheaves}
设 $R$ 为环，$M$ 为 $R$-模，并设 $\widetilde M$ 是一个
$\mathcal{O}_{\Spec(R)}$-模层，它与 $M$ 相伴。
\begin{enumerate}
\item 有 $\Gamma(\Spec(R), \mathcal{O}_{\Spec(R)}) = R$。
\item 有 $\Gamma(\Spec(R), \widetilde M) = M$，且这是作为 $R$-模的等式。
\item 对每个 $f \in R$，有
$\Gamma(D(f), \mathcal{O}_{\Spec(R)}) = R_f$。
\item 对每个 $f\in R$，有 $\Gamma(D(f), \widetilde M) = M_f$，
且这是作为 $R_f$-模的等式。
\item 只要 $D(g) \subset D(f)$，$\mathcal{O}_{\Spec(R)}$ 和
$\widetilde M$ 上的限制映射就是引理 \ref{schemes-lemma-standard-open} 中的映射
$R_f \to R_g$ 和 $M_f \to M_g$。
\item 设 $\mathfrak p$ 是 $R$ 的素理想，$x \in \Spec(R)$ 是相应的点。
则有 $\mathcal{O}_{\Spec(R), x} = R_{\mathfrak p}$。
\item 设 $\mathfrak p$ 是 $R$ 的素理想，$x \in \Spec(R)$ 是相应的点。
有 $\widetilde M_x = M_{\mathfrak p}$，且这是作为 $R_{\mathfrak p}$-模的等式。
\end{enumerate}
此外，所有这些识别都关于 $R$-模 $M$ 具有函子性。特别地，函子
$M \mapsto \widetilde M$ 是从 $R$-模范畴到
$\mathcal{O}_{\Spec(R)}$-模范畴的正合函子。
\end{lemma}

\begin{proof}
断言 (1)--(7) 由上面的讨论显然成立。函子 $M \mapsto \widetilde M$
的正合性来自以下两个事实：函子 $M \mapsto M_{\mathfrak p}$ 是正合的，
而短正合列的正合性可以在茎上检验；参见《模》的引理
\ref{modules-lemma-abelian}。
\end{proof}

\begin{definition}
\label{schemes-definition-affine-scheme}
{\it 仿射概形} 是这样的局部环层空间：它作为局部环层空间同构于
$\Spec(R)$，其中 $R$ 是某个环。{\it 仿射概形之间的态射} 是局部环层空间范畴中的态射。
\end{definition}

\noindent
仿射概形在所有局部环层空间中起着特殊作用；下一节将讨论这一点。


















\section{仿射概形范畴}
\label{schemes-section-category-affine-schemes}

\noindent
注意，若 $Y$ 是仿射概形，则存在典范的 $1-1$ 双射，
把它的点对应到 $\Gamma(Y, \mathcal{O}_Y)$ 中的素理想。

\begin{lemma}
\label{schemes-lemma-morphism-into-affine-where-point-goes}
设 $X$ 为局部环层空间，$Y$ 为仿射概形，并设
$f \in \Mor(X, Y)$ 是局部环层空间态射。给定点 $x \in X$，考虑环同态
$$
\Gamma(Y, \mathcal{O}_Y) \xrightarrow{f^\sharp}
\Gamma(X, \mathcal{O}_X) \to \mathcal{O}_{X, x}
$$
以 $\mathfrak p \subset \Gamma(Y, \mathcal{O}_Y)$ 表示 $\mathfrak m_x$ 的原像。
设 $y \in Y$ 是相应的点，则 $f(x) = y$。
\end{lemma}

\begin{proof}
考虑交换图
$$
\xymatrix{
\Gamma(X, \mathcal{O}_X) \ar[r] &
\mathcal{O}_{X, x} \\
\Gamma(Y, \mathcal{O}_Y) \ar[r] \ar[u] &
\mathcal{O}_{Y, f(x)} \ar[u]
}
$$
（参见《层》的定义 \ref{sheaves-definition-f-map} 下方关于 $f$-映射的讨论）。
由于右侧竖直箭头是局部同态，可知 $\mathfrak m_{f(x)}$ 是
$\mathfrak m_x$ 的原像。结论随即成立。
\end{proof}

\begin{lemma}
\label{schemes-lemma-f-open}
设 $X$ 为局部环层空间，$f \in \Gamma(X, \mathcal{O}_X)$。集合
$$
D(f) = \{x \in X \mid \text{像 }f \not\in \mathfrak m_x\}
$$
是开集。此外，$f|_{D(f)}$ 有逆元。
\end{lemma}

\begin{proof}
这是《模》的引理 \ref{modules-lemma-s-open} 的特殊情形；这里也给出直接证明。
设 $U \subset X$ 和 $V \subset X$ 是两个开子集，且 $f|_U$ 有逆元 $g$，
$f|_V$ 有逆元 $h$。显然 $g|_{U\cap V} = h|_{U\cap V}$。
因此，只须说明 $f$ 在任意 $x \in D(f)$ 的某个开邻域上可逆。
这是显然的，因为 $f \not \in \mathfrak m_x$ 意味着
$f \in \mathcal{O}_{X, x}$ 有逆元 $g \in \mathcal{O}_{X, x}$；
这又意味着存在某个开邻域 $x \in U \subset X$，使得
$g \in \mathcal{O}_X(U)$ 且 $g\cdot f|_U = 1$。
\end{proof}

\begin{lemma}
\label{schemes-lemma-f-open-affine}
在上面的引理 \ref{schemes-lemma-f-open} 中，若 $X$ 是仿射概形，
则开集 $D(f)$ 与前面定义的标准开集 $D(f)$ 相同
（见《代数》的定义 \ref{algebra-definition-spectrum-ring}）。
\end{lemma}

\begin{proof}
略。
\end{proof}

\begin{lemma}
\label{schemes-lemma-morphism-into-affine}
\begin{reference}
这一事实可参见 \cite[II, Err 1, Prop. 1.8.1]{EGA}，其中将其归于 J. Tate。
\end{reference}
设 $X$ 为局部环层空间，$Y$ 为仿射概形。映射
$$
\Mor(X, Y)
\longrightarrow
\Hom(\Gamma(Y, \mathcal{O}_Y), \Gamma(X, \mathcal{O}_X))
$$
把 $f$ 映到全局截面上的 $f^\sharp$；这个映射是双射。
\end{lemma}

\begin{proof}
由于 $Y$ 是仿射概形，有
$(Y, \mathcal{O}_Y) \cong (\Spec(R), \mathcal{O}_{\Spec(R)})$
，其中 $R$ 是某个环。在证明中，我们将使用关于 $Y$ 及其结构层的一些事实；它们直接来自已知的环谱性质，
例如引理 \ref{schemes-lemma-spec-sheaves}。

\medskip\noindent
受上述引理启发，我们构造逆映射。设
$\psi_Y : \Gamma(Y, \mathcal{O}_Y) \to \Gamma(X, \mathcal{O}_X)$
是环同态。首先定义相应的空间映射
$$
\Psi : X \longrightarrow Y
$$
，其规则由引理 \ref{schemes-lemma-morphism-into-affine-where-point-goes} 给出。
换言之，给定 $x \in X$，定义 $\Psi(x)$ 为 $Y$ 中与下述素理想相应的点：
该素理想位于 $\Gamma(Y, \mathcal{O}_Y)$ 中，并且是 $\mathfrak m_x$
在如下复合映射下的原像：
$
\Gamma(Y, \mathcal{O}_Y) \xrightarrow{\psi_Y}
\Gamma(X, \mathcal{O}_X) \to
\mathcal{O}_{X, x}
$.

\medskip\noindent
我们断言映射 $\Psi : X \to Y$ 连续。标准开集 $D(g)$（其中
$g \in \Gamma(Y, \mathcal{O}_Y)$）构成 $Y$ 的拓扑的一组基。
因此只须证明 $\Psi^{-1}(D(g))$ 是开集。按 $\Psi$ 的构造，原像
$\Psi^{-1}(D(g))$ 恰为集合 $D(\psi_Y(g)) \subset X$；
由引理 \ref{schemes-lemma-f-open}，它是开集。因此 $\Psi$ 连续。

\medskip\noindent
接着构造一个层的 $\Psi$-映射，它从 $\mathcal{O}_Y$ 到 $\mathcal{O}_X$。
由《层》的引理 \ref{sheaves-lemma-f-map-basis-below-structures}，
只须定义与限制映射相容的环同态
$\psi_{D(g)} : \Gamma(D(g), \mathcal{O}_Y) \to
\Gamma(\Psi^{-1}(D(g)), \mathcal{O}_X)$。
存在典范同构
$\Gamma(D(g), \mathcal{O}_Y) = \Gamma(Y, \mathcal{O}_Y)_g$，
因为 $Y$ 是仿射概形。
由于 $\psi_Y(g)$ 在 $D(\psi_Y(g))$ 上可逆，存在典范映射
$$
\Gamma(Y, \mathcal{O}_Y)_g
\longrightarrow
\Gamma(\Psi^{-1}(D(g)), \mathcal{O}_X)
=
\Gamma(D(\psi_Y(g)), \mathcal{O}_X)
$$
；由局部化的泛性质，它扩张映射 $\psi_Y$。注意，这里只能取这个典范映射！
将它与典范识别
$\Gamma(D(g), \mathcal{O}_Y) = \Gamma(Y, \mathcal{O}_Y)_g$ 结合，
定义为 $\psi_{D(g)}$。这与局部化相容，因为仿射概形上的限制映射同样由
局部化的泛性质定义；参见引理 \ref{schemes-lemma-spec-sheaves} 和
\ref{schemes-lemma-standard-open}。

\medskip\noindent
这样就定义了环层空间态射
$(\Psi, \psi) : (X, \mathcal{O}_X) \to (Y, \mathcal{O}_Y)$
，它在全局截面上给回 $\psi_Y$。要说明它是局部环层空间态射，
必须证明所诱导的局部环同态
$$
\psi_x : \mathcal{O}_{Y, \Psi(x)} \longrightarrow \mathcal{O}_{X, x}
$$
都是局部的。这立即来自引理
\ref{schemes-lemma-morphism-into-affine-where-point-goes} 的证明中的交换图以及
$\Psi$ 的定义。

\medskip\noindent
最后要说明，构造 $(\Psi, \psi) \mapsto \psi_Y$ 与构造
$\psi_Y \mapsto (\Psi, \psi)$ 互为逆构造。显然，
$\psi_Y \mapsto (\Psi, \psi) \mapsto \psi_Y$。因此只须证明：给定
$\psi_Y$，至多有一个二元组 $(\Psi, \psi)$ 产生它。
$\Psi$ 的唯一性已在引理 \ref{schemes-lemma-morphism-into-affine-where-point-goes}
中证明；给定 $\Psi$ 的唯一性后，映射 $\psi$ 的唯一性也已在上面的证明过程中指出。
\end{proof}

\begin{lemma}
\label{schemes-lemma-category-affine-schemes}
仿射概形范畴与环范畴的反范畴等价。这个等价由如下函子给出：
它把仿射概形映到其结构层的全局截面。
\end{lemma}

\begin{proof}
这由定义 \ref{schemes-definition-affine-scheme} 和引理
\ref{schemes-lemma-morphism-into-affine} 立即得到。
\end{proof}

\begin{lemma}
\label{schemes-lemma-standard-open-affine}
设 $Y$ 为仿射概形，$f \in \Gamma(Y, \mathcal{O}_Y)$。
开子空间 $D(f)$ 是仿射概形。
\end{lemma}

\begin{proof}
可以假设 $Y = \Spec(R)$ 且 $f \in R$。考虑仿射概形态射
$\phi : U = \Spec(R_f) \to \Spec(R) = Y$，它由环同态
$R \to R_f$ 诱导。由《代数》的引理
\ref{algebra-lemma-standard-open}，它同胚地映到 $D(f)$ 上。
另一方面，映射 $\phi^{-1}\mathcal{O}_Y \to \mathcal{O}_U$ 在茎上是同构，
因而本身是同构。所以 $\phi$ 是开浸入。由引理 \ref{schemes-lemma-open-immersion}，
$D(f)$ 同构于 $U$。
\end{proof}

\begin{lemma}
\label{schemes-lemma-fibre-product-affine-schemes}
仿射概形范畴有有限积和纤维积。换言之，它有有限极限。
此外，仿射概形范畴中的积和纤维积与局部环层空间范畴中的相同。
用公式表示，在局部环层空间范畴中有
$$
\Spec(R) \times \Spec(S) =
\Spec(R \otimes_{\mathbf{Z}} S)
$$
；而给定环同态 $R \to A$、$R \to B$，有
$$
\Spec(A) \times_{\Spec(R)} \Spec(B)
=
\Spec(A \otimes_R B).
$$
\end{lemma}

\begin{proof}
这只是引理 \ref{schemes-lemma-morphism-into-affine} 的一个应用。
首先，由该引理，仿射概形 $\Spec(\mathbf{Z})$ 是局部环层空间范畴中的终对象。
因此第一个展示公式由第二个推出。为证明第二个公式，注意对任意局部环层空间 $X$，有
\begin{eqnarray*}
\Mor(X, \Spec(A \otimes_R B))
& = &
\Hom(A \otimes_R B, \mathcal{O}_X(X)) \\
& = &
\Hom(A, \mathcal{O}_X(X))
\times_{\Hom(R, \mathcal{O}_X(X))}
\Hom(B, \mathcal{O}_X(X)) \\
& = &
\Mor(X, \Spec(A))
\times_{\Mor(X, \Spec(R))}
\Mor(X, \Spec(B))
\end{eqnarray*}
这就证明了该公式。有关定义参见《范畴》第
\ref{categories-section-fibre-products} 节。
\end{proof}

\begin{lemma}
\label{schemes-lemma-disjoint-union-affines}
设 $X$ 为局部环层空间。假设 $X = U \amalg V$，其中 $U$ 和 $V$ 是开集，
且 $U$、$V$ 都是仿射概形，则 $X$ 是仿射概形。
\end{lemma}

\begin{proof}
令 $R = \Gamma(X, \mathcal{O}_X)$。由层性质，
$R = \mathcal{O}_X(U) \times \mathcal{O}_X(V)$。
由引理 \ref{schemes-lemma-morphism-into-affine}，存在典范局部环层空间态射
$X \to \Spec(R)$。由《代数》的引理 \ref{algebra-lemma-spec-product}，
作为拓扑空间有
$\Spec(\mathcal{O}_X(U)) \amalg \Spec(\mathcal{O}_X(V)) =
\Spec(R)$
，其中的映射来自环同态 $R \to \mathcal{O}_X(U)$ 和
$R \to \mathcal{O}_X(V)$。这当然也意味着 $\Spec(R)$ 是局部环层空间范畴中的余积。
按假设，态射 $X \to \Spec(R)$ 诱导 $\Spec(\mathcal{O}_X(U))$ 与 $U$ 的同构，
对 $V$ 也同样如此。因此 $X \to \Spec(R)$ 是同构。
\end{proof}















\section{仿射概形上的拟凝聚层}
\label{schemes-section-quasi-coherent-affine}

\noindent
回顾，我们已在《模》的定义 \ref{modules-definition-quasi-coherent}
中定义了拟凝聚层这一抽象概念。本节将证明，仿射概形 $\Spec(R)$
上的任意拟凝聚层都对应于层 $\widetilde M$；这里它与某个 $R$-模 $M$
相关联。

\begin{lemma}
\label{schemes-lemma-compare-constructions}
设 $(X, \mathcal{O}_X) = (\Spec(R), \mathcal{O}_{\Spec(R)})$
为仿射概形。设 $M$ 为 $R$-模。层
$\widetilde M$（定义 \ref{schemes-definition-structure-sheaf}）由这个 $R$-模
$M$ 所关联；它与层 $\mathcal{F}_M$ 之间存在典范同构，后者由 $R$-模
$M$ 所关联
（《模》的定义 \ref{modules-definition-sheaf-associated}）。
该同构关于 $M$ 具有函子性。特别地，各层 $\widetilde M$ 都是拟凝聚的。
此外，它们由如下映射性质刻画：
$$
\Hom_{\mathcal{O}_X}(\widetilde M, \mathcal{F})
=
\Hom_R(M, \Gamma(X, \mathcal{F}))
$$
其中可取任意一个 $\mathcal{O}_X$-模层 $\mathcal{F}$。这里，映射
$\alpha : \widetilde M \to \mathcal{F}$ 对应于它在整体截面上诱导的映射。
\end{lemma}

\begin{proof}
由《模》的引理 \ref{modules-lemma-construct-quasi-coherent-sheaves}，
存在一个态射 $\mathcal{F}_M \to \widetilde M$，它对应于映射
$M \to \Gamma(X, \widetilde M) = M$。设 $x \in X$ 对应于素理想
$\mathfrak p \subset R$。在茎上诱导的映射是
$\mathcal{O}_{X, x} \otimes_R M \to M_{\mathfrak p}$；它们都是同构，
因为 $R_{\mathfrak p} \otimes_R M = M_{\mathfrak p}$。
所以映射 $\mathcal{F}_M \to \widetilde M$ 是同构。所述映射性质则来自
各层 $\mathcal{F}_M$ 的映射性质。
\end{proof}

\begin{lemma}
\label{schemes-lemma-widetilde-constructions}
设 $(X, \mathcal{O}_X) = (\Spec(R), \mathcal{O}_{\Spec(R)})$
为仿射概形。存在如下典范同构：
\begin{enumerate}
\item
$
\widetilde{M \otimes_R N}
\cong
\widetilde M \otimes_{\mathcal{O}_X} \widetilde N
$,
参见《模》第 \ref{modules-section-tensor-product} 节。
\item
$
\widetilde{\text{T}^n(M)}
\cong
\text{T}^n(\widetilde M)
$,
$
\widetilde{\text{Sym}^n(M)}
\cong
\text{Sym}^n(\widetilde M)
$，以及
$
\widetilde{\wedge^n(M)}
\cong
\wedge^n(\widetilde M)
$，
参见《模》第 \ref{modules-section-symmetric-exterior} 节。
\item 若 $M$ 是有限表示的 $R$-模，则
$
\SheafHom_{\mathcal{O}_X}(\widetilde M, \widetilde N)
\cong
\widetilde{\Hom_R(M,  N)}
$，
参见《模》第 \ref{modules-section-internal-hom} 节。
\end{enumerate}
\end{lemma}

\begin{proof}[第一种证明]
由引理 \ref{schemes-lemma-compare-constructions} 和《模》的引理
\ref{modules-lemma-construct-quasi-coherent-sheaves} 可知，函子
$M \mapsto \widetilde M$ 可以看成 $\pi^*$，其中 $\pi$ 是某个环层空间态射。
再由《模》的引理 \ref{modules-lemma-tensor-product-pullback} 和
\ref{modules-lemma-pullback-tensor-algebra}，模的拉回与张量构造相交换。
例如，态射 $\pi : (X, \mathcal{O}_X) \to (\{*\}, R)$ 是平坦的，
因为 $\mathcal{O}_X$ 的各个茎都是 $R$ 的局部化（引理
\ref{schemes-lemma-spec-sheaves}），因而在 $R$ 上平坦。因此，当第一个模是
有限表示的时，由《模》的引理 \ref{modules-lemma-pullback-internal-hom}，
沿 $\pi$ 的拉回与内部 Hom 相交换。
\end{proof}

\begin{proof}[第二种证明]
(1) 的证明。由引理 \ref{schemes-lemma-compare-constructions}，要给出从
$\widetilde{M \otimes_R N}$ 到
$\widetilde M \otimes_{\mathcal{O}_X} \widetilde N$ 的映射，
只须给出整体截面上的映射
$M \otimes_R N \to
\Gamma(X, \widetilde M \otimes_{\mathcal{O}_X} \widetilde N)$
它由模层张量积的定义给出。要证明该映射是同构，只须检验它在各个茎上
都是同构。这来自 $\widetilde{M}$ 的茎的描述（见引理
\ref{schemes-lemma-spec-sheaves}，或《模》的引理
\ref{modules-lemma-construct-quasi-coherent-sheaves}）、张量积与局部化
相交换这一事实（《代数》的引理
\ref{algebra-lemma-tensor-product-localization}），以及《模》的引理
\ref{modules-lemma-stalk-tensor-product}。

\medskip\noindent
(2) 的证明。这与 (1) 的证明类似，只须使用《代数》的引理
\ref{algebra-lemma-tensor-algebra-localization} 和《模》的引理
\ref{modules-lemma-stalk-tensor-algebra}。

\medskip\noindent
(3) 的证明。由于构造 $M \mapsto \widetilde{M}$ 具有函子性，存在一个
$R$-线性映射
$\Hom_R(M, N) \to \Hom_{\mathcal{O}_X}(\widetilde{M}, \widetilde{N})$.
该映射的目标是
$\SheafHom_{\mathcal{O}_X}(\widetilde M, \widetilde N)$ 的整体截面。
所以由引理 \ref{schemes-lemma-compare-constructions}，得到一个
$\mathcal{O}_X$-模映射 $\widetilde{\Hom_R(M,  N)} \to
\SheafHom_{\mathcal{O}_X}(\widetilde M, \widetilde N)$。
比较各茎即可检验它是同构。若 $M$ 作为 $R$-模是有限表示的，则
$\widetilde M$ 作为 $\mathcal{O}_X$-模具有整体有限表示。因此结论来自
《代数》的引理 \ref{algebra-lemma-hom-from-finitely-presented} 和《模》的
引理 \ref{modules-lemma-stalk-internal-hom}。
\end{proof}

\begin{proof}[(1) 的第三种证明]
对任意 $\mathcal{O}_X$-模 $\mathcal{F}$，有如下关于 $M$、$N$ 和
$\mathcal{F}$ 具有函子性的同构：
\begin{align*}
\Hom_{\mathcal{O}_X}(\widetilde{M} \otimes _{\mathcal{O} _X} \widetilde{N},
\mathcal{F})
& =
\Hom_{\mathcal{O}_X}(\widetilde{M},
\SheafHom_{\mathcal{O} _X} (\widetilde{N}, \mathcal{F})) \\
& =
\Hom_R(M, \Gamma(X,
\SheafHom_{\mathcal{O}_X}(\widetilde{N}, \mathcal{F}))) \\
& =
\Hom_R(M, \Hom_{\mathcal{O}_X}(\widetilde{N}, \mathcal{F})) \\
& =
\Hom_R(M, \Hom_R(N, \Gamma(X,\mathcal{F}))) \\
& =
\Hom_R(M \otimes_R N, \Gamma(X, \mathcal{F})) \\
& =
\Hom_{\mathcal{O}_X}(\widetilde{M \otimes_R N}, \mathcal{F})
\end{align*}
第一个等式是《模》的引理 \ref{modules-lemma-internal-hom}。
第二个等式是 $\widetilde{M}$ 的泛性质，参见引理
\ref{schemes-lemma-compare-constructions}。第三个等式由 $\SheafHom$ 的定义成立。
第四个等式是 $\widetilde{N}$ 的泛性质。第五个等式是《代数》的引理
\ref{algebra-lemma-hom-from-tensor-product}。最后一个等式是
$\widetilde{M \otimes_R N}$ 的泛性质。由米田引理（《范畴》的引理
\ref{categories-lemma-yoneda}）即得 (1)。
\end{proof}

\begin{lemma}
\label{schemes-lemma-widetilde-pullback}
设
$(X, \mathcal{O}_X) = (\Spec(S), \mathcal{O}_{\Spec(S)})$,
$(Y, \mathcal{O}_Y) = (\Spec(R), \mathcal{O}_{\Spec(R)})$
为仿射概形。设 $\psi : (X, \mathcal{O}_X) \to (Y, \mathcal{O}_Y)$
为仿射概形的态射，对应于环同态 $\psi^\sharp : R \to S$
（见引理 \ref{schemes-lemma-category-affine-schemes}）。
\begin{enumerate}
\item 有 $\psi^* \widetilde M = \widetilde{S \otimes_R M}$，且该等式关于
$R$-模 $M$ 具有函子性。
\item 有 $\psi_* \widetilde N = \widetilde{N_R}$，且该等式关于
$S$-模 $N$ 具有函子性。
\end{enumerate}
\end{lemma}

\begin{proof}
第一个断言来自引理 \ref{schemes-lemma-compare-constructions} 中的认同以及《模》的
引理 \ref{modules-lemma-restrict-quasi-coherent}。第二个断言来自等式
$\psi^{-1}(D(f)) = D(\psi^\sharp(f))$，因而
$$
\psi_* \widetilde N(D(f)) = \widetilde N(D(\psi^\sharp(f))) =
N_{\psi^\sharp(f)} = (N_R)_f = \widetilde{N_R}(D(f))
$$
正是所需的等式。
\end{proof}

\noindent
上面的引理 \ref{schemes-lemma-widetilde-pullback} 特别说明：把层
$\widetilde M$ 限制到标准仿射开子空间 $D(f)$ 上，所得就是
$\widetilde{M_f}$。下文将直接使用这一事实，不再另行说明。

\begin{lemma}
\label{schemes-lemma-quasi-coherent-affine}
设 $(X, \mathcal{O}_X) = (\Spec(R), \mathcal{O}_{\Spec(R)})$
为仿射概形，$\mathcal{F}$ 为拟凝聚 $\mathcal{O}_X$-模。那么
$\mathcal{F}$ 同构于由 $R$-模 $\Gamma(X, \mathcal{F})$ 所关联的层。
\end{lemma}

\begin{proof}
设 $\mathcal{F}$ 为拟凝聚 $\mathcal{O}_X$-模。由于每个标准开集
$D(f)$ 都是拟紧的，可见 $X$ 是局部拟紧的；也就是说，每一点都有一个
由拟紧邻域组成的基本系，参见《拓扑》的定义
\ref{topology-definition-locally-quasi-compact}。因此，由《模》的引理
\ref{modules-lemma-quasi-coherent-module}，对每个素理想
$\mathfrak p \subset R$ 及其对应的点 $x \in X$，都存在开邻域
$x \in U \subset X$，
使得 $\mathcal{F}|_U$ 同构于由某个 $\mathcal{O}_X(U)$-模 $M$ 所关联的
拟凝聚层。换言之，我们得到一个由具有这一性质的 $U$ 组成的开覆盖。
例如，由引理 \ref{schemes-lemma-standard-open}，可把这个覆盖加细为标准开覆盖。
于是得到覆盖 $\Spec(R) = \bigcup D(f_i)$、$R_{f_i}$-模 $M_i$，以及同构
$\varphi_i : \mathcal{F}|_{D(f_i)} \to \mathcal{F}_{M_i}$，其中再次把
相应的 $R_{f_i}$-模记为 $M_i$。在交叠处得到同构
$$
\xymatrix{
\mathcal{F}_{M_i}|_{D(f_if_j)}
\ar[rr]^{\varphi_i^{-1}|_{D(f_if_j)}}
& &
\mathcal{F}|_{D(f_if_j)}
\ar[rr]^{\varphi_j|_{D(f_if_j)}}
& &
\mathcal{F}_{M_j}|_{D(f_if_j)}.
}
$$
把这些同构记作 $\psi_{ij}$。显然，在三重交叠处有余循环条件
$$
\psi_{jk}|_{D(f_if_jf_k)}
\circ
\psi_{ij}|_{D(f_if_jf_k)}
=
\psi_{ik}|_{D(f_if_jf_k)}
$$
。

\medskip\noindent
回顾，每个开子空间 $D(f_i)$、$D(f_if_j)$、$D(f_if_jf_k)$ 都是仿射概形。
因此，各层 $\mathcal{F}_{M_i}$ 都同构于层 $\widetilde M_i$；这是上面的
引理 \ref{schemes-lemma-compare-constructions} 的结论。特别地，
$\mathcal{F}_{M_i}(D(f_if_j)) = (M_i)_{f_j}$，等等。再由上面的引理
\ref{schemes-lemma-compare-constructions} 可知，$\psi_{ij}$ 对应于唯一的
$R_{f_if_j}$-模同构
$$
\psi_{ij} : (M_i)_{f_j} \longrightarrow (M_j)_{f_i}
$$
即 $\psi_{ij}$ 在 $D(f_if_j)$ 上的截面所诱导的映射。此外，这些同构满足
如下余循环条件：
$$
\xymatrix{
(M_i)_{f_jf_k}
\ar[rd]_{\psi_{ij}}
\ar[rr]^{\psi_{ik}}
& &
(M_k)_{f_if_j} \\
&
(M_j)_{f_if_k} \ar[ru]_{\psi_{jk}}
}
$$
对任意三元组 $i, j, k$，该图交换。

\medskip\noindent
现在，《代数》的引理 \ref{algebra-lemma-glue-modules} 表明，存在一个
$R$-模 $M$，使得 $M_i = M_{f_i}$，并与各态射 $\psi_{ij}$ 相容。
考虑 $\mathcal{F}_M = \widetilde M$。此时，要证明 $\widetilde M$ 同构于
我们最初的拟凝聚层 $\mathcal{F}$，只剩一个形式性步骤。确切地说，
层 $\mathcal{F}$ 与 $\widetilde M$ 给出的 $\mathcal{O}_X$-模层粘合数据
彼此同构；这里所取覆盖为 $X = \bigcup D(f_i)$。参见《层》第
\ref{sheaves-section-glueing-sheaves} 节，尤其是引理
\ref{sheaves-lemma-mapping-property-glue}。具体到当前情形，这归结为如下论证：
我们构造一个 $R$-模映射
$$
M \longrightarrow \Gamma(X, \mathcal{F}).
$$
具体而言，给定 $m \in M$，得到 $m_i = m/1 \in M_{f_i} = M_i$；
这是 $M$ 的构造所致。由 $M_i$ 的构造，这对应于一个截面
$s_i \in \mathcal{F}(U_i)$（即 $\varphi^{-1}_i(m_i)$）。我们断言
$s_i|_{D(f_if_j)} = s_j|_{D(f_if_j)}$。这成立是因为，由 $M$ 的构造，
有 $\psi_{ij}(m_i) = m_j$，并且各 $\psi_{ij}$ 正是按上述方式构造的。
由 $\mathcal{F}$ 的层条件，这组截面唯一地给出一个截面 $s$；它是
$\mathcal{F}$ 在 $X$ 上的截面。留给读者验证 $m \mapsto s$ 是
$R$-模映射。由引理
\ref{schemes-lemma-compare-constructions}，得到相应的 $\mathcal{O}_X$-模映射
$$
\widetilde M \longrightarrow \mathcal{F}.
$$
按构造，该映射化为同构 $\varphi_i^{-1}$；在每个 $D(f_i)$ 上皆如此，
因而该映射本身是同构。
\end{proof}

\begin{lemma}
\label{schemes-lemma-equivalence-quasi-coherent}
设 $(X, \mathcal{O}_X) = (\Spec(R), \mathcal{O}_{\Spec(R)})$
为仿射概形。函子 $M \mapsto \widetilde M$ 与
$\mathcal{F} \mapsto \Gamma(X, \mathcal{F})$ 给出如下范畴之间互为拟逆的等价：
$$
\xymatrix{
\QCoh(\mathcal{O}_X)
\ar@<1ex>[r]
&
\text{Mod}_R
\ar@<1ex>[l]
}
$$
即拟凝聚 $\mathcal{O}_X$-模范畴与 $R$-模范畴之间的等价。
\end{lemma}

\begin{proof}
见上面的引理 \ref{schemes-lemma-compare-constructions} 和
\ref{schemes-lemma-quasi-coherent-affine}。
\end{proof}

\noindent
从现在起，我们不再区分仿射概形上的拟凝聚层与形如 $\widetilde M$ 的层。

\begin{lemma}
\label{schemes-lemma-kernel-cokernel-quasi-coherent}
设 $X = \Spec(R)$ 为仿射概形。拟凝聚 $\mathcal{O}_X$-模之间映射的核与余核
都是拟凝聚的。
\end{lemma}

\begin{proof}
这来自函子 $\widetilde{\ }$ 的正合性，因为由引理
\ref{schemes-lemma-compare-constructions} 可知，任意映射
$\psi : \widetilde{M} \to \widetilde{N}$ 都来自一个 $R$-模映射
$\varphi : M \to N$。（所以有
$\Ker(\psi) = \widetilde{\Ker(\varphi)}$ 和
$\Coker(\psi) = \widetilde{\Coker(\varphi)}$。）
\end{proof}

\begin{lemma}
\label{schemes-lemma-colimit-quasi-coherent}
设 $X = \Spec(R)$ 为仿射概形。$X$ 上任意一族拟凝聚层的直和仍是拟凝聚的；
余极限也同样如此。
\end{lemma}

\begin{proof}
设 $\mathcal{F}_i$（$i \in I$）是 $X$ 上的一族拟凝聚层。由上面的引理
\ref{schemes-lemma-equivalence-quasi-coherent}，可写成
$\mathcal{F}_i = \widetilde{M_i}$，其中它来自某个 $R$-模 $M_i$。
令 $M = \bigoplus M_i$，并考虑层
$\widetilde{M}$。对每个标准开集 $D(f)$，有
$$
\widetilde{M}(D(f)) = M_f =
\left(\bigoplus M_i\right)_f =
\bigoplus M_{i, f}.
$$
因此可见，拟凝聚 $\mathcal{O}_X$-模 $\widetilde{M}$ 正是各层
$\mathcal{F}_i$ 的直和。一般余极限的情形可用类似论证处理。
\end{proof}

\begin{lemma}
\label{schemes-lemma-extension-quasi-coherent}
设 $(X, \mathcal{O}_X) = (\Spec(R), \mathcal{O}_{\Spec(R)})$
为仿射概形。设
$$
0 \to
\mathcal{F}_1 \to
\mathcal{F}_2 \to
\mathcal{F}_3 \to
0
$$
是 $\mathcal{O}_X$-模层的短正合列。若三项中有两项是拟凝聚的，
则第三项也是拟凝聚的。
\end{lemma}

\begin{proof}
若 $\mathcal{F}_1$ 和 $\mathcal{F}_2$ 都是拟凝聚的，则断言显然成立，
因为函子 $M \mapsto \widetilde M$ 是正合的；见引理
\ref{schemes-lemma-spec-sheaves}。$\mathcal{F}_2$ 和 $\mathcal{F}_3$ 都拟凝聚的情形
也类似。现在设 $\mathcal{F}_1 = \widetilde M_1$ 和
$\mathcal{F}_3 = \widetilde M_3$ 都是拟凝聚的。令
$M_2 = \Gamma(X, \mathcal{F}_2)$。我们断言，只须证明序列
$$
0 \to M_1 \to M_2 \to M_3 \to 0
$$
正合。事实上，若它正合，则利用引理 \ref{schemes-lemma-compare-constructions}
的映射性质可得到交换图
$$
\xymatrix{
0 \ar[r] &
\widetilde M_1 \ar[r] \ar[d] &
\widetilde M_2 \ar[r] \ar[d] &
\widetilde M_3 \ar[r] \ar[d] &
0 \\
0 \ar[r] &
\mathcal{F}_1 \ar[r] &
\mathcal{F}_2 \ar[r] &
\mathcal{F}_3 \ar[r] &
0
}
$$
于是由蛇引理即得结论。

\medskip\noindent
这里的“标准”论证应当先证明 $H^1(X, \mathcal{F}) = 0$，其中
$\mathcal{F}$ 是任意拟凝聚层。这其实并不十分困难，但最好留到后面再做。
这里我们改用一个小技巧。

\medskip\noindent
取 $m \in M_3 = \Gamma(X, \mathcal{F}_3)$。考虑如下集合：
$$
I = \{ f \in R \mid \text{元素 }fm\text{ 来自 }M_2\}.
$$
显然这是一个理想。只须证明 $1 \in I$。因此，只须证明对任意素理想
$\mathfrak p$，都存在 $f \in I$，且 $f \not\in \mathfrak p$。
设 $x \in X$ 为对应于 $\mathfrak p$ 的点。由于满射性可以在茎上检验，
存在一个开邻域 $U$，它是 $x$ 的邻域，并使得 $m|_U$ 来自某个局部截面
$s \in \mathcal{F}_2(U)$。事实上，可以假定 $U = D(f)$ 是标准开集，
即 $f \in R$ 且 $f \not \in \mathfrak p$。我们将证明，对某个
$N \gg 0$ 有 $f^N \in I$，这便完成证明。

\medskip\noindent
任取一点 $z \in V(f)$，设它对应于素理想 $\mathfrak q \subset R$。
还可以找到 $g \in R$，满足 $g \not \in \mathfrak q$，使得
$m|_{D(g)}$ 可提升为某个 $s' \in \mathcal{F}_2(D(g))$。考虑差
$s|_{D(fg)} - s'|_{D(fg)}$。这是一个元素 $m'$，它属于
$\mathcal{F}_1(D(fg)) = (M_1)_{fg}$。对某个整数
$n = n(z)$，元素 $f^n m'$ 来自某个 $m'_1 \in (M_1)_g$。可见
$f^n s$ 延拓为一个截面 $\sigma$，它是 $\mathcal{F}_2$ 在
$D(f) \cup D(g)$ 上的截面；这是因为它与 $f^n s' + m'_1$ 的限制
在 $D(f) \cap D(g) = D(fg)$ 上相同。此外，$\sigma$ 映到 $f^n m$ 在
$D(f) \cup D(g)$ 上的限制。

\medskip\noindent
由于 $V(f)$ 是拟紧的，存在有限多个元素 $g_1, \ldots, g_m \in R$，
使得 $V(f) \subset \bigcup D(g_j)$；还存在一个整数 $n > 0$ 以及截面
$\sigma_j \in \mathcal{F}_2(D(f) \cup D(g_j))$，使得
$\sigma_j|_{D(f)} = f^n s$，并且 $\sigma_j$ 映到截面
$f^nm|_{D(f) \cup D(g_j)}$；后者是 $\mathcal{F}_3$ 的截面。考虑这些差：
$$
\sigma_j|_{D(f) \cup D(g_jg_k)}
-
\sigma_k|_{D(f) \cup D(g_jg_k)}.
$$
它们对应于 $\mathcal{F}_1$ 在 $D(f) \cup D(g_jg_k)$ 上的截面，
并且在 $D(f)$ 上为零。特别地，它们在
$\mathcal{F}_1(D(g_jg_k)) = (M_1)_{g_jg_k}$ 中的像，在
$(M_1)_{g_jg_kf}$ 中为零。因此，$f$ 的某个高次幂把每一个这样的像都
零化。换言之，各元素 $f^N \sigma_j$ 对某个 $N \gg 0$ 满足层性质的
粘合条件，并给出一个截面 $\sigma $；它属于 $\mathcal{F}_2$，定义在
$\bigcup (D(f) \cup D(g_j)) = X$ 上，正如所需。
\end{proof}







\section{仿射概形的闭子空间}
\label{schemes-section-closed-in-affine}

\begin{example}
\label{schemes-example-closed-immersion-affines}
设 $R$ 为环，$I \subset R$ 为理想。考虑仿射概形的态射
$i : Z = \Spec(R/I) \to \Spec(R) = X$。由《代数》的引理
\ref{algebra-lemma-spec-closed}，它把 $Z$ 同胚地映到 $X$ 的一个闭子集上。
此外，若 $I \subset \mathfrak p \subset R$ 是对应于点 $x = i(z)$ 的
素理想，其中 $x \in X$、$z \in Z$，则在茎上得到映射
$$
\mathcal{O}_{X, x} = R_{\mathfrak p}
\longrightarrow
R_{\mathfrak p}/IR_{\mathfrak p} = \mathcal{O}_{Z, z}
$$
由此可见，$i$ 是局部环层空间的闭浸入；见定义
\ref{schemes-definition-closed-immersion-locally-ringed-spaces}。显然，它同构于由
拟凝聚理想层 $\widetilde I$ 所关联的闭子空间，如例
\ref{schemes-example-closed-subspace} 所述。
\end{example}

\begin{lemma}
\label{schemes-lemma-closed-immersion-affine-case}
\begin{slogan}
对仿射概形，闭浸入对应于理想。
\end{slogan}
设 $(X, \mathcal{O}_X) = (\Spec(R), \mathcal{O}_{\Spec(R)})$
为仿射概形，$i : Z \to X$ 为局部环层空间的任意闭浸入。则存在唯一理想
$I \subset R$，使得态射 $i : Z \to X$ 可与闭浸入
$\Spec(R/I) \to \Spec(R)$ 认同；后一个闭浸入是在上面的例
\ref{schemes-example-closed-immersion-affines} 中构造的。
\end{lemma}

\begin{proof}
这几乎是显然的！事实上，由引理 \ref{schemes-lemma-closed-immersion}，可以把
$Z \to X$ 与由理想层 $\mathcal{I} \subset \mathcal{O}_X$ 所关联的闭子空间
认同，如定义 \ref{schemes-definition-closed-subspace} 和例
\ref{schemes-example-closed-subspace} 所述。依照我们的约定，该理想层作为
$\mathcal{O}_X$-模层局部地由截面生成。因此，商层
$\mathcal{O}_X / \mathcal{I}$ 在 $X$ 上局部地是某个映射
$\bigoplus_{j \in J} \mathcal{O}_U \to \mathcal{O}_U$ 的余核。
所以按定义，$\mathcal{O}_X / \mathcal{I}$ 是拟凝聚的。由第
\ref{schemes-section-quasi-coherent-affine} 节的结果，它形如 $\widetilde S$，
是由某个 $R$-模 $S$ 得到的。此外，由于
$\mathcal{O}_X = \widetilde R \to \widetilde S$ 是满射，由引理
\ref{schemes-lemma-extension-quasi-coherent} 可见，$\mathcal{I}$ 也是拟凝聚的，
记作 $\mathcal{I} = \widetilde I$。当然，$I \subset R$ 且 $S = R/I$，
一切都已清楚。
\end{proof}













\section{概形}
\label{schemes-section-schemes}

\begin{definition}
\label{schemes-definition-scheme}
\begin{history}
在 \cite{EGA1} 中，我们现在所谓的概形称为“pre-sch\'ema”，而“sch\'ema”
一词留给 Stacks Project 中所谓的分离概形。在第二版
\cite{EGA1-second} 中，术语改成了如今的标准用法。不过，偶尔仍会见到
“prescheme”这一术语，例如 \cite{Murre-lectures}。
\end{history}
{\it 概形}是满足如下性质的局部环层空间：每一点都有一个本身为仿射概形的
开邻域。{\it 概形态射}是局部环层空间的态射。概形范畴记作 $\Sch$。
\end{definition}

\noindent
设 $X$ 为概形。我们将作如下（极轻微的）语言滥用：若开子空间
$U \subset X$ 中的 $U$ 是仿射概形，就称它为{\it 仿射开集}或
{\it 开仿射集}。
我们常写 $U = \Spec(R)$，以表示 $U$ 同构于 $\Spec(R)$，并且还表示
我们将暂时认同 $U$ 与 $\Spec(R)$。

\begin{lemma}
\label{schemes-lemma-open-subspace-scheme}
设 $X$ 为概形，$j : U \to X$ 为局部环层空间的开浸入。则 $U$ 是概形。
特别地，$X$ 的任意开子空间都是概形。
\end{lemma}

\begin{proof}
设 $U \subset X$，并取 $u \in U$。选取仿射开邻域
$u \in V \subset X$。由于 $V$ 的标准开集构成 $V$ 上拓扑的一组基，
存在 $f\in \mathcal{O}_V(V)$，使得 $u \in D(f) \subset U$。而由引理
\ref{schemes-lemma-standard-open-affine}，$D(f)$ 是仿射概形。这就证明了 $U$ 的
每一点都有一个仿射开邻域。
\end{proof}

\noindent
显然，该引理（或其证明）表明，任意概形 $X$ 的拓扑都有一组由仿射开集
构成的基（见《拓扑》第 \ref{topology-section-bases} 节）。

\begin{example}
\label{schemes-example-not-affine}
设 $k$ 为域。开子空间
$U = \Spec(k[x, y]) \setminus \{ (x, y)\}$ 位于仿射概形
$X =\Spec(k[x, y])$ 中，并给出了一个非仿射概形的例子。
它由以下两个仿射开集覆盖：
\begin{center}
$D(x) = \Spec(k[x, y, 1/x])$，\qquad
$D(y) = \Spec(k[x, y, 1/y])$。
\end{center}
二者之交为 $D(xy) = \Spec(k[x, y, 1/xy])$。由 $\mathcal{O}_U$ 的层性质，存在正合列
$$
0 \to
\Gamma(U, \mathcal{O}_U) \to
k[x, y, 1/x] \times k[x, y, 1/y] \to
k[x, y, 1/xy]
$$
由此可知，映射 $k[x, y] \to \Gamma(U, \mathcal{O}_U)$（它来自态射
$U \to X$）是同构。因此 $U$ 不可能是仿射的；否则由引理
\ref{schemes-lemma-category-affine-schemes} 将有 $U \cong X$。
\end{example}











\section{概形的浸入}
\label{schemes-section-immersions}

\noindent
在引理 \ref{schemes-lemma-open-subspace-scheme} 中，我们看到概形的任意开子空间
仍是概形。下面将证明，概形的闭子空间也有同样的性质。

\medskip\noindent
注意，拟凝聚 $\mathcal{O}_X$-模层的概念对任意环层空间 $X$ 都有定义，
特别是在 $X$ 为概形时。由第 \ref{schemes-section-quasi-coherent-affine} 节的结果
可知，这样的层在任意仿射开集 $U \subset X$ 上都形如 $\widetilde M$，
其中它来自某个 $\mathcal{O}_X(U)$-模 $M$。

\begin{lemma}
\label{schemes-lemma-closed-subspace-scheme}
设 $X$ 为概形，$i : Z \to X$ 为局部环层空间的闭浸入。
\begin{enumerate}
\item 局部环层空间 $Z$ 是概形；
\item 核 $\mathcal{I}$ 是拟凝聚理想层；这里所取映射为
$\mathcal{O}_X \to i_*\mathcal{O}_Z$；
\item 对任意仿射开集 $U = \Spec(R)$，若它属于 $X$，则态射
$i^{-1}(U) \to U$ 可与 $\Spec(R/I) \to \Spec(R)$ 认同，其中
$I \subset R$ 是某个理想；
\item 有 $\mathcal{I}|_U = \widetilde I$。
\end{enumerate}
特别地，任意局部由截面生成的理想层都是拟凝聚理想层（反之亦然），
而且 $X$ 的任意闭子空间都是概形。
\end{lemma}

\begin{proof}
设 $i : Z \to X$ 为闭浸入，并取一点 $z \in Z$。任取一个仿射开邻域
$i(z) \in U \subset X$，记 $U = \Spec(R)$。由引理
\ref{schemes-lemma-closed-immersion-affine-case} 可知，$i^{-1}(U) \to U$ 可与
仿射概形的态射 $\Spec(R/I) \to \Spec(R)$ 认同。首先，这说明
$z \in i^{-1}(U) \subset Z$ 是 $z$ 的仿射邻域，因而 $Z$ 是概形。
其次，这说明 $\mathcal{I}|_U$ 就是 $\widetilde I$。换言之，对每一点
$x \in i(Z)$，都存在一个开邻域，使得 $\mathcal{I}$ 在该邻域中拟凝聚。
注意，$\mathcal{I}|_{X \setminus i(Z)}
\cong \mathcal{O}_{X \setminus i(Z)}$。因此，该理想层在
$X \setminus i(Z)$ 上的限制也是拟凝聚的。由此可知，$\mathcal{I}$ 拟凝聚。
\end{proof}

\begin{definition}
\label{schemes-definition-immersion}
设 $X$ 为概形。
\begin{enumerate}
\item 若一个概形态射作为局部环层空间态射是开浸入，就称它为
{\it 开浸入}（见定义 \ref{schemes-definition-immersion-locally-ringed-spaces}）。
\item $X$ 的{\it 开子概形}是 $X$ 的开子空间，其含义见定义
\ref{schemes-definition-open-subspace}；$X$ 的开子概形由引理
\ref{schemes-lemma-open-subspace-scheme} 可知是概形。
\item 若一个概形态射作为局部环层空间态射是闭浸入，就称它为
{\it 闭浸入}（见定义
\ref{schemes-definition-closed-immersion-locally-ringed-spaces}）。
\item $X$ 的{\it 闭子概形}是 $X$ 的闭子空间，其含义见定义
\ref{schemes-definition-closed-subspace}；由引理 \ref{schemes-lemma-closed-subspace-scheme}，
闭子概形是概形。
\item 若概形态射 $f : X \to Y$ 可分解为 $j \circ i$，其中 $i$ 是闭浸入、
$j$ 是开浸入，就称它为{\it 浸入}或{\it 局部闭浸入}。
\end{enumerate}
\end{definition}

\noindent
由第 \ref{schemes-section-open-immersion} 节和第
\ref{schemes-section-closed-immersion} 节中的引理可知，概形的任意开浸入
（相应地，闭浸入）都同构于目标中某个开子概形（相应地，闭子概形）的
包含映射。

\medskip\noindent
我们对闭浸入的定义介于 Hartshorne 与 EGA 的定义之间。Hartshorne 把
闭浸入定义为概形态射 $f : X \to Y$：它把 $X$ 同胚地映到 $Y$ 的一个
闭子集上，并且 $f^\# : \mathcal{O}_Y \to f_*\mathcal{O}_X$ 是满射；见
\cite[Page 85]{H}。我们将在引理
\ref{schemes-lemma-characterize-closed-immersions} 中证明，这与我们的概念等价。
在 \cite{EGA} 中，Grothendieck 和 Dieudonn\'e 首先利用拟凝聚理想层，
通过例 \ref{schemes-example-closed-subspace} 的构造定义闭子概形；然后把闭浸入
定义为一个态射 $f : X \to Y$，它诱导出与某个闭子概形的同构。由引理
\ref{schemes-lemma-closed-subspace-scheme} 可知，这与我们的概念一致。

\medskip\noindent
从教学角度说，上面的定义简直是一场灾难。给学生讲授这些材料时，我们发现
常常更方便的做法是把闭浸入定义为概形的仿射态射 $f : X \to Y$，并要求
$f^\# : \mathcal{O}_Y \to f_*\mathcal{O}_X$ 是满射。事实上，仿射态射这一
概念（《态射》第 \ref{morphisms-section-affine} 节）十分自然且容易理解。

\medskip\noindent
关于闭浸入的更多信息，建议读者参阅《态射》第
\ref{morphisms-section-closed-immersions} 节和第
\ref{morphisms-section-closed-immersions-quasi-coherent} 节。

\medskip\noindent
本节末尾将讨论局部闭子概形与浸入。

\begin{remark}
\label{schemes-remark-not-reverse-open-closed}
若 $f : X \to Y$ 是概形的浸入，则一般不能把 $f$ 分解成先开浸入、
后闭浸入的复合。见《态射》的例 \ref{morphisms-example-thibaut}。
\end{remark}

\begin{lemma}
\label{schemes-lemma-immersion-when-closed}
设 $f : Y \to X$ 为概形的浸入。则 $f$ 是闭浸入，当且仅当
$f(Y) \subset X$ 是闭子集。
\end{lemma}

\begin{proof}
若 $f$ 是闭浸入，则按定义 $f(Y)$ 为闭集。反过来，设 $f(Y)$ 为闭集。
按定义，存在开子概形 $U \subset X$，使得 $f$ 是闭浸入
$i : Y \to U$ 与开浸入 $j : U \to X$ 的复合。设
$\mathcal{I} \subset \mathcal{O}_U$ 为与闭浸入 $i$ 关联的拟凝聚理想层。
注意
$\mathcal{I}|_{U \setminus i(Y)}
= \mathcal{O}_{U \setminus i(Y)}
= \mathcal{O}_{X \setminus i(Y)}|_{U \setminus i(Y)}$.
所以可以把它们粘合（见《层》第 \ref{sheaves-section-glueing-sheaves} 节）：
将 $\mathcal{I}$ 与 $\mathcal{O}_{X \setminus i(Y)}$ 粘合成理想层
$\mathcal{J} \subset \mathcal{O}_X$。由于 $X$ 的每一点都有一个邻域，
使得 $\mathcal{J}$ 在该邻域上拟凝聚，可见 $\mathcal{J}$ 是拟凝聚的
（特别地，它局部由截面生成）。按构造，$\mathcal{O}_X/\mathcal{J}$
的支集包含在 $U$ 中，并且限制到那里就等于
$\mathcal{O}_U/\mathcal{I}$。因此，与 $\mathcal{I}$ 和 $\mathcal{J}$
关联的闭子空间典范同构；见例 \ref{schemes-example-closed-subspace}。特别地，
$U$ 中与 $\mathcal{I}$ 关联的闭子空间同构于 $X$ 的一个闭子空间。
由于 $Y \to U$ 与由 $\mathcal{I}$ 关联的闭子空间认同（见引理
\ref{schemes-lemma-closed-immersion}），所以 $Y \to U \to X$ 是闭浸入。
\end{proof}

\noindent
设 $f : Y \to X$ 为浸入。令
$Z = \overline{f(Y)} \setminus f(Y)$；这是 $X$ 的闭子集。再令
$U = X \setminus Z$。该引理说明，$U$ 是 $X$ 的最大开子空间，使得
$f : Y \to X$ 可通过一个进入 $U$ 的闭浸入分解。我们把 $X$ 的一个
{\it 局部闭子概形}定义为二元组 $(Z, U)$，其中 $Z$ 是开子概形 $U$ 的
闭子概形，而后者又是 $X$ 的开子概形，并且还满足
$\overline{Z} \cup U = X$。通常只说“设 $Z$ 为 $X$ 的局部闭子概形”，
因为可以恢复 $U$，所依据的是态射 $Z \to X$。上面的论述表明，任意浸入
$f : Y \to X$ 唯一地分解为 $Y \to Z \to X$，其中 $Z$ 是 $X$ 的
局部闭子空间，而 $Y \to Z$ 是同构。

\medskip\noindent
这一做法的意义在于，$X$ 的所有局部闭子概形构成一个集合。可以在该集合上
定义一个偏序；出于显然的理由，我们称之为包含关系。具体地，若
$Z \to X$ 和 $Z' \to X$ 是 $X$ 的两个局部闭子概形，那么我们称
{\it $Z$ 包含于 $Z'$}，当且仅当态射 $Z \to X$ 通过 $Z'$ 分解。若确实如此，
则 $Z$ 当然与 $Z'$ 的唯一一个局部闭子概形认同，依此类推。

\medskip\noindent
关于浸入的更多信息，参见《态射》第 \ref{morphisms-section-immersions} 节。










\section{概形的 Zariski 拓扑}
\label{schemes-section-topology}

\noindent
关于适用于概形 Zariski 拓扑的一些基础拓扑材料，见《拓扑》第
\ref{topology-section-introduction} 节。

\begin{lemma}
\label{schemes-lemma-scheme-sober}
设 $X$ 为概形。$X$ 的任意不可约闭子集都有唯一泛点。换言之，
$X$ 是清醒拓扑空间；见《拓扑》的定义
\ref{topology-definition-generic-point}。
\end{lemma}

\begin{proof}
设 $Z \subset X$ 为不可约闭子集。对每个仿射开集 $U \subset X$，
$U = \Spec(R)$，都有 $Z \cap U = V(I)$，其中 $I \subset R$ 是唯一根理想。
注意，$Z \cap U$ 或为空，或不可约。在后一种情形（它至少对一个 $U$
发生），可见 $I = \mathfrak p$ 是素理想；它给出一个泛点 $\xi$，
属于 $Z \cap U$。于是 $Z = \overline{\{\xi\}}$；换言之，
$\xi$ 是 $Z$ 的泛点。若 $\xi'$ 是另一个泛点，则
$\xi' \in Z \cap U$，从而立即得到 $\xi' = \xi$。
\end{proof}

\begin{lemma}
\label{schemes-lemma-basis-affine-opens}
设 $X$ 为概形。$X$ 的所有仿射开集构成 $X$ 上拓扑的一组基。
\end{lemma}

\begin{proof}
这来自第 \ref{schemes-section-schemes} 节中关于开子概形的讨论。
\end{proof}

\begin{remark}
\label{schemes-remark-intersection-affine-opens}
一般而言，$X$ 中两个仿射开集的交不再是仿射开集。见例
\ref{schemes-example-affine-space-zero-doubled}。
\end{remark}

\begin{lemma}
\label{schemes-lemma-locally-quasi-compact}
任意概形的底拓扑空间都是局部拟紧的；见《拓扑》的定义
\ref{topology-definition-locally-quasi-compact}。
\end{lemma}

\begin{proof}
这来自上面的引理 \ref{schemes-lemma-basis-affine-opens}，以及环的谱是拟紧的这一
事实；见《代数》的引理 \ref{algebra-lemma-quasi-compact}。
\end{proof}

\begin{lemma}
\label{schemes-lemma-standard-open-two-affines}
设 $X$ 为概形。设 $U, V$ 为 $X$ 的仿射开集，且 $x \in U \cap V$。
存在一个仿射开邻域 $W$，它是 $x$ 的邻域，并使得 $W$ 同时是 $U$ 和
$V$ 的标准开集。
\end{lemma}

\begin{proof}
写成 $U = \Spec(A)$ 和 $V = \Spec(B)$。设 $x$ 对应于素理想
$\mathfrak p \subset A$ 以及素理想 $\mathfrak q \subset B$。可以选取
$f \in A$，$f \not \in \mathfrak p$，使得 $D(f) \subset U \cap V$。
注意，$D(f)$ 的任意标准开集都是 $\Spec(A) = U$ 的标准开集。因此可以
假定 $U \subset V$。换言之，现在可以把 $U$ 看成 $V$ 的一个仿射开集。
接着选取 $g \in B$，$g \not \in \mathfrak q$，使得 $D(g) \subset U$。
此时可见 $D(g) = D(g_A)$，其中 $g_A \in A$ 表示 $g$ 在映射
$B \to A$ 下的像。引理得证。
\end{proof}

\begin{lemma}
\label{schemes-lemma-good-subcover}
设 $X$ 为概形。设 $X = \bigcup_i U_i$ 为仿射开覆盖，并设
$V \subset X$ 为仿射开集。存在标准开覆盖
$V = \bigcup_{j = 1, \ldots, m} V_j$（见定义
\ref{schemes-definition-standard-covering}），使得每个 $V_j$ 都是某个 $U_i$
中的标准开集。
\end{lemma}

\begin{proof}
取 $v \in V$。则有 $v \in U_i$，其中 $i$ 是某个指标。由上面的引理
\ref{schemes-lemma-standard-open-two-affines}，存在开集
$v \in W_v \subset V \cap U_i$，使得 $W_v$ 同时是 $V$ 和 $U_i$
中的标准开集。由于 $V$ 拟紧，引理得证。
\end{proof}

\begin{lemma}
\label{schemes-lemma-sheaf-on-affines}
设 $X$ 为概形。设 $\mathcal{B}$ 为 $X$ 的所有仿射开集构成的集合。
设 $\mathcal{F}$ 为 $\mathcal{B}$ 上的集合预层，见《层》的定义
\ref{sheaves-definition-presheaf-basis}。以下条件等价：
\begin{enumerate}
\item $\mathcal{F}$ 是 $X$ 上某个层限制到 $\mathcal{B}$ 的结果；
\item $\mathcal{F}$ 是 $\mathcal{B}$ 上的层；
\item $\mathcal{F}(\emptyset)$ 是单点集，并且每当
$U = V \cup W$，其中 $U, V, W \in \mathcal{B}$ 且
$V, W \subset U$ 是标准开集（《代数》的定义
\ref{algebra-definition-Zariski-topology}），映射
$$
\mathcal{F}(U) \longrightarrow \mathcal{F}(V) \times \mathcal{F}(W)
$$
是单射，其像由所有二元组 $(s, t)$ 构成，并且这些二元组满足
$s|_{V \cap W} = t|_{V \cap W}$。
\end{enumerate}
\end{lemma}

\begin{proof}
(1) 与 (2) 的等价性就是《层》的引理
\ref{sheaves-lemma-restrict-basis-equivalence}。显然，(2) 蕴含 (3)。
所以只须证明 (3) 蕴含 (2)。由《层》的引理
\ref{sheaves-lemma-cofinal-systems-coverings-standard-case} 和引理
\ref{schemes-lemma-standard-open}，只须对标准开覆盖（定义
\ref{schemes-definition-standard-covering}）证明层条件；这里覆盖的是
$\mathcal{B}$ 中的元素。设
$U = U_1 \cup \ldots \cup U_n$ 是标准开覆盖，其中 $U \subset X$ 为
仿射开集。我们对 $n$ 归纳证明该覆盖满足层条件。若 $n = 0$，则 $U$
为空，层条件由假设成立。若 $n = 1$，则无事可证。若 $n = 2$，
这正是条件 (3)。若 $n > 2$，则写 $U_i = D(f_i)$，其中
$f_i \in A = \mathcal{O}_X(U)$。设 $s_i \in \mathcal{F}(U_i)$ 是截面，
且有 $s_i|_{U_i \cap U_j} = s_j|_{U_i \cap U_j}$，这对所有
$1 \leq i < j \leq n$ 都成立。由于 $U = U_1 \cup \ldots \cup U_n$，有
$1 = \sum_{i = 1, \ldots, n} a_i f_i$；该等式在 $A$ 中成立，其中某个
$a_i \in A$。见
《代数》的引理 \ref{algebra-lemma-Zariski-topology}。令
$g = \sum_{i = 1, \ldots, n - 1} a_if_i$。于是
$U = D(g) \cup D(f_n)$。注意，
$D(g) = D(gf_1) \cup \ldots \cup D(gf_{n - 1})$ 是标准开覆盖。
由归纳假设，存在唯一截面 $s' \in \mathcal{F}(D(g))$，它与
$s_i|_{D(gfi)}$ 对 $i = 1, \ldots, n - 1$ 的每个值都相符。我们断言
$s'$ 与 $s_n$ 在 $D(gf_n)$ 上的限制相同。这来自归纳假设以及覆盖
$D(gf_n) = D(gf_nf_1) \cup \ldots \cup D(gf_nf_{n - 1})$。因此存在唯一截面
$s \in \mathcal{F}(U)$，它在 $D(g)$ 上的限制是 $s'$，在 $D(f_n)$ 上的
限制是 $s_n$。我们略去如下验证：$s$ 的限制是 $s_i$，相应开集为
$D(f_i)$，其中 $i = 1, \ldots, n - 1$；也略去 $s$ 唯一性的验证。
\end{proof}

\begin{lemma}
\label{schemes-lemma-scheme-finite-discrete-affine}
设 $X$ 为概形，其底拓扑空间是有限离散集。则 $X$ 是仿射的。
\end{lemma}

\begin{proof}
写成 $X = \{x_1, \ldots, x_n\}$。则 $U_i = \{x_i\}$ 是 $x_i$ 的开邻域。
由引理 \ref{schemes-lemma-basis-affine-opens}，它是仿射的。因此 $X$ 是有限多个
仿射概形的不交并，进而由引理 \ref{schemes-lemma-disjoint-union-affines} 可知它是
仿射的。
\end{proof}

\begin{example}
\label{schemes-example-scheme-without-closed-points}
存在没有闭点的概形。事实上，设 $R$ 为局部整环，其谱形如
$(0) = \mathfrak p_0 \subset \mathfrak p_1 \subset \mathfrak p_2
\subset \ldots \subset \mathfrak m$。那么开子概形
$\Spec(R) \setminus \{\mathfrak m\}$ 没有闭点。为说明这样的环 $R$
确实存在，使用如下事实：给定任意全序群 $(\Gamma, \geq)$，存在赋值环
$A$，其值群为 $(\Gamma, \geq)$；见 \cite{Krull}。记号见《代数》第
\ref{algebra-section-valuation-rings} 节。我们取
$\Gamma = \mathbf{Z}x_1 \oplus \mathbf{Z}x_2 \oplus \mathbf{Z}x_3 \oplus
\ldots$，并定义 $\sum_i a_i x_i \geq 0$ 当且仅当第一个非零的 $a_i$
满足 $> 0$，或所有 $a_i = 0$。于是
$x_1 \geq x_2 \geq x_3 \geq  \ldots \geq 0$。各子集
$x_i + \Gamma_{\geq 0}$ 是 $(\Gamma, \geq)$ 的素理想；参见《代数》的
引理 \ref{algebra-lemma-ideals-valuation-ring} 上方的记号。它们连同
$\emptyset$ 与 $\Gamma_{\geq 0}$ 是全部素理想。因此，$A$ 是一个环的
例子，其谱具有上述结构；这来自《代数》的引理
\ref{algebra-lemma-ideals-valuation-ring}。
\end{example}



















\section{约化概形}
\label{schemes-section-reduced}

\begin{definition}
\label{schemes-definition-reduced}
设 $X$ 为概形。称 $X$ 是{\it 约化的}，如果每个局部环
$\mathcal{O}_{X, x}$ 都是约化的。
\end{definition}

\begin{lemma}
\label{schemes-lemma-reduced}
概形 $X$ 是约化的，当且仅当 $\mathcal{O}_X(U)$ 是约化环；这里
$U \subset X$ 是任意开集。
\end{lemma}

\begin{proof}
先设 $X$ 约化。取截面 $f \in \mathcal{O}_X(U)$，并设 $f^n = 0$。
则 $f$ 在 $\mathcal{O}_{U, u}$ 中的像为零；这对所有 $u \in U$ 都成立。
所以 $f$ 为零；见《层》的引理
\ref{sheaves-lemma-sheaf-subset-stalks}。反过来，设 $\mathcal{O}_X(U)$
约化；这对所有开集 $U$ 都成立。任取非零元素
$f \in \mathcal{O}_{X, x}$。任意代表 $(U, f \in \mathcal{O}(U))$，
只要它代表 $f$，就非零，因而不是幂零元。
所以 $f$ 在 $\mathcal{O}_{X, x}$ 中不是幂零元。
\end{proof}

\begin{lemma}
\label{schemes-lemma-affine-reduced}
仿射概形 $\Spec(R)$ 约化，当且仅当 $R$ 约化。
\end{lemma}

\begin{proof}
正向蕴含立即来自上面的引理 \ref{schemes-lemma-reduced}。反向蕴含则来自如下事实：
约化环的任意局部化仍约化，特别地，约化环的各局部环都约化。
\end{proof}

\begin{lemma}
\label{schemes-lemma-reduced-closed-subscheme}
设 $X$ 为概形，$T \subset X$ 为闭子集。存在唯一闭子概形
$Z \subset X$，满足：(a) $Z$ 的底拓扑空间等于 $T$；(b) $Z$ 约化。
\end{lemma}

\begin{proof}
设 $\mathcal{I} \subset \mathcal{O}_X$ 为由下式定义的子预层：
$$
\mathcal{I}(U) = \{f \in \mathcal{O}_X(U) \mid
f(t) = 0\text{ 对所有 }t \in T\cap U\}
$$
这里用 $f(t)$ 表示 $f$ 在剩余域 $\kappa(t)$ 中的像；该域是 $X$ 于
$t$ 处的剩余域。
由于该条件具有局部性，显然 $\mathcal{I}$ 是理想层。此外，设
$U = \Spec(R)$ 为仿射开集。可以写成 $T \cap U = V(I)$，其中
$I \subset R$ 是唯一根理想。给定素理想 $\mathfrak p \in V(I)$，
它对应于 $t \in T \cap U$；再给定元素 $f \in R$，则有
$f(t) = 0 \Leftrightarrow f \in \mathfrak p$。所以，
$\mathcal{I}(U) = \bigcap_{\mathfrak p \in V(I)} \mathfrak p
= I$；这是《代数》的引理 \ref{algebra-lemma-Zariski-topology} 的结论。
此外，对任意标准开集 $D(g) \subset \Spec(R) = U$，同样的论证给出
$\mathcal{I}(D(g)) = I_g$。所以 $\widetilde I$ 与 $\mathcal{I}|_U$
在一组开集基上作为理想相同，因而二者相等。因此，$\mathcal{I}$ 是拟凝聚
理想层。

\medskip\noindent
现在可以把 $Z$ 定义为由理想层 $\mathcal{I}$ 所关联的闭子空间。
对每个仿射开集 $U = \Spec(R)$，若它属于 $X$，则有
$Z \cap U = \Spec(R/I)$，其中 $I$ 是根理想，因而 $Z$ 约化
（由上面的引理 \ref{schemes-lemma-affine-reduced}）。按构造，$Z$ 的底闭子集就是
$T$。因此已经找到满足 (a)、(b) 的闭子概形。

\medskip\noindent
设 $Z' \subset X$ 是另一个满足 (a)、(b) 的闭子概形。对每个仿射开集
$U = \Spec(R)$，若它属于 $X$，则有
$Z' \cap U = \Spec(R/I')$，其中 $I' \subset R$ 是某个理想。由引理
\ref{schemes-lemma-affine-reduced}，环 $R/I'$ 约化，因而 $I'$ 是根理想。由于
$V(I') = T \cap U = V(I)$，得到 $I = I'$；这是《代数》的引理
\ref{algebra-lemma-Zariski-topology} 的结论。所以 $Z'$ 与 $Z$
由同一个理想层定义，因而相等。
\end{proof}

\begin{definition}
\label{schemes-definition-reduced-induced-scheme}
设 $X$ 为概形，$Z \subset X$ 为闭子集。$Z$ 上的一个{\it 概形结构}
由闭子概形 $Z'$ 给出；它属于 $X$，且其底集合等于 $Z$。我们常说“设
$(Z, \mathcal{O}_Z)$ 是 $Z$ 上的一个概形结构”来表达这一点。$Z$ 上的
{\it 约化诱导概形结构}是引理 \ref{schemes-lemma-reduced-closed-subscheme}
中构造的结构。{\it 约化 $X_{red}$，其对象为 $X$}，是 $X$ 本身上的
约化诱导概形结构。
\end{definition}

\noindent
我们说“设 $Z \subset X$ 为 $X$ 的一个不可约分支”时，通常使用约化诱导
概形结构，把 $Z$ 看成 $X$ 的约化闭子概形。

\begin{remark}
\label{schemes-remark-reduced-induced-locally-closed}
设 $X$ 为概形，$T \subset X$ 为局部闭子集。在这种情形下，我们有时也使用
“$T$ 上的约化诱导概形结构”这一说法。它指如下结构：按照定义
\ref{schemes-definition-reduced-induced-scheme}，把 $T$ 看成开子概形
$X \setminus \partial T$ 中的闭子集；后者是 $X$ 的开子概形。这里
$\partial T = \overline{T} \setminus T$ 是 $T$ 在 $X$ 的拓扑空间中的“边界”。
\end{remark}

\begin{lemma}
\label{schemes-lemma-map-into-reduction}
设 $X$ 为概形，$Z \subset X$ 为闭子概形，$Y$ 为约化概形。态射
$f : Y \to X$ 通过 $Z$ 分解，当且仅当 $f(Y) \subset Z$（在集合论意义下）。
特别地，任意态射 $Y \to X$ 都分解为 $Y \to X_{red} \to X$。
\end{lemma}

\begin{proof}
设 $f(Y) \subset Z$（在集合论意义下）。设
$\mathcal{I} \subset \mathcal{O}_X$ 为 $Z$ 的理想层。对任意仿射开集
$U \subset X$、$\Spec(B) = V \subset Y$，若 $f(V) \subset U$，并任取
$g \in \mathcal{I}(U)$，则其拉回
$b = f^\sharp(g) \in \Gamma(V, \mathcal{O}_Y) = B$ 在任意 $y \in V$
的剩余域中都映为零。换言之，
$b \in \bigcap_{\mathfrak p \subset B} \mathfrak p$。这推出 $b = 0$，
因为 $B$ 约化（引理 \ref{schemes-lemma-reduced} 以及《代数》的引理
\ref{algebra-lemma-Zariski-topology}）。因此，$f$ 通过 $Z$ 分解；
这是引理 \ref{schemes-lemma-characterize-closed-subspace} 的结论。
\end{proof}










\section{概形的点}
\label{schemes-section-points}

\noindent
给定概形 $X$，可以定义函子
$$
h_X : \Sch^{opp}
\longrightarrow
\textit{Sets}, \quad
T \longmapsto \Mor(T, X).
$$
见《范畴》的例 \ref{categories-example-hom-functor}。它称为 $X$ 的
{\it 点函子}。概形理论中颇有趣的一项工作，是描述 $X$ 的内部几何；
这里所用的是函子 $h_X$。本节将给出描述 $X$ 的点的一种简单方法。

\medskip\noindent
设 $X$ 为概形，$R$ 为局部环，其极大理想为 $\mathfrak m \subset R$。
设 $f : \Spec(R) \to X$ 为概形态射。令 $x \in X$ 为闭点
$\mathfrak m \in \Spec(R)$ 的像。于是得到局部环之间的局部同态
$$
f^\sharp :
\mathcal{O}_{X, x}
\longrightarrow
\mathcal{O}_{\Spec(R), \mathfrak m} = R.
$$

\begin{lemma}
\label{schemes-lemma-morphism-from-spec-local-ring}
设 $X$ 为概形，$R$ 为局部环。上述构造在态射 $\Spec(R) \to X$ 与二元组
$(x, \varphi)$ 之间给出一一对应；后二者由点 $x \in X$ 和局部环的局部同态
$\varphi : \mathcal{O}_{X, x} \to R$ 组成。
\end{lemma}

\begin{proof}
设 $A$ 为环。对任意环同态 $\psi : A \to R$，存在唯一素理想
$\mathfrak p \subset A$ 以及分解 $A \to A_{\mathfrak p} \to R$，
使得后一个映射是局部环的局部同态。事实上，
$\mathfrak p = \psi^{-1}(\mathfrak m)$。由引理
\ref{schemes-lemma-morphism-into-affine}，这就证明了当 $X$ 是仿射概形时的断言。

\medskip\noindent
现设 $X$ 为一般概形。任意 $x \in X$ 都包含在某个仿射开集
$U \subset X$ 中。由仿射情形可知，每个二元组 $(x, \varphi)$ 都由引理前的
构造产生。

\medskip\noindent
要完成证明，只须说明任意态射 $f : \Spec(R) \to X$ 的像都包含在每个
含有闭点像 $x$ 的仿射开集中；这里闭点属于 $\Spec(R)$。事实上，设
$x \in V \subset X$ 是任意含 $x$ 的开邻域。则
$f^{-1}(V) \subset \Spec(R)$ 是包含唯一闭点的开集，因而等于 $\Spec(R)$。
\end{proof}

\noindent
作为上面引理的特例，对每一点 $x$，若它属于概形 $X$，便得到典范态射
\begin{equation}
\label{schemes-equation-canonical-morphism}
\Spec(\mathcal{O}_{X, x}) \longrightarrow X
\end{equation}
它对应于 $X$ 在 $x$ 处局部环上的恒等映射。上面的引理也可以重述为：
对任意态射 $f : \Spec(R) \to X$，存在唯一点 $x \in X$，使得 $f$ 分解为
$\Spec(R) \to \Spec(\mathcal{O}_{X, x}) \to X$
其中第一个映射来自局部同态 $\mathcal{O}_{X, x} \to R$。

\medskip\noindent
若有概形态射 $f : X \to S$，且点 $x$ 映到点 $s \in S$，则得到交换图
$$
\xymatrix{
\Spec(\mathcal{O}_{X, x}) \ar[r] \ar[d] & X \ar[d] \\
\Spec(\mathcal{O}_{S, s}) \ar[r] & S
}
$$
其中左侧竖直映射对应于局部环同态
$f^\sharp_x : \mathcal{O}_{S, s} \to \mathcal{O}_{X, x}$。

\begin{lemma}
\label{schemes-lemma-specialize-points}
设 $X$ 为概形。设 $x, x' \in X$ 为 $X$ 的点。则 $x' \in X$ 是 $x$ 的泛化，
当且仅当 $x'$ 属于典范态射 $\Spec(\mathcal{O}_{X, x}) \to X$ 的像。
\end{lemma}

\begin{proof}
连续映射保持特化／泛化关系。由于 $\Spec(\mathcal{O}_{X, x})$ 的每个点
都是闭点的泛化，可见 $\Spec(\mathcal{O}_{X, x}) \to X$ 像中的每一点
都是 $x$ 的泛化。反过来，设 $x'$ 是 $x$ 的泛化。选取仿射开邻域
$U = \Spec(R)$，使它成为 $x$ 的邻域。则 $x' \in U$。设素理想
$\mathfrak p \subset R$ 与
$\mathfrak p' \subset R$ 分别对应于 $x$ 和 $x'$。由于 $x'$ 是 $x$ 的
泛化，可见 $\mathfrak p' \subset \mathfrak p$。这意味着 $\mathfrak p'$
属于态射
$\Spec(\mathcal{O}_{X, x}) = \Spec(R_{\mathfrak p})
\to \Spec(R) = U \subset X$ as desired.
\end{proof}

\noindent
现在讨论从域的谱出发的态射。设 $(R, \mathfrak m, \kappa)$ 为局部环，
其极大理想为 $\mathfrak m$，剩余域为 $\kappa$。设 $K$ 为域。按定义，
局部同态 $R \to K$ 分解为 $R \to \kappa \to K$；也就是说，它等同于态射
$\kappa \to K$。所以可见，态射
$$
\Spec(K) \longrightarrow X
$$
对应于二元组 $(x, \kappa(x) \to K)$。可以在从域的谱到 $X$ 的态射上
定义预序：若 $\Spec(K) \to X$ 通过 $\Spec(L) \to X$ 分解，
就说 $\Spec(K) \to X$ 支配 $\Spec(L) \to X$。这提示如下概念：暂且称
两个态射 $p : \Spec(K) \to X$ 与 $q : \Spec(L) \to X$ {\it 等价}，
如果存在第三个域 $\Omega$ 以及交换图
$$
\xymatrix{
\Spec(\Omega) \ar[r] \ar[d] &
\Spec(L) \ar[d]^q \\
\Spec(K) \ar[r]^p &
X
}
$$
当然，这立即说明三个概形 $\Spec(K)$、$\Spec(L)$ 和 $\Spec(\Omega)$ 的
唯一点都映到同一个 $x \in X$。因此，由上面的讨论，这样的图对应于一点
$x \in X$ 以及如下域的交换图：
$$
\xymatrix{
\Omega &
L \ar[l] \\
K \ar[u] &
\kappa(x) \ar[l] \ar[u]
}
$$
这定义了一个等价关系，因为给定任意一族域扩张 $K_i/\kappa$，都存在某个
域扩张 $\Omega/\kappa$，使得所有域扩张 $K_i$ 都包含在扩张 $\Omega$ 中。

\begin{lemma}
\label{schemes-lemma-characterize-points}
设 $X$ 为概形。$X$ 的点与从域的谱到 $X$ 的态射之等价类一一对应。
此外，每个等价类都包含一个最小元 $\Spec(\kappa(x)) \to X$，
它在唯一同构意义下唯一。
\end{lemma}

\begin{proof}
这来自上面的讨论。
\end{proof}

\noindent
当然，态射 $\Spec(\kappa(x)) \to X$ 通过典范态射
$\Spec(\mathcal{O}_{X, x}) \to X$ 分解。在当前语境下，引理
\ref{schemes-lemma-specialize-points} 的内容是：态射 $\Spec(\kappa(x')) \to X$
分解为
$\Spec(\kappa(x')) \to \Spec(\mathcal{O}_{X, x}) \to X$
只要 $x'$ 是 $x$ 的泛化。若有概形态射 $f : X \to S$，且点 $x$ 映到点
$s \in S$，则得到交换图
$$
\xymatrix{
\Spec(\kappa(x)) \ar[r] \ar[d] &
\Spec(\mathcal{O}_{X, x}) \ar[r] \ar[d] &
X \ar[d] \\
\Spec(\kappa(s)) \ar[r] &
\Spec(\mathcal{O}_{S, s}) \ar[r] &
S.
}
$$










\section{概形的粘合}
\label{schemes-section-glueing-schemes}

\noindent
设 $I$ 为集合。对每个 $i \in I$，设 $(X_i, \mathcal{O}_i)$ 为局部环层空间。
（事实上，下面的构造对环层空间同样适用。）对每一对 $i, j \in I$，设
$U_{ij} \subset X_i$ 为开子空间。对每一对 $i, j \in I$，设
$$
\varphi_{ij} : U_{ij} \to U_{ji}
$$
为局部环层空间的同构。为方便起见，假设 $U_{ii} = X_i$。对每个三元组
$i, j, k \in I$，假设：
\begin{enumerate}
\item 有
$\varphi_{ij}^{-1}(U_{ji} \cap U_{jk}) =  U_{ij} \cap U_{ik}$；
\item 下图
$$
\xymatrix{
U_{ij} \cap U_{ik} \ar[rr]_{\varphi_{ik}} \ar[rd]_{\varphi_{ij}} & &
U_{ki} \cap U_{kj} \\
& U_{ji} \cap U_{jk} \ar[ru]_{\varphi_{jk}}
}
$$
交换。
\end{enumerate}
（特别地，由此可得 $\varphi_{ii} = \text{id}_{X_i}$。）称一个族
$(I, (X_i)_{i\in I}, (U_{ij})_{i, j\in I}, (\varphi_{ij})_{i, j\in I})$
若满足上述条件，就是一组粘合数据。

\begin{lemma}
\label{schemes-lemma-glue}
\begin{slogan}
若有两个局部环层空间，并且第一个的一个子空间同构于另一个的一个子空间，
就可以把它们粘合成一个大的局部环层空间。
\end{slogan}
给定局部环层空间的任意粘合数据，存在局部环层空间 $X$、开子空间
$U_i \subset X$，以及局部环层空间的同构 $\varphi_i : X_i \to U_i$，使得
\begin{enumerate}
\item $X=\bigcup_{i\in I} U_i$,
\item $\varphi_i(U_{ij}) = U_i \cap U_j$, and
\item $\varphi_{ij} =
\varphi_j^{-1}|_{U_i \cap U_j} \circ \varphi_i|_{U_{ij}}$.
\end{enumerate}
局部环层空间 $X$ 由如下映射性质刻画：给定局部环层空间 $Y$，有
\begin{eqnarray*}
\Mor(X, Y) & = & \{ (f_i)_{i\in I} \mid
f_i : X_i \to Y, \ f_j \circ \varphi_{ij} = f_i|_{U_{ij}}\} \\
f & \mapsto & (f|_{U_i} \circ \varphi_i)_{i \in I} \\
\Mor(Y, X) & = &
\left\{
\begin{matrix}
\text{开覆盖 }Y = \bigcup\nolimits_{i \in I} V_i\text{ 且 }
(g_i : V_i \to X_i)_{i \in I}
\text{ 使得}\\
g_i^{-1}(U_{ij}) = V_i \cap V_j
\text{ 且 }
g_j|_{V_i \cap V_j} = \varphi_{ij} \circ g_i|_{V_i \cap V_j}
\end{matrix}
\right\} \\
g & \mapsto &
V_i = g^{-1}(U_i), \ g_i = \varphi_i^{-1} \circ g|_{V_i}
\end{eqnarray*}
\end{lemma}

\begin{proof}
分阶段构造 $X$。作为集合，取
$$
X = (\coprod X_i) / \sim.
$$
这里，给定 $x \in X_i$ 与 $x' \in X_j$，规定 $x \sim x'$ 当且仅当
$x \in U_{ij}$、$x' \in U_{ji}$ 且 $\varphi_{ij}(x) = x'$。这是等价关系：
若 $x \in X_i$、$x' \in X_j$、$x'' \in X_k$，并且 $x \sim x'$、
$x' \sim x''$，则 $x' \in U_{ji} \cap U_{jk}$；于是由粘合数据的条件 (1)，
还有 $x \in U_{ij} \cap U_{ik}$ 与 $x'' \in U_{ki} \cap U_{kj}$，再由条件
(2) 得到 $\varphi_{ik}(x) = x''$。（自反性与对称性来自假设
$U_{ii} = X_i$ 和 $\varphi_{ii} = \text{id}_{X_i}$。）把自然映射记作
$\varphi_i : X_i \to X$，并记 $U_i = \varphi_i(X_i) \subset X$。
注意，$\varphi_i : X_i \to U_i$ 是双射。

\medskip\noindent
$X$ 上的拓扑按如下规则定义：$U \subset X$ 为开集，当且仅当
$\varphi_i^{-1}(U)$ 对所有 $i$ 都是开集。留给读者验证，这确实定义了拓扑。
特别地，注意 $U_i$ 是开集，因为 $\varphi_j^{-1}(U_i)
= U_{ji}$ 是 $X_j$ 中的开集；这对所有 $j$ 都成立。此外，对任意开集
$W \subset X_i$，其像
$\varphi_i(W) \subset U_i$ 是开集，因为
$\varphi_j^{-1}(\varphi_i(W)) = \varphi_{ji}^{-1}(W \cap U_{ij})$。
因此，$\varphi_i : X_i \to U_i$ 是同胚。

\medskip\noindent
要得到局部环层空间，还须构造环层 $\mathcal{O}_X$。我们通过粘合各环层
$\mathcal{O}_{U_i} := \varphi_{i, *} \mathcal{O}_i$ 来完成。具体而言，
在交换图
$$
\xymatrix{
U_{ij} \ar[rr]_{\varphi_{ij}} \ar[rd]_{\varphi_i|_{U_{ij}}}
& &
U_{ji} \ar[ld]^{\varphi_j|_{U_{ji}}} \\
& U_i \cap U_j &
}
$$
中，上方箭头是环层空间的同构，因而得到唯一的环层同构
$$
\mathcal{O}_{U_i}|_{U_i \cap U_j}
\longrightarrow
\mathcal{O}_{U_j}|_{U_i \cap U_j}.
$$
它们满足《层》第 \ref{sheaves-section-glueing-sheaves} 节中的余循环条件。
由该节的结果，得到环层 $\mathcal{O}_X$；它定义在 $X$ 上，并使得
$\mathcal{O}_X|_{U_i}$ 与 $\mathcal{O}_{U_i}$ 同构，并与上面展示的粘合映射
相容。特别地，$(X, \mathcal{O}_X)$ 是局部环层空间，因为
$\mathcal{O}_X$ 的茎等于对应点处 $\mathcal{O}_i$ 的茎。

\medskip\noindent
映射性质的证明略去。
\end{proof}

\begin{lemma}
\label{schemes-lemma-glue-schemes}
\begin{slogan}
概形可以粘合成新的概形。
\end{slogan}
在上面的引理 \ref{schemes-lemma-glue} 中，假设所有 $X_i$ 都是概形。
则所得局部环层空间 $X$ 是概形。
\end{lemma}

\begin{proof}
这是显然的，因为每个 $U_i$ 都是概形，因而每个 $x \in X$ 都有仿射邻域。
\end{proof}

\noindent
通常把 $X_i$ 看成 $X$ 的开子空间；所借助的是同构 $\varphi_i$。
下面两个例子中都采用这一约定。

\begin{example}[零点加倍的仿射空间]
\label{schemes-example-affine-space-zero-doubled}
设 $k$ 为域，$n \geq 1$。令 $X_1 = \Spec(k[x_1, \ldots, x_n])$，
令 $X_2 = \Spec(k[y_1, \ldots, y_n])$。令 $0_1 \in X_1$ 为对应于极大理想
$(x_1, \ldots, x_n) \subset k[x_1, \ldots, x_n]$ 的点。令 $0_2 \in X_2$
为对应于极大理想 $(y_1, \ldots, y_n) \subset k[y_1, \ldots, y_n]$ 的点。
令 $U_{12} = X_1 \setminus \{0_1\}$，并令
$U_{21} = X_2 \setminus \{0_2\}$。设 $\varphi_{12} : U_{12} \to U_{21}$
为来自 $k$-代数同构
$k[y_1, \ldots, y_n] \to k[x_1, \ldots, x_n]$ 的同构；该同构把 $y_i$
映到 $x_i$（并诱导出 $X_1 \cong X_2$，把 $0_1$ 映到 $0_2$）。
设 $X$ 为由粘合数据
$(X_1, X_2, U_{12}, U_{21}, \varphi_{12},
\varphi_{21} = \varphi_{12}^{-1})$
所得的概形。按照例子上方引入的轻微记号滥用，把 $X_1, X_2 \subset X$
看成开子概形。存在态射 $f : X \to \Spec(k[t_1, \ldots, t_n])$，
它在 $X_1$ 上（相应地，在 $X_2$ 上）对应于 $k$-代数同态
$k[t_1, \ldots, t_n] \to k[x_1, \ldots, x_n]$
(resp.\ $k[t_1, \ldots, t_n] \to k[y_1, \ldots, y_n]$)
它把 $t_i$ 映到 $x_i$（相应地，把 $t_i$ 映到 $y_i$）。容易看出，
该态射把 $k[t_1, \ldots, t_n]$ 与 $\Gamma(X, \mathcal{O}_X)$ 认同。
由于 $f(0_1) = f(0_2)$，可见 $X$ 不是仿射的。

\medskip\noindent
注意，$X_1$ 与 $X_2$ 都是 $X$ 的仿射开集。但若 $n = 2$，则
$X_1 \cap X_2$ 是例 \ref{schemes-example-not-affine} 中描述的概形，因而不是仿射的。
所以一般而言，概形的仿射开集之交不再仿射。（更一般地，该事实对任意
$n > 1$ 都成立。）

\medskip\noindent
该例还有如下奇特性质。若 $n > 1$，则存在许多不可约闭子集 $T \subset X$
（例如取 $X_1$ 中任意非闭点的闭包）。但除非 $T = \{0_1\}$ 或
$T = \{0_2\}$，总有 $0_1 \in T \Leftrightarrow 0_2 \in T$。证明略去。
\end{example}

\begin{example}[射影直线]
\label{schemes-example-projective-line}
设 $k$ 为域。令 $X_1 = \Spec(k[x])$，令 $X_2 = \Spec(k[y])$。
令 $0 \in X_1$ 为对应于极大理想 $(x) \subset k[x]$ 的点。
令 $\infty \in X_2$ 为对应于极大理想 $(y) \subset k[y]$ 的点。
令 $U_{12} = X_1 \setminus \{0\} = D(x) = \Spec(k[x, 1/x])$，并令
$U_{21} = X_2 \setminus \{\infty\} = D(y) = \Spec(k[y, 1/y])$。
设 $\varphi_{12} : U_{12} \to U_{21}$ 为来自 $k$-代数同构
$k[y, 1/y] \to k[x, 1/x]$ 的同构；该同构把 $y$ 映到 $1/x$。
设 $\mathbf{P}^1_k$ 为由粘合数据
$(X_1, X_2, U_{12}, U_{21}, \varphi_{12},
\varphi_{21} = \varphi_{12}^{-1})$
所得的概形。按照例子上方引入的轻微记号滥用，把
$X_i \subset \mathbf{P}^1_k$ 看成开子概形。在此情形下，
$\Gamma(\mathbf{P}^1_k, \mathcal{O}) = k$，因为满足如下条件的多项式
$g(x)$ 只有常数多项式：它是 $x$ 中的多项式，$g(1/y)$ 也是 $y$ 中的
多项式。由于
$\mathbf{P}^1_k$ 是无限的，可见 $\mathbf{P}^1_k$ 不是仿射的。

\medskip\noindent
我们断言存在仿射开集 $U \subset \mathbf{P}^1_k$，它同时包含 $0$ 与
$\infty$。事实上，令 $U = \mathbf{P}^1_k \setminus \{1\}$，其中 $1$ 是
$X_1$ 中对应于极大理想 $(x - 1)$ 的点，也是 $X_2$ 中对应于极大理想
$(y - 1)$ 的点。容易看出
$s = 1/(x - 1) = y/(1 - y) \in \Gamma(U, \mathcal{O}_U)$。事实上，可以证明
$\Gamma(U, \mathcal{O}_U)$ 等于多项式环 $k[s]$，且相应态射
$U \to \Spec(k[s])$ 是概形同构。细节略去。
\end{example}




\section{一个可表性判据}
\label{schemes-section-representable}

\noindent
本节用函子的语言重新表述第 \ref{schemes-section-glueing-schemes} 节的粘合引理。
先回顾《范畴》第 \ref{categories-section-opposite} 节的一些材料。
回顾，给定概形 $X$，可以定义函子
$$
h_X : \Sch^{opp}
\longrightarrow
\textit{Sets}, \quad
T \longmapsto \Mor(T, X).
$$
它称为 $X$ 的{\it 点函子}。

\medskip\noindent
设 $F$ 为从概形范畴到集合范畴的反变函子。写成公式即
$$
F : \Sch^{opp}
\longrightarrow
\textit{Sets}.
$$
我们采用《位点》第 \ref{sites-section-presheaves} 节中的相同术语。
具体地，给定概形 $T$、元素 $\xi \in F(T)$ 以及态射 $f : T' \to T$，
用 $f^*\xi$ 表示元素 $F(f)(\xi)$，有时甚至采用记号 $\xi|_{T'}$。

\begin{definition}
\label{schemes-definition-representable-functor}
（见《范畴》的定义 \ref{categories-definition-representable-functor}。）
设 $F$ 为从概形范畴到集合范畴的反变函子（如上）。称 $F$
{\it 由概形表示}，或称它{\it 可表}，如果存在概形 $X$，使得
$h_X \cong F$。
\end{definition}

\noindent
设 $F$ 由概形 $X$ 表示，且 $s : h_X \to F$ 是同构。由《范畴》的米田引理
\ref{categories-lemma-yoneda}，二元组 $(X, s : h_X \to F)$ 若存在，
就在唯一同构意义下唯一。此外，米田引理说明，对上述任意反变函子 $F$
和任意概形 $Y$，都有双射
$$
\Mor_{\text{Fun}(\Sch^{opp}, \textit{Sets})} (h_Y, F)
\longrightarrow
F(Y), \quad
s \longmapsto s(\text{id}_Y).
$$
反向构造如下：给定任意 $\xi \in F(Y)$，函子变换
$s_\xi : h_Y \to F$ 把任意态射 $f : T \to Y$ 送到元素
$f^*\xi \in F(T)$。

\medskip\noindent
特别地，当 $F$ 可表时，存在概形 $X$ 和元素 $\xi \in F(X)$，使得相应态射
$h_X \to F$ 是同构。在这种情形下，也说{\it 二元组 $(X, \xi)$ 表示 $F$}。
元素 $\xi \in F(X)$ 常称为{\it “泛族”}；等到讨论代数栈时，这一名称的
理由会更加清楚（此处将来插入参考）。目前只须注意：若二元组 $(X, \xi)$
表示 $F$，则每个元素 $\xi' \in F(T)$（其中 $T$ 任意）都形如
$\xi' = f^*\xi$，其中 $f : T \to X$ 唯一。

\begin{example}
\label{schemes-example-global-sections}
考虑如下规则：把每个概形 $T$ 送到集合
$F(T) = \Gamma(T, \mathcal{O}_T)$。
对态射 $f : T' \to T$ 使用拉回映射
$f^\sharp : \Gamma(T, \mathcal{O}_T) \to \Gamma(T', \mathcal{O}_{T'})$，
即可把它变成反变函子。给定环 $R$ 和元素 $t \in R$，存在唯一环同态
$\mathbf{Z}[x] \to R$，把 $x$ 映到 $t$。因此，利用引理
\ref{schemes-lemma-morphism-into-affine}，可见
$$
\Mor(T, \Spec(\mathbf{Z}[x])) =
\Hom(\mathbf{Z}[x], \Gamma(T, \mathcal{O}_T)) =
\Gamma(T, \mathcal{O}_T).
$$
这确实给出同构 $h_{\Spec(\mathbf{Z}[x])} \to F$。“泛族” $\xi$ 是什么？
要得到它，必须把上面的认同应用于
$\text{id}_{\Spec(\mathbf{Z}[x])}$。显然，在这些认同下得到
$\xi = x \in \Gamma(\Spec(\mathbf{Z}[x]),
\mathcal{O}_{\Spec(\mathbf{Z}[x])}) = \mathbf{Z}[x]$
，正如预期。
\end{example}

\begin{definition}
\label{schemes-definition-representable-by-open-immersions}
设 $F$ 为概形范畴上取值于集合的反变函子。
\begin{enumerate}
\item 称 $F$ {\it 满足 Zariski 拓扑的层性质}，如果对每个概形 $T$、
每个开覆盖 $T = \bigcup_{i \in I} U_i$，
以及任意一族元素 $\xi_i \in F(U_i)$，只要
$\xi_i|_{U_i \cap U_j} =
\xi_j|_{U_i \cap U_j}$，就存在唯一元素
$\xi \in F(T)$，使得 $\xi_i = \xi|_{U_i}$ 在 $F(U_i)$ 中成立。
\item {\it 子函子 $H \subset F$} 是如下规则：对每个概形 $T$，赋予子集
$H(T) \subset F(T)$，使得映射 $F(f) : F(T) \to F(T')$ 把 $H(T)$
映入 $H(T')$；这对所有概形态射 $f : T' \to T$ 都成立。
\item 设 $H \subset F$ 为子函子。称 $H \subset F$ {\it 由开浸入表示}，
如果对所有二元组 $(T, \xi)$，其中 $T$ 是概形且 $\xi \in F(T)$，
都存在具有如下性质的开子概形 $U_\xi \subset T$：
\begin{itemize}
\item[(*)] 态射 $f : T' \to T$ 通过 $U_\xi$ 分解，当且仅当
$f^*\xi \in H(T')$。
\end{itemize}
\item 设 $I$ 为集合。对每个 $i \in I$，设 $H_i \subset F$ 为子函子。
称族 $(H_i)_{i \in I}$ {\it 覆盖 $F$}，当且仅当对每个
$\xi \in F(T)$，都存在开覆盖 $T = \bigcup U_i$，使得
$\xi|_{U_i} \in H_i(U_i)$。
\end{enumerate}
\end{definition}

\noindent
在条件 (4) 中，若 $H_i \subset F$ 对所有 $i$ 都由开浸入表示，
则要检验 $(H_i)_{i \in I}$ 覆盖 $F$，只须检验
$F(T) = \bigcup H_i(T)$；这里 $T$ 是域的谱。

\begin{lemma}
\label{schemes-lemma-glue-functors}
设 $F$ 为概形范畴上取值于集合范畴的反变函子。假设：
\begin{enumerate}
\item $F$ 满足 Zariski 拓扑的层性质；
\item 存在集合 $I$ 以及一族子函子 $F_i \subset F$，使得
\begin{enumerate}
\item 每个 $F_i$ 可表；
\item 每个 $F_i \subset F$ 由开浸入表示；
\item 族 $(F_i)_{i \in I}$ 覆盖 $F$。
\end{enumerate}
\end{enumerate}
则 $F$ 可表。
\end{lemma}

\begin{proof}
令 $X_i$ 为表示 $F_i$ 的概形，并令
$\xi_i \in F_i(X_i) \subset F(X_i)$ 为“泛族”。
由于 $F_j \subset F$ 由开浸入表示，
存在开子集 $U_{ij} \subset X_i$，使得
$T \to X_i$ 经由 $U_{ij}$ 分解，当且仅当
$\xi_i|_T \in F_j(T)$。特别地，
$\xi_i|_{U_{ij}} \in F_j(U_{ij})$，因而得到典范态射
$\varphi_{ij} : U_{ij} \to X_j$，满足
$\varphi_{ij}^*\xi_j = \xi_i|_{U_{ij}}$。按照 $U_{ji}$ 的定义，
这表明 $\varphi_{ij}$ 经由 $U_{ji}$ 分解。
由于 $(\varphi_{ij} \circ \varphi_{ji})^*\xi_j
=\varphi_{ji}^*(\varphi_{ij}^*\xi_j) =
\varphi_{ji}^*\xi_i = \xi_j$，我们断定
$\varphi_{ij} \circ \varphi_{ji} = \text{id}_{U_{ji}}$，
这是因为二元组 $(X_j, \xi_j)$ 表示 $F_j$。
特别地，映射 $\varphi_{ij} : U_{ij} \to U_{ji}$
都是概形同构。
接着须证明
$\varphi_{ij}^{-1}(U_{ji} \cap U_{jk}) = U_{ij} \cap U_{ik}$。
这是因为：(a) $U_{ji} \cap U_{jk}$ 是 $U_{ji}$ 中使
$\xi_j$ 限制为 $F_k$ 的一个元素的最大开子集；
(b) $U_{ij} \cap U_{ik}$ 是 $U_{ij}$ 中使
$\xi_i$ 限制为 $F_k$ 的一个元素的最大开子集；
(c) $\varphi_{ij}^*\xi_j = \xi_i$。此外，
第 \ref{schemes-section-glueing-schemes} 节中的余圈条件成立，因为
$\varphi_{jk}|_{U_{ji} \cap U_{jk}} \circ
\varphi_{ij}|_{U_{ij} \cap U_{ik}}$ 与
$\varphi_{ik}|_{U_{ij} \cap U_{ik}}$ 都把 $\xi_k$
拉回为元素 $\xi_i$。
因此可以应用引理 \ref{schemes-lemma-glue-schemes}，
得到概形 $X$、开覆盖
$X = \bigcup U_i$，以及同构 $\varphi_i : X_i \to U_i$，
它们具有引理 \ref{schemes-lemma-glue} 中的性质。
令 $\xi_i' = (\varphi_i^{-1})^* \xi_i$。
引理 \ref{schemes-lemma-glue} 的条件蕴含
$\xi_i'|_{U_i \cap U_j} = \xi_j'|_{U_i \cap U_j}$。
因此，由于 $F$ 满足 Zariski 拓扑的层条件，
可知存在元素
$\xi' \in F(X)$，使得
$\xi_i = \varphi_i^*\xi'|_{U_i}$ 对所有 $i$ 都成立。
由于 $\varphi_i$ 是同构，还可得到
$(U_i, \xi'|_{U_i})$ 表示函子 $F_i$。

\medskip\noindent
我们声称二元组 $(X, \xi')$ 表示函子 $F$。
为证此点，令 $T$ 为概形，并令 $\xi \in F(T)$。
我们将构造唯一态射 $g : T \to X$，使得
$g^*\xi' = \xi$。确切地说，由子函子
$F_i$ 覆盖 $F$ 这一条件，存在开覆盖 $T = \bigcup V_i$，
使得对每个 $i$，限制 $\xi|_{V_i} \in F_i(V_i)$。
此外，由于每个包含 $F_i \subset F$ 都由开浸入表示，
可假设每个 $V_i \subset T$ 都是具有此性质的最大开子集。
由于 $(U_i, \xi'|_{U_i})$ 表示函子 $F_i$，
得到唯一态射 $g_i : V_i \to U_i$，使得
$g_i^*\xi'|_{U_i} = \xi|_{V_i}$。在交集 $V_i \cap V_j$ 上，
态射 $g_i$ 与 $g_j$ 相同；例如，这是因为它们都把
$\xi'|_{U_i \cap U_j} \in F_i(U_i \cap U_j)$
拉回为同一个元素。因此态射 $g_i$ 粘合成唯一态射
$T \to X$，正合所需。
\end{proof}

\begin{remark}
\label{schemes-remark-representable-locally-ringed}
设函子 $F$ 定义在所有局部环层空间上，
并把引理 \ref{schemes-lemma-glue-functors} 的条件替换为以下条件：
\begin{enumerate}
\item $F$ 在局部环层空间范畴上满足层性质；
\item 存在集合 $I$ 以及一族子函子
$F_i \subset F$，使得
\begin{enumerate}
\item 每个 $F_i$ 都由概形表示；
\item 每个 $F_i \subset F$ 在局部环层空间范畴上都由开浸入表示；
\item 族 $(F_i)_{i \in I}$ 在局部环层空间范畴上的函子意义下覆盖 $F$。
\end{enumerate}
\end{enumerate}
我们把进一步写出细节留给读者。
最终结论是：函子 $F$ 在局部环层空间范畴中可表，
并且表示对象是一个概形。
\end{remark}



\section{概形纤维积的存在性}
\label{schemes-section-existence-fibre-products}

\noindent
一个非常基本的问题是：概形范畴中是否存在乘积与纤维积。
我们先抽象地证明乘积与纤维积存在，下一节再说明如何以一种
合理的方式理解概形的纤维积。

\begin{lemma}
\label{schemes-lemma-fibre-products}
概形范畴具有终对象、乘积和纤维积。
换言之，概形范畴具有有限极限；见《范畴》，
引理 \ref{categories-lemma-finite-limits-exist}。
\end{lemma}

\begin{proof}
请跳过这个证明。更重要的是学习如何使用下一节所说明的纤维积。

\medskip\noindent
由引理 \ref{schemes-lemma-morphism-into-affine}，
概形 $\Spec(\mathbf{Z})$ 是局部环层空间范畴中的终对象。
因此只须证明纤维积存在。

\medskip\noindent
令 $f : X \to S$ 与 $g : Y \to S$ 为概形态射。须证明函子
\begin{eqnarray*}
F : \Sch^{opp} & \longrightarrow & \textit{Sets} \\
T & \longmapsto &
\Mor(T, X) \times_{\Mor(T, S)} \Mor(T, Y)
\end{eqnarray*}
可表。我们声称引理 \ref{schemes-lemma-glue-functors}
适用于函子 $F$。证明此点便证得本引理。

\medskip\noindent
首先证明 $F$ 满足 Zariski 拓扑的层性质。具体而言，设 $T$ 为概形，
$T = \bigcup_{i \in I} U_i$ 为开覆盖，并且
$\xi_i \in F(U_i)$ 满足
$\xi_i|_{U_i \cap U_j} =  \xi_j|_{U_i \cap U_j}$；这对所有二元组
$i, j$ 都成立。按照定义，$\xi_i$ 对应于二元组 $(a_i, b_i)$，其中
$a_i : U_i \to X$ 与 $b_i : U_i \to Y$ 是概形态射，满足
$f \circ a_i = g \circ b_i$。粘合条件说明
$a_i|_{U_i \cap U_j} =  a_j|_{U_i \cap U_j}$
and
$b_i|_{U_i \cap U_j} =  b_j|_{U_i \cap U_j}$.
因此，粘合态射 $a_i$ 得到一个局部环层空间态射
（即概形态射）$a : T \to X$；同样得到 $b : T \to Y$
（例如见引理 \ref{schemes-lemma-glue} 的映射性质）。此外，在某个开覆盖的
各成员上，复合 $f \circ a$ 与 $g \circ b$ 相同。因此
$f \circ a = g \circ b$，二元组 $(a, b)$ 定义了 $F(T)$ 的一个元素，
它限制为二元组 $(a_i, b_i)$，在每个 $U_i$ 上皆如此。层条件得证。

\medskip\noindent
接着构造子函子族。选取由仿射开子集构成的开覆盖
$S = \bigcup\nolimits_{i \in I} U_i$。
对每个 $i \in I$，选取由仿射开子集构成的开覆盖
$f^{-1}(U_i) = \bigcup\nolimits_{j \in J_i} V_j$，以及
$g^{-1}(U_i) = \bigcup\nolimits_{k \in K_i} W_k$。
注意，$X = \bigcup_{i \in I} \bigcup_{j \in J_i} V_j$
是开覆盖；$Y$ 也同样。对任意 $i \in I$ 以及每一对
$(j, k) \in J_i \times K_i$，都有交换图
$$
\xymatrix{
    & W_k \ar[d] \ar[rd] &   \\
V_j \ar[rd] \ar[r] & U_i \ar[rd] & Y \ar[d] \\
    & X \ar[r]  & S
}
$$
其中所有斜箭头都是开浸入。对这样的三元组，得到函子
\begin{eqnarray*}
F_{i, j, k} : \Sch^{opp} & \longrightarrow & \textit{Sets} \\
T & \longmapsto &
\Mor(T, V_j) \times_{\Mor(T, U_i)} \Mor(T, W_k).
\end{eqnarray*}
存在显然的函子变换 $F_{i, j, k} \to F$
（来自上面的巨大交换图），它是单射，因此可把 $F_{i, j, k}$
视为 $F$ 的子函子。

\medskip\noindent
检验引理 \ref{schemes-lemma-glue-functors} 的条件 (2)(a)。
这直接由引理 \ref{schemes-lemma-fibre-product-affine-schemes} 得出。
（注意，这里用到了：仿射概形范畴中的纤维积也是整个局部环层空间
范畴中的纤维积。）

\medskip\noindent
检验引理 \ref{schemes-lemma-glue-functors} 的条件 (2)(b)。
令 $T$ 为概形，并令 $\xi \in F(T)$。换言之，
$\xi = (a, b)$，其中 $a : T \to X$ 与 $b : T \to Y$ 是概形态射，
满足 $f \circ a = g \circ b$。
令 $V_{i, j, k} = a^{-1}(V_j) \cap b^{-1}(W_k)$。对任意另一个态射
$h : T' \to T$，都有
$h^*\xi = (a \circ h, b \circ h)$。因此可知，
$h^*\xi \in F_{i, j, k}(T')$ 当且仅当
$a(h(T')) \subset V_j$ 且 $b(h(T')) \subset W_k$。
换言之，当且仅当 $h(T') \subset V_{i, j, k}$。这证明了条件 (2)(b)。

\medskip\noindent
检验引理 \ref{schemes-lemma-glue-functors} 的条件 (2)(c)。
令 $T$ 为概形，并令 $\xi = (a, b) \in F(T)$ 如上。
仍令 $V_{i, j, k} = a^{-1}(V_j) \cap b^{-1}(W_k)$。
条件 (2)(c) 正是说 $T = \bigcup V_{i, j, k}$，这显然成立。
因此引理得证，纤维积存在。
\end{proof}

\begin{remark}
\label{schemes-remark-fibre-product-schemes-locally-ringed}
利用注 \ref{schemes-remark-representable-locally-ringed}，
可以证明：概形态射的纤维积在局部环层空间范畴中存在，
并且是概形。
\end{remark}




\section{概形的纤维积}
\label{schemes-section-fibre-products}

\noindent
尽管已经证明了概形的纤维积存在，这里仍回顾其一般定义。

\begin{definition}
\label{schemes-definition-fibre-product}
给定概形态射 $f : X \to S$ 与 $g : Y \to S$，
{\it 纤维积}是概形 $X \times_S Y$ 连同投影态射
$p : X \times_S Y \to X$ 与 $q : X \times_S Y \to Y$，
它们处于下列交换图中：
$$
\xymatrix{
X \times_S Y \ar[r]_q \ar[d]_p & Y \ar[d]^g \\
X \ar[r]^f & S
}
$$
并且此图在所有这类图中具有泛性质；见《范畴》，
定义 \ref{categories-definition-fibre-products}。
\end{definition}

\noindent
换言之，给定任意由概形态射组成的实线交换图
$$
\xymatrix{
T \ar[rrrd] \ar@{-->}[rrd] \ar[rrdd]
&
&
\\
&
&
X \times_S Y \ar[d] \ar[r]
&
Y \ar[d]
\\
&
&
X \ar[r]
&
S
}
$$
存在唯一的虚线箭头使该图交换。我们将证明若干引理，
以说明应如何理解纤维积。

\begin{lemma}
\label{schemes-lemma-fibre-product-affines}
令 $f : X \to S$ 与 $g : Y \to S$ 为具有相同目标的概形态射。
若 $X, Y, S$ 都是仿射的，则 $X \times_S Y$ 是仿射的。
\end{lemma}

\begin{proof}
设 $X = \Spec(A)$、$Y = \Spec(B)$，且
$S = \Spec(R)$。由引理 \ref{schemes-lemma-fibre-product-affine-schemes}，
仿射概形 $\Spec(A \otimes_R B)$ 是局部环层空间范畴中的纤维积
$X \times_S Y$。因而它当然也是概形范畴中的纤维积。
\end{proof}

\begin{lemma}
\label{schemes-lemma-open-fibre-product}
令 $f : X \to S$ 与 $g : Y \to S$ 为具有相同目标的概形态射。
令 $X \times_S Y$、$p$、$q$ 为该纤维积。设
$U \subset S$，$V \subset X$、$W \subset Y$ 是开子概形，满足
$f(V) \subset U$ 且 $g(W) \subset U$。则典范态射
$V \times_U W \to X \times_S Y$ 是开浸入，它把
$V \times_U W$ 等同于 $p^{-1}(V) \cap q^{-1}(W)$。
\end{lemma}

\begin{proof}
令 $T$ 为概形。设 $a : T \to V$ 与 $b : T \to W$ 为态射，
并且有 $f \circ a = g \circ b$，这里把二者视为到 $U$ 的态射。
那么它们作为到 $S$ 的态射也相同。由纤维积的泛性质，
得到唯一态射 $T \to X \times_S Y$。当然，该态射的像包含在开集
$p^{-1}(V) \cap q^{-1}(W)$ 中。因此
$p^{-1}(V) \cap q^{-1}(W)$ 是 $V$ 与 $W$ 在 $U$ 上的纤维积。
结论由纤维积的唯一性得出；见《范畴》，
第 \ref{categories-section-fibre-products} 节。
\end{proof}

\noindent
特别地，这说明在该引理的情形下
$V \times_U W = V \times_S W$。此外，若 $U, V, W$ 都是仿射的，
则已知 $V \times_U W$ 是仿射的。当然，可以覆盖 $X \times_S Y$，
所用的是这样的仿射开集 $V \times_U W$。将此表述为一个引理。

\begin{lemma}
\label{schemes-lemma-affine-covering-fibre-product}
\begin{slogan}
纤维积的直接构造：概形纤维积的仿射开覆盖可以由各部分相容的
仿射开覆盖拼装而成。
\end{slogan}
令 $f : X \to S$ 与 $g : Y \to S$ 为具有相同目标的概形态射。
令 $S = \bigcup U_i$ 为 $S$ 的任意仿射开覆盖。对每个 $i \in I$，
令 $f^{-1}(U_i) = \bigcup_{j \in J_i} V_j$ 为 $f^{-1}(U_i)$ 的仿射开覆盖，
并令 $g^{-1}(U_i) = \bigcup_{k \in K_i} W_k$ 为 $g^{-1}(U_i)$ 的仿射开覆盖。
则
$$
X \times_S Y =
\bigcup\nolimits_{i \in I}
\bigcup\nolimits_{j \in J_i, \ k \in K_i}
V_j \times_{U_i} W_k
$$
是 $X \times_S Y$ 的仿射开覆盖。
\end{lemma}

\begin{proof}
见本引理上方的讨论。
\end{proof}

\noindent
换言之，我们本可以利用前一个引理，通过粘合仿射概形直接构造纤维积。
（当然，这恰好正是引理 \ref{schemes-lemma-fibre-products} 的证明中所做的。）
下面给出描述概形纤维积点集的一种方式。

\begin{lemma}
\label{schemes-lemma-points-fibre-product}
令 $f : X \to S$ 与 $g : Y \to S$ 为具有相同目标的概形态射。
点 $z$（视为 $X \times_S Y$ 的点）与下列四元组一一对应：
$$
(x, y, s, \mathfrak p)
$$
其中 $x \in X$、$y \in Y$、$s \in S$ 是满足
$f(x) = s$、$g(y) = s$ 的点，而 $\mathfrak p$ 是环
$\kappa(x) \otimes_{\kappa(s)} \kappa(y)$ 的素理想。
$z$ 的剩余域对应于素理想 $\mathfrak p$ 的剩余域。
\end{lemma}

\begin{proof}
令 $z$ 为 $X \times_S Y$ 的一个点，并构造上述四元组。
回忆可以把 $z$ 看成态射
$\Spec(\kappa(z)) \to X \times_S Y$；见引理
\ref{schemes-lemma-characterize-points}。此态射对应于态射
$a : \Spec(\kappa(z)) \to X$ 与
$b : \Spec(\kappa(z)) \to Y$，它们满足
$f \circ a = g \circ b$。再次应用同一引理，得到点
$x \in X$、$y \in Y$，它们位于同一点 $s \in S$ 上方；还得到域映射
$\kappa(x) \to \kappa(z)$、$\kappa(y) \to \kappa(z)$，使得复合
$\kappa(s) \to \kappa(x) \to \kappa(z)$ 与
$\kappa(s) \to \kappa(y) \to \kappa(z)$ 相同。换言之，得到环映射
$\kappa(x) \otimes_{\kappa(s)} \kappa(y) \to \kappa(z)$。
令 $\mathfrak p$ 为此映射的核。

\medskip\noindent
反过来，给定四元组 $(x, y, s, \mathfrak p)$，得到实线交换图
$$
\xymatrix{
X \times_S Y
\ar@/_/[dddr] \ar@/^/[rrrd]
& & & \\
&
\Spec(\kappa(x) \otimes_{\kappa(s)} \kappa(y)/\mathfrak p)
\ar[r] \ar[d] \ar@{-->}[lu]
&
\Spec(\kappa(y)) \ar[d] \ar[r] &
Y \ar[dd] \\
&
\Spec(\kappa(x)) \ar[r] \ar[d] &
\Spec(\kappa(s)) \ar[rd] &
\\
&
X \ar[rr] &
&
S
}
$$
见第 \ref{schemes-section-points} 节中的讨论。因而得到虚线箭头。
相应的点 $z$（作为 $X \times_S Y$ 的点）是下列概形泛点的像：
$\Spec(\kappa(x) \otimes_{\kappa(s)} \kappa(y)/\mathfrak p)$。
我们略去验证这两个构造互逆的过程。
\end{proof}

\begin{lemma}
\label{schemes-lemma-fibre-product-immersion}
令 $f : X \to S$ 与 $g : Y \to S$ 为具有相同目标的概形态射。
\begin{enumerate}
\item 若 $f : X \to S$ 是闭浸入，则
$X \times_S Y \to Y$ 是闭浸入。此外，若 $X \to S$ 对应于拟凝聚理想层
$\mathcal{I} \subset \mathcal{O}_S$，则 $X \times_S Y \to Y$ 对应于理想层
$\Im(g^*\mathcal{I} \to \mathcal{O}_Y)$。
\item 若 $f : X \to S$ 是开浸入，则
$X \times_S Y \to Y$ 是开浸入。
\item 若 $f : X \to S$ 是浸入，则
$X \times_S Y \to Y$ 是浸入。
\end{enumerate}
\end{lemma}

\begin{proof}
假设 $X \to S$ 是对应于拟凝聚理想层
$\mathcal{I} \subset \mathcal{O}_S$ 的闭浸入。由引理
\ref{schemes-lemma-restrict-map-to-closed}，闭子空间 $Z \subset Y$，
它由理想层 $\Im(g^*\mathcal{I} \to \mathcal{O}_Y)$ 定义，
是局部环层空间范畴中的纤维积。
由引理 \ref{schemes-lemma-closed-subspace-scheme}，$Z$ 是概形。
因此 $Z = X \times_S Y$，第一项得证。例如，第二项由引理
\ref{schemes-lemma-open-fibre-product} 得出。第三项是前两项的结合。
\end{proof}

\begin{definition}
\label{schemes-definition-inverse-image-closed-subscheme}
令 $f : X \to Y$ 为概形态射。令 $Z \subset Y$ 为 $Y$ 的闭子概形。
所谓{\it 逆像 $f^{-1}(Z)$，即闭子概形 $Z$ 的逆像}，是闭子概形
$Z \times_Y X$（视为 $X$ 的闭子概形）。见上面的引理
\ref{schemes-lemma-fibre-product-immersion}。
\end{definition}

\noindent
有时我们也把这个术语用于局部闭子概形和开子概形。









\section{代数几何中的基变换}
\label{schemes-section-base-change}

\noindent
引入概形语言的一个动机是：它精确刻画了在某个给定域上定义一个簇
是什么意思。例如，簇 $X$ 在 $\mathbf{Q}$ 上，等价于
（《簇》，定义 \ref{varieties-definition-variety}）一个有限型、分离、
不可约且约化的态射 $X \to \Spec(\mathbf{Q})$\footnote{当然，代数几何学家
仍会争论是否应当要求 $X$ 在 $\mathbf{Q}$ 上几何不可约。}。
无论如何，更一般的思想是在给定的{\it 基底概形}上研究概形，
这个基底通常记为 $S$。我们说“令 $X$ 为 $S$ 上的概形”，
仅仅是指 $X$ 配备了一个态射 $X \to S$。在图中，我们会尽量把
{\it 结构态射} $X \to S$ 画成从 $X$ 向下指向 $S$ 的箭头。
我们往往更关心 $X$ 相对于 $S$ 的性质，而不是 $X$ 的内在几何。
例如，我们希望了解态射 $X \to S$ 的纤维、$X$ 在基变换后会发生什么，等等。

\medskip\noindent
下面引入一些惯用术语。当然，这些术语只是研究一个范畴中某个给定对象
之上对象的范畴这一思想的特例；见《范畴》，
例 \ref{categories-example-category-over-X}。

\begin{definition}
\label{schemes-definition-base-change}
令 $S$ 为概形。
\begin{enumerate}
\item 我们称 $X$ 为{\it $S$ 上的概形}，是指 $X$ 配备一个概形态射
$X \to S$。态射 $X \to S$ 有时称为{\it 结构态射}。
\item 若 $R$ 为环，我们说 $X$ 是{\it $R$ 上的概形}，
而不说 $X$ 是 $\Spec(R)$ 上的概形。
\item 所谓{\it 态射 $f : X \to Y$，即 $S$ 上概形的态射}，
是指一个概形态射，使得复合 $X \to Y \to S$
（由 $f$ 与 $Y$ 的结构态射组成）等于 $X$ 的结构态射。
\item 记 $\Mor_S(X, Y)$ 为所有从 $X$ 到 $Y$ 且位于 $S$ 上的态射所成的集合。
\item 令 $X$ 为 $S$ 上的概形，并令 $S' \to S$ 为概形态射。
$X$ 的{\it 基变换}是概形 $X_{S'} = S' \times_S X$，视为 $S'$ 上的概形。
\item 令 $f : X \to Y$ 为 $S$ 上的概形态射，并令 $S' \to S$ 为概形态射。
$f$ 的{\it 基变换}是诱导态射 $f' : X_{S'} \to Y_{S'}$
（即态射 $\text{id}_{S'} \times_{\text{id}_S} f$）。
\item 令 $R$ 为环。令 $X$ 为 $R$ 上的概形。令 $R \to R'$ 为环同态。
{\it 基变换} $X_{R'}$ 是概形 $\Spec(R') \times_{\Spec(R)} X$，
视为 $R'$ 上的概形。
\end{enumerate}
\end{definition}

\noindent
下面是一个典型结论。

\begin{lemma}
\label{schemes-lemma-base-change-immersion}
令 $S$ 为概形。令 $f : X \to Y$ 为 $S$ 上概形的浸入
（分别地，闭浸入、开浸入）。则 $f$ 的任意基变换仍为浸入
（分别地，闭浸入、开浸入）。
\end{lemma}

\begin{proof}
可以把 $f$ 沿态射 $S' \to S$ 的基变换看作下列交换图中
左上方的竖直箭头：
$$
\xymatrix{
X_{S'} \ar[r] \ar[d] & X \ar[d] \ar@/^4ex/[dd] \\
Y_{S'} \ar[r] \ar[d] & Y \ar[d] \\
S' \ar[r] & S
}
$$
此图蕴含 $X_{S'} \cong Y_{S'} \times_Y X$，
本引理由引理 \ref{schemes-lemma-fibre-product-immersion} 得出。
\end{proof}

\noindent
事实上，这类结论如此典型，以至于有一个专门术语来表述它。

\begin{definition}
\label{schemes-definition-preserved-by-base-change}
性质与基变换。
\begin{enumerate}
\item 令 $\mathcal{P}$ 为基底上概形的一种性质。我们称
$\mathcal{P}$ 在{\it 任意基变换下保持}，或简称
$\mathcal{P}$ 在{\it 基变换下保持}，如果每当 $X/S$ 具有
$\mathcal{P}$ 时，任意基变换 $X_{S'}/S'$ 都具有 $\mathcal{P}$。
\item 令 $\mathcal{P}$ 为基底上概形态射的一种性质。我们称
$\mathcal{P}$ 在{\it 任意基变换下保持}，或简称在{\it 基变换下保持}，
如果每当 $f : X \to Y$ 在 $S$ 上具有 $\mathcal{P}$ 时，任意基变换
$f' : X_{S'} \to Y_{S'}$ 在 $S'$ 上都具有 $\mathcal{P}$。
\end{enumerate}
\end{definition}

\noindent
现在可以说，“是闭浸入”这一性质在任意基变换下保持。

\begin{definition}
\label{schemes-definition-fibre}
令 $f : X \to S$ 为概形态射。令 $s \in S$ 为一点。
所谓{\it 概形论纤维 $X_s$，即 $f$ 在 $s$ 上方的纤维}，
或简称{\it $f$ 在 $s$ 上方的纤维}，是处于下列纤维积图中的概形：
$$
\xymatrix{
X_s = \Spec(\kappa(s)) \times_S X \ar[r] \ar[d] &
X \ar[d] \\
\Spec(\kappa(s)) \ar[r] &
S
}
$$
我们始终把纤维 $X_s$ 看作 $\kappa(s)$ 上的概形。
\end{definition}

\begin{lemma}
\label{schemes-lemma-fibre-topological}
令 $f : X \to S$ 为概形态射。考虑下列两图：
$$
\vcenter{
\xymatrix{
X_s \ar[r] \ar[d] & X \ar[d] \\
\Spec(\kappa(s)) \ar[r] & S
}
}
\quad\text{且}\quad
\vcenter{
\xymatrix{
\Spec(\mathcal{O}_{S, s}) \times_S X \ar[r] \ar[d] & X \ar[d] \\
\Spec(\mathcal{O}_{S, s}) \ar[r] & S
}
}
$$
在两种情形中，这些图都诱导拓扑空间的纤维方块；特别地，
上方的水平箭头是到其像的同胚。
\end{lemma}

\begin{proof}
选取仿射开集 $U \subset S$，使其包含 $s$。
下方的水平态射经由 $U$ 分解；例如见引理
\ref{schemes-lemma-morphism-from-spec-local-ring}。
因此可假设 $S$ 是仿射的。若 $X$ 也是仿射的，则结论由
《代数》，注 \ref{algebra-remark-fundamental-diagram} 得出。
一般情形可通过用仿射开集覆盖 $X$ 得出。
\end{proof}

\begin{lemma}
\label{schemes-lemma-local-ring-fibre}
令 $f : X \to S$ 为概形态射。令 $x \in X$ 的像为 $s \in S$。
把 $x$ 看作 $X_s$ 的一点（见引理 \ref{schemes-lemma-fibre-topological}），则有
$$
\mathcal{O}_{X_s, x} \cong
\mathcal{O}_{X, x}/\mathfrak m_s\mathcal{O}_{X, x} \cong
\mathcal{O}_{X, x} \otimes_{\mathcal{O}_{S, s}} \kappa(s)
$$
\end{lemma}

\begin{proof}
与引理 \ref{schemes-lemma-fibre-topological} 的证明相同，这归结为 $X$ 与 $S$
均仿射的情形。在仿射情形中，该陈述转化为《代数》，
注 \ref{algebra-remark-local-ring-fibre}。
\end{proof}







\section{拟紧态射}
\label{schemes-section-quasi-compact}

\noindent
若一个概形的底拓扑空间拟紧，则称该概形{\it 拟紧}。
相对版本的概念定义如下。

\begin{definition}
\label{schemes-definition-quasi-compact}
若一个概形态射的底拓扑空间映射拟紧，则称该态射{\it 拟紧}；
见《拓扑》，定义 \ref{topology-definition-quasi-compact}。
\end{definition}

\begin{lemma}
\label{schemes-lemma-quasi-compact-affine}
令 $f : X \to S$ 为概形态射。下列条件等价：
\begin{enumerate}
\item $f : X \to S$ 拟紧；
\item 每个仿射开集的逆像拟紧；
\item 存在某个仿射开覆盖 $S = \bigcup_{i \in I} U_i$，
使得 $f^{-1}(U_i)$ 对所有 $i$ 都拟紧。
\end{enumerate}
\end{lemma}

\begin{proof}
设给定 (3) 中的覆盖 $S = \bigcup_{i \in I} U_i$。
首先，令 $U \subset S$ 为任意仿射开集。对任意 $u \in U$，
可以找到指标 $i(u) \in I$，使得 $u \in U_{i(u)}$。
由于标准开集构成 $U_{i(u)}$ 的拓扑基，可以找到
$W_u \subset U \cap U_{i(u)}$，它是 $U_{i(u)}$ 中的标准开集。
由拟紧性，可以找到有限多个点 $u_1, \ldots, u_n \in U$，使得
$U = \bigcup_{j = 1}^n W_{u_j}$。对每个 $j$，把
$f^{-1}U_{i(u_j)} = \bigcup_{k \in K_j} V_{jk}$ 写成有限多个仿射开集的并。
由于 $W_{u_j} \subset U_{i(u_j)}$ 是标准开集，可知
$f^{-1}(W_{u_j}) \cap V_{jk}$ 是 $V_{jk}$ 的标准开集；见《代数》，
引理 \ref{algebra-lemma-spec-functorial}。因此
$f^{-1}(W_{u_j}) \cap V_{jk}$ 是仿射的，从而
$f^{-1}(W_{u_j})$ 是有限多个仿射概形的并。这证明任意仿射开集的
逆像都是有限多个仿射开集的并。

\medskip\noindent
接着，假设每个仿射开集的逆像都是有限多个仿射开集的并。
令 $K \subset S$ 为任意拟紧开集。由于 $S$ 具有由仿射开集组成的拓扑基，
可知 $K$ 是有限多个仿射开集的并。因此 $K$ 的逆像是有限多个仿射开集的并，
所以 $f$ 拟紧。

\medskip\noindent
最后，假设 $f$ 拟紧。此时，前一段的论证说明任意仿射开集的逆像
都是有限多个仿射开集的并。
\end{proof}

\begin{lemma}
\label{schemes-lemma-quasi-compact-preserved-base-change}
“拟紧”是基底上概形态射的一种性质，并且在任意基变换下保持。
\end{lemma}

\begin{proof}
略。
\end{proof}

\begin{lemma}
\label{schemes-lemma-composition-quasi-compact}
拟紧态射的复合拟紧。
\end{lemma}

\begin{proof}
这由定义及《拓扑》，
引理 \ref{topology-lemma-composition-quasi-compact} 得出。
\end{proof}

\begin{lemma}
\label{schemes-lemma-closed-immersion-quasi-compact}
闭浸入拟紧。
\end{lemma}

\begin{proof}
这由定义及《拓扑》，
引理 \ref{topology-lemma-closed-in-quasi-compact} 得出。
\end{proof}

\begin{example}
\label{schemes-example-open-immersion-not-quasi-compact}
一般而言，开浸入不拟紧。标准例子是开子空间
$U \subset X$，其中 $X = \Spec(k[x_1, x_2, x_3, \ldots])$，
$U$ 是 $X \setminus \{0\}$，而 $0$ 是 $X$ 中对应于极大理想
$(x_1, x_2, x_3, \ldots)$ 的点。
\end{example}

\begin{lemma}
\label{schemes-lemma-image-quasi-compact-closed}
令 $f : X \to S$ 为拟紧概形态射。下列条件等价：
\begin{enumerate}
\item $f(X) \subset S$ 是闭集；
\item $f(X) \subset S$ 在特化下稳定。
\end{enumerate}
\end{lemma}

\begin{proof}
由《拓扑》，引理 \ref{topology-lemma-open-closed-specialization}，
有 (1) $\Rightarrow$ (2)。假设 (2)。令 $U \subset S$ 为仿射开集。
只须证明 $f(X) \cap U$ 闭。由于 $U \cap f(X)$ 在 $U$ 中的特化下稳定，
问题归结为 $S$ 仿射的情形。因为 $f$ 拟紧，可推出
$X = f^{-1}(S)$ 拟紧，这是由于 $S$ 仿射。因此可以写成
$X = \bigcup_{i = 1}^n U_i$，其中 $U_i \subset X$ 是仿射开集。
设 $S = \Spec(R)$，并令 $U_i = \Spec(A_i)$，其中
$R$-代数为 $A_i$。于是
$f(X) = \Im(\Spec(A_1 \times \ldots \times A_n)
\to \Spec(R))$。因此本引理由《代数》，
引理 \ref{algebra-lemma-image-stable-specialization-closed} 得出。
\end{proof}

\begin{lemma}
\label{schemes-lemma-quasi-compact-closed}
令 $f : X \to S$ 为拟紧概形态射。则 $f$ 是闭映射，当且仅当特化可沿
$f$ 提升；见《拓扑》，定义 \ref{topology-definition-lift-specializations}。
\end{lemma}

\begin{proof}
由《拓扑》，引理 \ref{topology-lemma-closed-open-map-specialization}，
若 $f$ 是闭映射，则特化可沿 $f$ 提升。反过来，假设特化可沿 $f$ 提升。
令 $Z \subset X$ 为闭子集。赋予 $Z$ 约化诱导概形结构，便可把它看作概形；
见定义 \ref{schemes-definition-reduced-induced-scheme}。由于 $Z \subset X$ 闭，
$f$ 在 $Z$ 上的限制仍拟紧。此外，特化也可沿 $Z \to S$ 提升；
见《拓扑》，引理 \ref{topology-lemma-lift-specialization-composition}。
因此，只须证明 $f(X)$ 闭，前提是特化可沿 $f$ 提升。
特别地，$f(X)$ 在特化下稳定；见《拓扑》，
引理 \ref{topology-lemma-lift-specializations-images}。于是由引理
\ref{schemes-lemma-image-quasi-compact-closed}，$f(X)$ 闭。
\end{proof}










\section{万有闭性的赋值判据}
\label{schemes-section-valuative-criterion-universal-closedness}

\noindent
《拓扑》第 \ref{topology-section-proper} 节讨论了固有映射：
可以把它看作所有纤维都拟紧的拓扑空间闭映射，
也可以看作所有基变换都是闭映射的映射。下面给出代数几何中的对应概念。

\begin{definition}
\label{schemes-definition-universally-closed}
若概形态射 $f : X \to S$ 的每个基变换
$f' : X_{S'} \to S'$ 都是闭映射，则称它{\it 万有闭}。
\end{definition}

\noindent
事实上，形容词“万有”经常这样使用。换言之，给定态射的一种性质
$\mathcal{P}$，我们说“$X \to S$ 万有地具有 $\mathcal{P}$”，
当且仅当每个基变换 $X_{S'} \to S'$ 都具有 $\mathcal{P}$。

\medskip\noindent
关于万有闭态射性质的更详细讨论，请见《态射》，
第 \ref{morphisms-section-proper} 节。本节只讨论万有闭态射与满足
赋值判据存在性部分的态射之间的关系。

\begin{lemma}
\label{schemes-lemma-specializations-lift}
令 $f : X \to S$ 为概形态射。
\begin{enumerate}
\item 若 $f$ 万有闭，则特化可沿 $f$ 的任意基变换提升；
见《拓扑》，定义 \ref{topology-definition-lift-specializations}。
\item 若 $f$ 拟紧，且特化可沿 $f$ 的任意基变换提升，则 $f$ 万有闭。
\end{enumerate}
\end{lemma}

\begin{proof}
(1) 直接由《拓扑》，引理
\ref{topology-lemma-closed-open-map-specialization} 得出。
(2) 由引理 \ref{schemes-lemma-quasi-compact-closed} 与
\ref{schemes-lemma-quasi-compact-preserved-base-change} 得出。
\end{proof}

\begin{definition}
\label{schemes-definition-valuative-criterion}
令 $f : X \to S$ 为概形态射。称 $f${\it 满足赋值判据的存在性部分}，
如果给定任意实线交换图
$$
\xymatrix{
\Spec(K) \ar[r] \ar[d] & X \ar[d] \\
\Spec(A) \ar[r] \ar@{-->}[ru] & S
}
$$
其中 $A$ 是以 $K$ 为分式域的赋值环，虚线箭头总存在。称 $f$
{\it 满足赋值判据的唯一性部分}，如果对任意上述图，虚线箭头
至多有一个（当然不要求存在）。
\end{definition}

\noindent
{\it 赋值环}是其分式域中支配关系下极大的局部整环；见《代数》，
定义 \ref{algebra-definition-valuation-ring}。因此，赋值环的谱有唯一泛点
$\eta$ 和唯一闭点 $0$，当然还有特化 $\eta \leadsto 0$。
赋值环的重要性在于：任意概形中点的任何特化，都是某个赋值环之谱的
某个态射下 $\eta \leadsto 0$ 的像。精确结论如下。

\begin{lemma}
\label{schemes-lemma-points-specialize}
\begin{slogan}
特化由赋值环见证。
\end{slogan}
令 $S$ 为概形。令 $s' \leadsto s$ 为 $S$ 中点的特化。则：
\begin{enumerate}
\item 存在赋值环 $A$ 和态射 $f : \Spec(A) \to S$，使得泛点
$\eta$（属于 $\Spec(A)$）映到 $s'$，而特殊点映到 $s$；
\item 给定域扩张 $K/\kappa(s')$，可以安排使扩张
$\kappa(\eta)/\kappa(s')$（由 $f$ 诱导）同构于给定扩张。
\end{enumerate}
\end{lemma}

\begin{proof}
令 $s' \leadsto s$ 为 $S$ 中的特化，并令 $K/\kappa(s')$ 为域扩张。
由引理 \ref{schemes-lemma-specialize-points} 及引理
\ref{schemes-lemma-characterize-points} 后面的讨论，得到环映射
$\mathcal{O}_{S, s} \to \kappa(s') \to K$。
令 $A \subset K$ 为任意赋值环，其分式域为 $K$，并且它支配
$\mathcal{O}_{S, s} \to K$ 的像；见《代数》，引理
\ref{algebra-lemma-dominate}。环同态 $\mathcal{O}_{S, s} \to A$
诱导态射 $f : \Spec(A) \to S$；见引理
\ref{schemes-lemma-morphism-from-spec-local-ring}。按照构造，此态射具有所有所需性质。
\end{proof}

\begin{lemma}
\label{schemes-lemma-lift-specializations-valuative}
令 $f : X \to S$ 为概形态射。下列条件等价：
\begin{enumerate}
\item 特化可沿 $f$ 的任意基变换提升；
\item 态射 $f$ 满足赋值判据的存在性部分。
\end{enumerate}
\end{lemma}

\begin{proof}
假设 (1) 成立。给定定义 \ref{schemes-definition-valuative-criterion} 中的实线图。
为了找到虚线箭头，可以把 $X \to S$ 替换为
$X_{\Spec(A)} \to \Spec(A)$，因为该假设在基变换下保持。
因此可假设 $S = \Spec(A)$。令 $x' \in X$ 为
$\Spec(K) \to X$ 的像，于是 $\kappa(x') \subset K$；
见引理 \ref{schemes-lemma-characterize-points}。
由假设，存在特化 $x' \leadsto x$，它发生在 $X$ 中，使得 $x$ 映到
$S = \Spec(A)$ 的闭点。得到局部环同态
$A \to \mathcal{O}_{X, x}$ 以及环同态
$\mathcal{O}_{X, x} \to \kappa(x')$；见引理
\ref{schemes-lemma-specialize-points} 及引理 \ref{schemes-lemma-characterize-points}
后面的讨论。复合
$A \to \mathcal{O}_{X, x} \to \kappa(x') \to K$ 是给定的单射
$A \to K$。由于 $A \to \mathcal{O}_{X, x}$ 是局部同态，
$\mathcal{O}_{X, x} \to K$ 的像支配 $A$，因而等于 $A$；
见《代数》，定义 \ref{algebra-definition-valuation-ring}。
因此得到环同态 $\mathcal{O}_{X, x} \to A$，进而得到态射
$\Spec(A) \to X$（见引理 \ref{schemes-lemma-morphism-from-spec-local-ring}
及其后的讨论）。这证明了 (2)。

\medskip\noindent
反过来，假设 (2) 成立。考察下列交换图，立即可知赋值判据的存在性部分
对任意基变换 $X_{S'} \to S'$ 都成立；这里该基变换来自 $f$：
$$
\xymatrix{
\Spec(K) \ar[r] \ar[d] & X_{S'} \ar[r] \ar[d] & X \ar[d] \\
\Spec(A) \ar[r] \ar@{-->}[ru] \ar@{-->}[rru] & S' \ar[r] & S
}
$$
具体而言，由纤维积的定义，较为水平的虚线箭头会给出另一个虚线箭头。
因此显然只须证明特化可沿 $f$ 提升。
令 $s' \leadsto s$ 为 $S$ 中的特化，并令 $x' \in X$ 为位于 $s'$ 上方的点。
将引理 \ref{schemes-lemma-points-specialize} 应用于 $s' \leadsto s$ 以及域扩张
$K = \kappa(x')/\kappa(s')$，得到交换图
$$
\xymatrix{
\Spec(K) \ar[rr] \ar[d] & & X \ar[d] \\
\Spec(A) \ar[r] \ar@{-->}[rru] &
\Spec(\mathcal{O}_{S, s}) \ar[r] & S
}
$$
由条件 (2) 得到虚线箭头。闭点的像 $x$，这里闭点属于 $\Spec(A)$，
而该像位于 $X$ 中，便是所求解。也就是说，$x$ 是 $x'$ 的特化，
并映到 $s$。
\end{proof}

\begin{proposition}[万有闭性的赋值判据]
\label{schemes-proposition-characterize-universally-closed}
令 $f$ 为拟紧概形态射。则 $f$ 万有闭，当且仅当 $f$
满足赋值判据的存在性部分。
\end{proposition}

\begin{proof}
这是上面的引理 \ref{schemes-lemma-specializations-lift} 与
\ref{schemes-lemma-lift-specializations-valuative} 的形式推论。
\end{proof}

\begin{example}
\label{schemes-example-projective-line-universally-closed}
令 $k$ 为域。考虑结构态射
$p : \mathbf{P}^1_k \to \Spec(k)$，它来自 $k$ 上的射影直线；
见例 \ref{schemes-example-projective-line}。
利用上面的赋值判据证明 $p$ 万有闭。
按照构造，$\mathbf{P}^1_k$ 由两个仿射开集覆盖，因而 $p$ 拟紧。
给定交换图
$$
\xymatrix{
\Spec(K) \ar[r]_\xi \ar[d] & \mathbf{P}^1_k \ar[d] \\
\Spec(A) \ar[r]^\varphi & \Spec(k)
}
$$
其中 $A$ 是赋值环，$K$ 是其分式域。回忆 $\mathbf{P}^1_k$ 由
$\Spec(k[x])$ 与 $\Spec(k[y])$ 粘合而得，其中把 $D(x)$ 与 $D(y)$
经由 $x = y^{-1}$ 粘合（或更对称地写成 $xy = 1$）。
为了证明存在态射 $\Spec(A) \to \mathbf{P}^1_k$ 斜向嵌入上图，
由对称性可假设 $\xi$ 映入开集 $\Spec(k[x])$。这给出下列环的交换图：
$$
\xymatrix{
K & k[x] \ar[l]^{\xi^\sharp} \\
A \ar[u] & k \ar[u] \ar[l]_{\varphi^\sharp}
}
$$
由《代数》，引理 \ref{algebra-lemma-valuation-ring-x-or-x-inverse}，
可知或者 $\xi^\sharp(x) \in A$，或者 $\xi^\sharp(x)^{-1} \in A$。
在第一种情形，得到环同态
$$
k[x] \to A,
\ \lambda \mapsto \varphi^\sharp(\lambda),
\  x \mapsto \xi^\sharp(x)
$$
它嵌入上面的环图中，结论得证。在第二种情形，得到环同态
$$
k[y] \to A,
\ \lambda \mapsto \varphi^\sharp(\lambda),
\ y \mapsto \xi^\sharp(x)^{-1}.
$$
这给出态射
$\Spec(A) \to \Spec(k[y]) \to \mathbf{P}^1_k$，
它斜向嵌入本例开头的交换图中（验证略）。
\end{example}





























\section{分离公理}
\label{schemes-section-separation-axioms}

\noindent
拓扑空间 $X$ 是豪斯多夫的，当且仅当对角线
$\Delta \subset X \times X$ 是闭子集。代数几何中的类似做法是：
给定概形 $X$，它位于基底概形 $S$ 上，考虑对角态射
$$
\Delta_{X/S} : X \longrightarrow X \times_S X.
$$
这是满足
$\text{pr}_1 \circ \Delta_{X/S} = \text{id}_X$ 以及
$\text{pr}_2 \circ \Delta_{X/S} = \text{id}_X$ 的唯一概形态射
（它在任何具有纤维积的范畴中都存在）。

\begin{lemma}
\label{schemes-lemma-diagonal-affines-closed}
仿射概形之间态射的对角态射是闭浸入。
\end{lemma}

\begin{proof}
与态射 $\Spec(S) \to \Spec(R)$ 相关联的对角态射，是由环同态
$S \otimes_R S \to S$、$a \otimes b \mapsto ab$ 所对应的谱上态射。
此映射显然满射，故有 $S \cong S \otimes_R S/J$，其中
$J \subset S \otimes_R S$ 是某个理想。因此 $\Delta$ 是闭浸入；
见例 \ref{schemes-example-closed-immersion-affines}。
\end{proof}

\begin{lemma}
\label{schemes-lemma-diagonal-immersion}
\begin{slogan}
相对概形的对角态射是浸入。
\end{slogan}
令 $X$ 为 $S$ 上的概形。对角态射 $\Delta_{X/S}$ 是浸入。
\end{lemma}

\begin{proof}
回忆，若 $V \subset X$ 是仿射开集并映入仿射开集 $U \subset S$，
则 $V \times_U V$ 是 $X \times_S X$ 中的仿射开集；见引理
\ref{schemes-lemma-fibre-product-affines} 与 \ref{schemes-lemma-open-fibre-product}。
考虑开子概形 $W$（属于 $X \times_S X$），它是这些仿射开集
$V \times_U V$ 的并。由引理 \ref{schemes-lemma-closed-local-target}，
只须证明每个态射
$\Delta_{X/S}^{-1}(V \times_U V) \to V \times_U V$ 都是闭浸入。
由于 $V = \Delta_{X/S}^{-1}(V \times_U V)$，这正是在检验
$\Delta_{V/U}$ 是闭浸入，而这就是引理
\ref{schemes-lemma-diagonal-affines-closed}。
\end{proof}

\begin{definition}
\label{schemes-definition-separated}
令 $f : X \to S$ 为概形态射。
\begin{enumerate}
\item 称 $f${\it 分离}，如果对角态射 $\Delta_{X/S}$ 是闭浸入。
\item 称 $f${\it 拟分离}，如果对角态射 $\Delta_{X/S}$ 拟紧。
\item 称概形 $Y${\it 分离}，如果态射 $Y \to \Spec(\mathbf{Z})$ 分离。
\item 称概形 $Y${\it 拟分离}，如果态射 $Y \to \Spec(\mathbf{Z})$ 拟分离。
\end{enumerate}
\end{definition}

\noindent
由引理 \ref{schemes-lemma-diagonal-immersion} 与 \ref{schemes-lemma-immersion-when-closed}，
可知 $\Delta_{X/S}$ 是闭浸入，当且仅当
$\Delta_{X/S}(X) \subset X \times_S X$ 是闭子集。此外，由引理
\ref{schemes-lemma-closed-immersion-quasi-compact}，分离态射必拟分离。
引入拟分离态射的原因是：研究代数簇时（尤其在模空间、代数栈等问题中）
自然会出现非分离态射，但它们多数仍然拟分离。

\begin{example}
\label{schemes-example-not-quasi-separated}
下面给出一个非拟分离态射的例子。设
$X = X_1 \cup X_2 \to S = \Spec(k)$，其中
$X_1 = X_2 = \Spec(k[t_1, t_2, t_3, \ldots])$，
并沿 $\{0\} = \{(t_1, t_2, t_3, \ldots)\}$ 的补集粘合
（粘合法同例 \ref{schemes-example-affine-space-zero-doubled}）。
此时，仿射概形 $X_1 \times_S X_2$ 在 $\Delta_{X/S}$ 下的逆像是概形
$\Spec(k[t_1, t_2, t_3, \ldots]) \setminus \{0\}$，它不拟紧。
\end{example}

\begin{lemma}
\label{schemes-lemma-where-are-they-equal}
令 $X$、$Y$ 为 $S$ 上的概形。令 $a, b : X \to Y$ 为 $S$ 上的概形态射。
存在最大的局部闭子概形 $Z \subset X$，使得 $a|_Z = b|_Z$。
事实上，$Z$ 是 $(a, b)$ 的等化子。此外，若 $Y$ 在 $S$ 上分离，
则 $Z$ 是闭子概形。
\end{lemma}

\begin{proof}
由范畴论理由，$(a, b)$ 的等化子是下图中的纤维积 $Z$：
$$
\xymatrix{
Z = Y \times_{(Y \times_S Y)} X \ar[r] \ar[d] &
 X \ar[d]^{(a , b)} \\
Y \ar[r]^-{\Delta_{Y/S}} & Y \times_S Y
}
$$
因此本引理由引理 \ref{schemes-lemma-base-change-immersion}、
\ref{schemes-lemma-diagonal-immersion} 及定义 \ref{schemes-definition-separated} 得出。
\end{proof}

\begin{lemma}
\label{schemes-lemma-characterize-quasi-separated}
令 $f : X \to S$ 为概形态射。下列条件等价：
\begin{enumerate}
\item 态射 $f$ 拟分离。
\item 对每一对仿射开集 $U, V \subset X$，若它们映入 $S$ 的同一个
仿射开集，则交集 $U \cap V$ 是 $X$ 的有限多个仿射开集的并。
\item 存在仿射开覆盖 $S = \bigcup_{i \in I} U_i$，并且对每个 $i$
存在仿射开覆盖 $f^{-1}U_i = \bigcup_{j \in I_i} V_j$，使得对每个 $i$
以及每一对 $j, j' \in I_i$，交集 $V_j \cap V_{j'}$ 都是 $X$ 的
有限多个仿射开集的并。
\end{enumerate}
\end{lemma}

\begin{proof}
先证明 (3) 蕴含 (1)。由引理 \ref{schemes-lemma-affine-covering-fibre-product}，覆盖
$X \times_S X = \bigcup_i \bigcup_{j, j'} V_j \times_{U_i} V_{j'}$
是 $X \times_S X$ 的仿射开覆盖。此外，
$\Delta_{X/S}^{-1}(V_j \times_{U_i} V_{j'}) = V_j \cap V_{j'}$。
故该蕴含由引理 \ref{schemes-lemma-quasi-compact-affine} 得出。

\medskip\noindent
(1) $\Rightarrow$ (2) 来自如下事实：在 (2) 的假设下，纤维积
$U \times_S V$ 是 $X \times_S X$ 的仿射开集。
(2) $\Rightarrow$ (3) 是平凡的。
\end{proof}

\begin{lemma}
\label{schemes-lemma-characterize-separated}
令 $f : X \to S$ 为概形态射。
\begin{enumerate}
\item 若 $f$ 分离，则对每一对仿射开集 $(U, V)$（属于 $X$），
只要它们映入 $S$ 的同一个仿射开集，就有：
\begin{enumerate}
\item 交集 $U \cap V$ 是仿射的；
\item 环同态
$\mathcal{O}_X(U) \otimes_{\mathbf{Z}} \mathcal{O}_X(V)
\to \mathcal{O}_X(U \cap V)$
满射。
\end{enumerate}
\item 若任意一对点 $x_1, x_2 \in X$ 位于同一点 $s \in S$ 上方，
都分别包含在仿射开集 $x_1 \in U$、$x_2 \in V$ 中，且这些开集映入
$S$ 的同一个仿射开集并满足 (a)、(b)，则 $f$ 分离。
\end{enumerate}
\end{lemma}

\begin{proof}
假设 $f$ 分离。设 $(U, V)$ 是 (1) 中的一对开集。
令 $W = \Spec(R)$ 为 $S$ 的仿射开集，并同时包含 $f(U)$ 与 $f(V)$。
写成 $U = \Spec(A)$、$V = \Spec(B)$，其中用到 $R$-代数 $A$ 与 $B$。
由引理 \ref{schemes-lemma-open-fibre-product}，
$U \times_S V = U \times_W V = \Spec(A \otimes_R B)$
是 $X \times_S X$ 的仿射开集。因而由引理
\ref{schemes-lemma-closed-subspace-scheme} 可知，
$\Delta^{-1}(U \times_S V) \to U \times_S V$ 可等同于
$\Spec((A \otimes_R B)/J) \to \Spec(A \otimes_R B)$，
其中 $J \subset A \otimes_R B$ 是某个理想。
因此 $U \cap V = \Delta^{-1}(U \times_S V)$ 是仿射的。
断言 (1)(b) 成立，因为
$A \otimes_{\mathbf{Z}} B \to (A \otimes_R B)/J$ 满射。

\medskip\noindent
假设 (2) 中表述的假设成立。显然，仿射开集 $U \times_S V$，
对应于 (2) 中各对 $(U, V)$，构成 $X \times_S X$ 的仿射开覆盖
（例如见引理 \ref{schemes-lemma-affine-covering-fibre-product}）。
因此，只须证明每个态射
$U \cap V = \Delta_{X/S}^{-1}(U \times_S V) \to U \times_S V$
都是闭浸入；见引理 \ref{schemes-lemma-closed-local-target}。
由假设 (a)，有 $U \cap V = \Spec(C)$，其中 $C$ 为某个环。
选取仿射开集 $W = \Spec(R)$（属于 $S$），使 $U$ 与 $V$ 都映入其中，
并写成 $U = \Spec(A)$、$V = \Spec(B)$，则假设 (b) 是说复合
$$
A \otimes_{\mathbf{Z}} B \to A \otimes_R B \to C
$$
满射。因此 $A \otimes_R B \to C$ 满射，从而
$\Spec(C) \to \Spec(A \otimes_R B)$ 是闭浸入。
\end{proof}

\begin{example}
\label{schemes-example-projective-line-separated}
令 $k$ 为域。考虑结构态射
$p : \mathbf{P}^1_k \to \Spec(k)$，它来自 $k$ 上的射影直线；
见例 \ref{schemes-example-projective-line}。利用上面的引理证明 $p$ 分离。
按照构造，$\mathbf{P}^1_k$ 由两个仿射开集
$U = \Spec(k[x])$ 与 $V = \Spec(k[y])$ 覆盖，其交集为
$U \cap V = \Spec(k[x, y]/(xy - 1))$（采用显然的记号）。
因此，只须检验引理 \ref{schemes-lemma-characterize-separated} 的条件 (2)(a)、(2)(b)
对仿射开集对 $(U, U)$、$(U, V)$、$(V, U)$ 及 $(V, V)$ 成立。
对 $(U, U)$ 与 $(V, V)$ 而言这是平凡的。对 $(U, V)$ 而言，
这相当于证明 $U \cap V$ 仿射（确实如此），并且环同态
$$
k[x] \otimes_{\mathbf{Z}} k[y] \longrightarrow k[x, y]/(xy - 1)
$$
满射。这是显然的，因为右端的任意元素都可写成关于 $x$ 的多项式
与关于 $y$ 的多项式之和。
\end{example}

\begin{lemma}
\label{schemes-lemma-fibre-product-after-map}
令 $f : X \to T$ 与 $g : Y \to T$ 为具有相同目标的概形态射，
并令 $h : T \to S$ 为概形态射。则诱导态射
$i : X \times_T Y \to X \times_S Y$ 是浸入。若 $T \to S$ 分离，
则 $i$ 是闭浸入；若 $T \to S$ 拟分离，则 $i$ 拟紧。
\end{lemma}

\begin{proof}
由一般范畴论，下图
$$
\xymatrix{
X \times_T Y \ar[r] \ar[d] & X \times_S Y \ar[d] \\
T \ar[r]^{\Delta_{T/S}} \ar[r] & T \times_S T
}
$$
是纤维积图。本引理由引理 \ref{schemes-lemma-diagonal-immersion}、
\ref{schemes-lemma-fibre-product-immersion} 与
\ref{schemes-lemma-quasi-compact-preserved-base-change} 得出。
\end{proof}

\begin{lemma}
\label{schemes-lemma-semi-diagonal}
令 $g : X \to Y$ 为 $S$ 上的概形态射。态射
$i : X \to X \times_S Y$ 是浸入。若 $Y$ 在 $S$ 上分离，
则它是闭浸入；若 $Y$ 在 $S$ 上拟分离，则它拟紧。
\end{lemma}

\begin{proof}
这是把引理 \ref{schemes-lemma-fibre-product-after-map} 应用于态射
$X = X \times_Y Y \to X \times_S Y$ 所得的特殊情形。
\end{proof}

\begin{lemma}
\label{schemes-lemma-section-immersion}
令 $f : X \to S$ 为概形态射。令 $s : S \to X$ 为 $f$ 的截面
（用公式写为 $f \circ s = \text{id}_S$）。则 $s$ 是浸入。
若 $f$ 分离，则 $s$ 是闭浸入；若 $f$ 拟分离，则 $s$ 拟紧。
\end{lemma}

\begin{proof}
这是把引理 \ref{schemes-lemma-semi-diagonal} 应用于 $g =s$ 所得的特殊情形；
此时态射为 $i = s : S \to S \times_S X$。
\end{proof}

\begin{lemma}
\label{schemes-lemma-separated-permanence}
恒久性质。
\begin{enumerate}
\item 分离态射的复合分离。
\item 拟分离态射的复合拟分离。
\item 分离态射的基变换分离。
\item 拟分离态射的基变换拟分离。
\item 分离态射的（纤维）积分离。
\item 拟分离态射的（纤维）积拟分离。
\end{enumerate}
\end{lemma}

\begin{proof}
令 $X \to Y \to Z$ 为态射。假设 $X \to Y$ 与 $Y \to Z$ 分离。复合
$$
X \to X \times_Y X \to X \times_Z X
$$
是闭浸入，因为第一个态射依假设是闭浸入，第二个态射则由引理
\ref{schemes-lemma-fibre-product-after-map} 是闭浸入。同一论证也适用于“拟分离”
（使用相同的引理）。

\medskip\noindent
令 $f : X \to Y$ 为基底 $S$ 上的概形态射。令 $S' \to S$ 为概形态射。
令 $f' : X_{S'} \to Y_{S'}$ 为 $f$ 的基变换。则 $f'$ 的对角态射为
$$
\Delta_{f'} :
X_{S'} = S' \times_S X
\longrightarrow
X_{S'} \times_{Y_{S'}} X_{S'} = S' \times _S (X \times_Y X)
$$
容易看出，这是 $\Delta_f$ 的基变换。因此，(3)、(4) 由如下事实得出：
闭浸入与拟紧态射在任意基变换下保持（引理
\ref{schemes-lemma-fibre-product-immersion} 与
\ref{schemes-lemma-quasi-compact-preserved-base-change}）。

\medskip\noindent
若 $f : X \to Y$ 与 $g : U \to V$ 是基底 $S$ 上的概形态射，则
$f \times g$ 是 $X \times_S U \to X \times_S V$（$g$ 的基变换）与
$X \times_S V \to Y \times_S V$（$f$ 的基变换）的复合。
因此 (5)、(6) 由 (1)--(4) 得出。
\end{proof}

\begin{lemma}
\label{schemes-lemma-compose-after-separated}
\begin{slogan}
分离态射与拟分离态射满足消去律。
\end{slogan}
令 $f : X \to Y$ 与 $g : Y \to Z$ 为概形态射。
若 $g \circ f$ 分离，则 $f$ 也分离；
若 $g \circ f$ 拟分离，则 $f$ 也拟分离。
\end{lemma}

\begin{proof}
假设 $g \circ f$ 分离。考虑其对角态射的分解
$X \to X \times_Y X \to X \times_Z X$；这里所说的对角态射属于
$g \circ f$。由引理 \ref{schemes-lemma-fibre-product-after-map}，后一个态射是浸入。
依假设，$X$ 在 $X \times_Z X$ 中的像闭，因而它在
$X \times_Y X$ 中也闭。因此由引理 \ref{schemes-lemma-immersion-when-closed}，
$X \to X \times_Y X$ 是闭浸入。

\medskip\noindent
假设 $g \circ f$ 拟分离。令 $V \subset Y$ 为映入 $Z$ 的某个仿射开集
的仿射开集。令 $U_1, U_2 \subset X$ 为映入 $V$ 的仿射开集。于是
$U_1 \cap U_2$ 是有限多个仿射开集的并，因为 $U_1, U_2$ 映入
$Z$ 的同一个仿射开集。又因为可以覆盖 $Y$，所用仿射开集形如 $V$，
本引理由引理
\ref{schemes-lemma-characterize-quasi-separated} 得出。
\end{proof}

\begin{lemma}
\label{schemes-lemma-quasi-compact-permanence}
令 $f : X \to Y$ 与 $g : Y \to Z$ 为概形态射。
若 $g \circ f$ 拟紧且 $g$ 拟分离，则 $f$ 拟紧。
\end{lemma}

\begin{proof}
这是因为 $f$ 等于复合
$(1, f) : X \to X \times_Z Y \to Y$。第一个映射由引理
\ref{schemes-lemma-section-immersion} 拟紧，因为它是拟分离态射
$X \times_Z Y \to X$ 的截面（该态射是 $g$ 的基变换；
见引理 \ref{schemes-lemma-separated-permanence}）。第二个映射也拟紧，
因为它是 $g \circ f$ 的基变换；见引理
\ref{schemes-lemma-quasi-compact-preserved-base-change}。最后，拟紧态射的复合拟紧；
见引理 \ref{schemes-lemma-composition-quasi-compact}。
\end{proof}

\begin{lemma}
\label{schemes-lemma-affine-separated}
仿射概形分离。从仿射概形到另一个概形的态射分离。
\end{lemma}

\begin{proof}
令 $U = \Spec(A)$ 为仿射概形。由引理
\ref{schemes-lemma-diagonal-affines-closed}，$U \to \Spec(\mathbf{Z})$ 的对角闭。
因此由定义 \ref{schemes-definition-separated}，$U$ 分离。
若 $U \to X$ 是概形态射，则可把引理
\ref{schemes-lemma-compose-after-separated} 应用于态射
$U \to X \to \Spec(\mathbf{Z})$，从而得出 $U \to X$ 分离。
\end{proof}

\noindent
读者或许会问：为了断定一对仿射开集的交仍仿射，
是否真的必须只考虑其像包含在某个仿射开集中的那些对？往往并不需要！

\begin{lemma}
\label{schemes-lemma-curiosity}
令 $f : X \to S$ 为态射。假设 $f$ 分离且 $S$ 是分离概形。
设 $U \subset X$ 与 $V \subset X$ 为仿射开集。则 $U \cap V$ 仿射
（并且是 $U \times V$ 的闭子概形）。
\end{lemma}

\begin{proof}
在此情形，由引理 \ref{schemes-lemma-separated-permanence}，$X$ 分离。
因此 $U \cap V$ 仿射；这是把引理 \ref{schemes-lemma-characterize-separated}
应用于态射 $X \to \Spec(\mathbf{Z})$ 所得的结论。
\end{proof}

\noindent
另一方面，下例说明不能指望仿射概形的像包含在某个仿射开集中。

\begin{example}
\label{schemes-example-image-affine-projective}
考虑非仿射概形 $U = \Spec(k[x, y]) \setminus \{(x, y)\}$，
见例 \ref{schemes-example-not-affine}。另一方面，考虑概形
$$
\mathbf{GL}_{2, k} = \Spec(k[a, b, c, d, 1/ad - bc]).
$$
存在态射 $\mathbf{GL}_{2, k} \to U$，它对应于环同态
$x \mapsto a$、$y \mapsto b$。容易看出该态射满射，故其像不包含在
$U$ 的任何仿射开集中。事实上，仿射概形 $\mathbf{GL}_{2, k}$ 也满射到
$\mathbf{P}^1_k$，而 $\mathbf{P}^1_k$ 甚至不能浸入{\it 任何}仿射概形。
\end{example}

\begin{remark}
\label{schemes-remark-quasi-compact-and-quasi-separated}
令 $P$ 为以下 $4$ 种性质之一：“拟分离”、“分离”、
“拟紧且拟分离”或“拟紧且分离”。则下列陈述成立：
\begin{enumerate}
\item 任意仿射概形都具有 $P$。
\item 给定概形态射 $f : X \to Y$，若 $Y$ 具有 $P$，则
$f$ 具有 $P$ 当且仅当 $X$ 具有 $P$。
\item 若 $X \to Y$ 与 $Z \to Y$ 为概形态射，并且
$X$、$Y$、$Z$ 都具有 $P$，则 $X \times_Y Z$ 具有 $P$。
\end{enumerate}
陈述 (1) 显然。令 $f : X \to Y$ 如 (2)。若 $f$ 具有 $P$，则 $X$ 也具有，
这是把引理 \ref{schemes-lemma-separated-permanence} 与
\ref{schemes-lemma-composition-quasi-compact} 应用于
$X \to Y \to \Spec(\mathbf{Z})$ 所得。若 $X$ 具有 $P$，则 $f$ 也具有，
这是把引理 \ref{schemes-lemma-compose-after-separated} 与
\ref{schemes-lemma-quasi-compact-permanence} 应用于
$X \to Y \to \Spec(\mathbf{Z})$ 所得。
令 $X \to Y$ 与 $Z \to Y$ 如 (3)。投影 $X \times_Y Z \to Z$ 具有 $P$，
因为它是箭头 $X \to Y$ 的基变换，而该箭头由 (2) 具有 $P$；
见引理 \ref{schemes-lemma-separated-permanence} 与
\ref{schemes-lemma-quasi-compact-preserved-base-change}。因此由 (2)，
$X \times_Y Z$ 具有 $P$。
\end{remark}





\section{分离性的赋值判据}
\label{schemes-section-valuative-separatedness}

\begin{lemma}
\label{schemes-lemma-separated-implies-valuative}
令 $f : X \to S$ 为概形态射。若 $f$ 分离，则 $f$ 满足
赋值判据的唯一性部分。
\end{lemma}

\begin{proof}
给定定义 \ref{schemes-definition-valuative-criterion} 中的一个图。假设有两个态射
$a, b : \Spec(A) \to X$ 嵌入该图。令 $Z \subset \Spec(A)$ 为 $a$ 与 $b$
的等化子。由引理 \ref{schemes-lemma-where-are-they-equal}，它是 $\Spec(A)$ 的
闭子概形。依假设，它包含 $\Spec(A)$ 的泛点。由于 $A$ 是整环，
这蕴含 $Z = \Spec(A)$。故如所需，$a = b$。
\end{proof}

\begin{lemma}[分离性的赋值判据]
\label{schemes-lemma-valuative-criterion-separatedness}
\begin{reference}
\cite[II Proposition 7.2.3]{EGA}
\end{reference}
令 $f : X \to S$ 为态射。假设：
\begin{enumerate}
\item 态射 $f$ 拟分离；
\item 态射 $f$ 满足赋值判据的唯一性部分。
\end{enumerate}
则 $f$ 分离。
\end{lemma}

\begin{proof}
由假设 (1)、命题 \ref{schemes-proposition-characterize-universally-closed}，
以及引理 \ref{schemes-lemma-diagonal-immersion}、\ref{schemes-lemma-immersion-when-closed}，
可知只须证明态射
$\Delta_{X/S} : X \to X \times_S X$ 满足赋值判据的存在性部分。
给定实线交换图
$$
\xymatrix{
\Spec(K) \ar[r] \ar[d] & X \ar[d] \\
\Spec(A) \ar[r] \ar@{-->}[ru] & X \times_S X
}
$$
右下箭头对应于一对态射 $a, b : \Spec(A) \to X$，它们位于 $S$ 上。
由 (2) 可知 $a = b$。因此取 $a$ 为虚线箭头即可。
\end{proof}





\section{单态射}
\label{schemes-section-monomorphisms}

\begin{definition}
\label{schemes-definition-monomorphism}
若一个概形态射是概形范畴中的单态射，则称它为{\it 单态射}；
见《范畴》，定义 \ref{categories-definition-mono-epi}。
\end{definition}

\begin{lemma}
\label{schemes-lemma-monomorphism}
\begin{slogan}
概形态射是单态射，当且仅当其对角态射是同构。
\end{slogan}
令 $j : X \to Y$ 为概形态射。则 $j$ 是单态射，当且仅当对角态射
$\Delta_{X/Y} : X \to X \times_Y X$ 是同构。
\end{lemma}

\begin{proof}
这在任何具有纤维积的范畴中都成立。
\end{proof}

\begin{lemma}
\label{schemes-lemma-monomorphism-separated}
概形单态射分离。
\end{lemma}

\begin{proof}
这是因为同构是闭浸入，并应用上面的引理 \ref{schemes-lemma-monomorphism}。
\end{proof}

\begin{lemma}
\label{schemes-lemma-composition-monomorphism}
单态射的复合是单态射。
\end{lemma}

\begin{proof}
这在任何范畴中都成立。
\end{proof}

\begin{lemma}
\label{schemes-lemma-base-change-monomorphism}
单态射的基变换是单态射。
\end{lemma}

\begin{proof}
这在任何具有纤维积的范畴中都成立。
\end{proof}

\begin{lemma}
\label{schemes-lemma-injective-points}
令 $j : X \to Y$ 为概形态射。若 $j$ 在点上单射，则 $j$ 分离。
\end{lemma}

\begin{proof}
令 $z$ 为 $X \times_Y X$ 的点。则 $x = \text{pr}_1(z)$ 与
$\text{pr}_2(z)$ 相同，因为 $j$ 把这两点映到同一点 $y$（属于 $Y$）。
于是可选取仿射开邻域 $V \subset Y$，它包含 $y$；以及仿射开邻域
$U \subset X$，它包含 $x$，并使得 $j(U) \subset V$。于是
$z \in U \times_V U \subset X \times_Y X$。
因此 $X \times_Y X$ 是仿射开集 $U \times_V U$ 的并。由于
$\Delta_{X/Y}^{-1}(U \times_V U) = U$，并且
$U \to U \times_V U$ 是闭浸入，可断定 $\Delta_{X/Y}$ 是闭浸入
（见引理 \ref{schemes-lemma-diagonal-immersion} 证明中的论证）。
\end{proof}

\begin{lemma}
\label{schemes-lemma-injective-points-surjective-stalks}
令 $j : X \to Y$ 为概形态射。若
\begin{enumerate}
\item $j$ 在点上单射；
\item 对任意 $x \in X$，环同态
$j^\sharp_x : \mathcal{O}_{Y, j(x)} \to \mathcal{O}_{X, x}$
满射，
\end{enumerate}
则 $j$ 是单态射。
\end{lemma}

\begin{proof}
令 $a, b : Z \to X$ 为两个概形态射，满足
$j \circ a  = j \circ b$。由 (1)，作为底拓扑空间映射有 $a = b$。
对任意 $z \in Z$，有
$a^\sharp_z \circ j^\sharp_{a(z)} = b^\sharp_z \circ j^\sharp_{b(z)}$，
这是从 $\mathcal{O}_{Y, j(a(z))} \to \mathcal{O}_{Z, z}$ 的映射。
映射 $j^\sharp_x$ 的满射性迫使
$a^\sharp_z = b^\sharp_z$，$\forall z \in Z$。这蕴含
$a^\sharp = b^\sharp$，因此如所需，作为概形态射有 $a = b$。
\end{proof}

\begin{lemma}
\label{schemes-lemma-immersions-monomorphisms}
概形浸入是单态射。特别地，任意浸入都分离。
\end{lemma}

\begin{proof}
可以通过检验引理 \ref{schemes-lemma-injective-points-surjective-stalks}
的判据适用来证明。或许更优雅的是，利用引理
\ref{schemes-lemma-restrict-map-to-opens} 与
\ref{schemes-lemma-characterize-closed-subspace}：它们蕴含开浸入和闭浸入
都是单态射，因而任意浸入（它是这类态射的复合）都是单态射。
\end{proof}

\begin{lemma}
\label{schemes-lemma-subscheme-of-separated-scheme}
令 $f : X \to S$ 为分离态射。任意局部闭子概形
$Z \subset X$ 在 $S$ 上分离。
\end{lemma}

\begin{proof}
这由引理 \ref{schemes-lemma-immersions-monomorphisms} 以及分离态射的复合分离
这一事实得出（引理 \ref{schemes-lemma-separated-permanence}）。
\end{proof}

\begin{example}
\label{schemes-example-Q-over-Z}
态射 $\Spec(\mathbf{Q}) \to \Spec(\mathbf{Z})$ 是单态射，
因为 $\mathbf{Q} \otimes_{\mathbf{Z}} \mathbf{Q} = \mathbf{Q}$。
更一般地，对任意概形 $S$ 及任意点 $s \in S$，典范态射
$$
\Spec(\mathcal{O}_{S, s}) \longrightarrow S
$$
是单态射。
\end{example}

\begin{lemma}
\label{schemes-lemma-mono-towards-spec-field}
令 $k_1, \ldots, k_n$ 为域。对任意概形单态射
$X \to \Spec(k_1 \times \ldots \times k_n)$，存在子集
$I \subset \{1, \ldots, n\}$，使得
$X \cong \Spec(\prod_{i \in I} k_i)$；这里把它们视为
$\Spec(k_1 \times \ldots \times k_n)$ 上的概形。
更一般地，若 $X = \coprod_{i \in I} \Spec(k_i)$ 是域的谱的不交并，
且 $Y \to X$ 是单态射，则存在子集 $J \subset I$，使得
$Y = \coprod_{i \in J} \Spec(k_i)$。
\end{lemma}

\begin{proof}
先归结为 $n = 1$（或 $\# I = 1$）的情形；方法是取开闭子概形
$\Spec(k_i)$ 的逆像。此时 $X$ 只有一个点，因而是仿射的。
相应的代数问题如下：若 $k \to R$ 为代数同态，并且
$R \otimes_k R \cong R$，则 $R \cong k$ 或 $R = 0$。
这由维数理由成立。另见《代数》，
引理 \ref{algebra-lemma-epimorphism-over-field}。
\end{proof}










\section{拟凝聚模的函子性}
\label{schemes-section-quasi-coherent}

\noindent
令 $X$ 为概形。记 $\QCoh(\mathcal{O}_X)$ 为拟凝聚
$\mathcal{O}_X$-模所成的范畴，定义见《模》，
定义 \ref{modules-definition-quasi-coherent}。我们已在第
\ref{schemes-section-quasi-coherent-affine} 节看到，范畴 $\QCoh(\mathcal{O}_X)$
在 $X$ 仿射时具有许多良好性质。由于拟凝聚性在 $X$ 上是局部的，
这些性质也由任意概形 $X$ 上的拟凝聚层范畴继承。现列举如下。
\begin{enumerate}
\item $\mathcal{O}_X$-模层 $\mathcal{F}$ 拟凝聚，当且仅当
$\mathcal{F}$ 在每个仿射开集 $U = \Spec(R)$ 上的限制都形如
$\widetilde M$，其中它来自某个 $R$-模 $M$。
\item $\mathcal{O}_X$-模层 $\mathcal{F}$ 拟凝聚，当且仅当
$\mathcal{F}$ 在某个仿射开覆盖的每个成员上的限制拟凝聚。
\item 拟凝聚层的任意直和拟凝聚。
\item 拟凝聚层的任意余极限拟凝聚。
\item
\label{schemes-item-kernel}
拟凝聚层态射的核与余核拟凝聚。
\item 给定 $\mathcal{O}_X$-模的短正合列
$0 \to \mathcal{F}_1 \to \mathcal{F}_2 \to \mathcal{F}_3 \to 0$
，若其中两个拟凝聚，则第三个也拟凝聚。
\item 给定概形态射 $f : Y \to X$，拟凝聚 $\mathcal{O}_X$-模的拉回是
拟凝聚 $\mathcal{O}_Y$-模。见《模》，引理
\ref{modules-lemma-pullback-quasi-coherent}。
\item 给定两个拟凝聚 $\mathcal{O}_X$-模，其张量积拟凝聚；见《模》，
引理 \ref{modules-lemma-tensor-product-permanence}。
\item 给定拟凝聚 $\mathcal{O}_X$-模 $\mathcal{F}$，其上的张量代数、
对称代数与外代数 $\mathcal{F}$ 都拟凝聚；见《模》，引理
\ref{modules-lemma-whole-tensor-algebra-permanence}。
\item 给定两个拟凝聚 $\mathcal{O}_X$-模 $\mathcal{F}$、$\mathcal{G}$，
并且 $\mathcal{F}$ 是有限表示的，则内部 Hom
$\SheafHom_{\mathcal{O}_X}(\mathcal{F}, \mathcal{G})$
拟凝聚；见《模》，引理
\ref{modules-lemma-internal-hom-locally-kernel-direct-sum} 及上面的
(\ref{schemes-item-kernel})。
\end{enumerate}
另一方面，拟凝聚模的推前一般并不拟凝聚。下面给出它确实拟凝聚的一种情形。

\begin{lemma}
\label{schemes-lemma-push-forward-quasi-coherent}
令 $f : X \to S$ 为概形态射。若 $f$ 拟紧且拟分离，则
$f_*$ 把拟凝聚 $\mathcal{O}_X$-模变为拟凝聚 $\mathcal{O}_S$-模。
\end{lemma}

\begin{proof}
此问题在 $S$ 上是局部的，故可假设 $S$ 仿射。由于 $X$ 拟紧，可以写成
$X = \bigcup_{i = 1}^n U_i$，其中每个 $U_i$ 都是仿射开集。
由于 $f$ 拟分离，可以把
$U_i \cap U_j = \bigcup_{k = 1}^{n_{ij}} U_{ijk}$ 写成某些仿射开集
$U_{ijk}$ 的并；见引理 \ref{schemes-lemma-characterize-quasi-separated}。
记 $f_i : U_i \to S$ 与 $f_{ijk} : U_{ijk} \to S$ 为 $f$ 的限制。
对任意开集 $V$（属于 $S$）以及任意层 $\mathcal{F}$（定义在 $X$ 上），有
\begin{eqnarray*}
f_*\mathcal{F}(V) & = & \mathcal{F}(f^{-1}V) \\
& = &
\Ker\left(
\bigoplus\nolimits_i \mathcal{F}(f^{-1}V \cap U_i)
\to
\bigoplus\nolimits_{i, j, k} \mathcal{F}(f^{-1}V \cap U_{ijk})\right) \\
& = &
\Ker\left(
\bigoplus\nolimits_i f_{i, *}(\mathcal{F}|_{U_i})(V)
\to
\bigoplus\nolimits_{i, j, k} f_{ijk, *}(\mathcal{F}|_{U_{ijk}})(V)\right) \\
& = &
\Ker\left(
\bigoplus\nolimits_i f_{i, *}(\mathcal{F}|_{U_i})
\to
\bigoplus\nolimits_{i, j, k} f_{ijk, *}(\mathcal{F}|_{U_{ijk}})\right)(V)
\end{eqnarray*}
换言之，有层的正合列
$$
0 \to f_*\mathcal{F}
\to \bigoplus f_{i, *}\mathcal{F}_i
\to \bigoplus f_{ijk, *}\mathcal{F}_{ijk}
$$
其中 $\mathcal{F}_i, \mathcal{F}_{ijk}$ 表示 $\mathcal{F}$ 在相应开集上的限制。
若 $\mathcal{F}$ 是拟凝聚 $\mathcal{O}_X$-模，则 $\mathcal{F}_i$ 是拟凝聚
$\mathcal{O}_{U_i}$-模，而 $\mathcal{F}_{ijk}$ 是拟凝聚
$\mathcal{O}_{U_{ijk}}$-模。因此由引理 \ref{schemes-lemma-widetilde-pullback}，
正合列的第二、第三项都是拟凝聚 $\mathcal{O}_S$-模。故
$f_*\mathcal{F}$ 是拟凝聚 $\mathcal{O}_S$-模。
\end{proof}

\noindent
利用这一点，可以如下刻画概形的（闭）浸入。

\begin{lemma}
\label{schemes-lemma-characterize-closed-immersions}
令 $f : X \to Y$ 为概形态射。假设：
\begin{enumerate}
\item $f$ 诱导 $X$ 到 $Y$ 的某个闭子集上的同胚；
\item $f^\sharp : \mathcal{O}_Y \to f_*\mathcal{O}_X$
满射。
\end{enumerate}
则 $f$ 是概形的闭浸入。
\end{lemma}

\begin{proof}
假设 (1)、(2)。由 (1)，态射 $f$ 拟紧
（见《拓扑》，引理 \ref{topology-lemma-closed-in-quasi-compact}）。
条件 (1)、(2) 蕴含引理 \ref{schemes-lemma-injective-points-surjective-stalks}
的条件 (1)、(2)。因此 $f : X \to Y$ 是单态射。特别地，$f$ 分离；
见引理 \ref{schemes-lemma-monomorphism-separated}。于是可以应用上面的引理
\ref{schemes-lemma-push-forward-quasi-coherent}，得出 $f_*\mathcal{O}_X$ 是拟凝聚
$\mathcal{O}_Y$-模。因此，$\mathcal{O}_Y \to f_*\mathcal{O}_X$ 的核拟凝聚；
见引理 \ref{schemes-lemma-extension-quasi-coherent}。由于拟凝聚层局部由截面生成
（见《模》，定义 \ref{modules-definition-quasi-coherent}），这蕴含 $f$
是闭浸入；见定义 \ref{schemes-definition-closed-immersion-locally-ringed-spaces}。
\end{proof}

\noindent
可以用这个引理证明下面的引理。

\begin{lemma}
\label{schemes-lemma-composition-immersion}
概形浸入的复合是浸入；概形闭浸入的复合是闭浸入；
概形开浸入的复合是开浸入。
\end{lemma}

\begin{proof}
开浸入的情形显然，因为开子空间的开子空间仍是开子空间。

\medskip\noindent
设 $a : Z \to Y$ 与 $b : Y \to X$ 为概形闭浸入。我们验证
$c = b \circ a$ 也是闭浸入。假设蕴含 $a$ 与 $b$ 都是到闭子集上的同胚，
因而 $c = b \circ a$ 也是到闭子集上的同胚。此外，映射
$\mathcal{O}_X \to c_*\mathcal{O}_Z$ 满射，因为它分解为满射
$\mathcal{O}_X \to b_*\mathcal{O}_Y$ 与
$b_*\mathcal{O}_Y \to b_*a_*\mathcal{O}_Z$ 的复合（后者满射是因为
$b_*$ 正合；见《模》，引理 \ref{modules-lemma-i-star-exact}）。
因此由上面的引理 \ref{schemes-lemma-characterize-closed-immersions}，$c$ 是闭浸入。

\medskip\noindent
最后处理一般浸入。设 $a : Z \to Y$ 与 $b : Y \to X$ 为概形浸入。
这意味着存在开子概形 $V \subset Y$ 与 $U \subset X$，使得
$a(Z) \subset V$、$b(Y) \subset U$，且 $a : Z \to V$ 与
$b : Y \to U$ 都是闭浸入。由于 $Y$ 上的拓扑由 $U$ 上的拓扑诱导，
可以找到开集 $U' \subset U$，使得 $V = b^{-1}(U')$。于是
$Z \to V = b^{-1}(U') \to U'$ 是闭浸入的复合，因而是闭浸入。
这证明 $Z \to X$ 是浸入，结论得证。
\end{proof}
